% Seção 4: Discussão dos Resultados
% Conteúdo incluído via % Seção 4: Discussão dos Resultados
% Conteúdo incluído via % Seção 4: Discussão dos Resultados
% Conteúdo incluído via % Seção 4: Discussão dos Resultados
% Conteúdo incluído via \input{resultados} no main.tex

A literatura empírica sobre os instrumentos da Política Nacional de Desenvolvimento Regional apresenta expressiva assimetria em termos de volume de evidências: os Fundos Constitucionais concentram a quase totalidade dos estudos, enquanto os Fundos de Desenvolvimento e os Incentivos Fiscais constituem vertentes consideravelmente menos desenvolvidas. Os resultados obtidos variam substancialmente conforme a especificação econométrica adotada, o período amostral, a escala geográfica de análise e o controle de características locais não observáveis. Esta seção apresenta os principais achados da literatura, organizados por instrumento e abordagem metodológica, buscando identificar para cada um deles convergências, divergências e condições que modulam a eficácia da política.

\subsection{Avaliações de Impacto dos Fundos Constitucionais sobre o PIB}
\label{subsec:fc-pib}

Os Fundos Constitucionais são os instrumentos da PNDR mais amplamente avaliados pela literatura empírica: dentre os 34 estudos aprovados nesta revisão, 28 investigam os efeitos desses fundos sobre variáveis de resultado relacionadas ao produto interno bruto ou ao mercado de trabalho. A questão central é se esses fundos --- que operam via crédito subsidiado de ampla capilaridade territorial --- efetivamente contribuem para o crescimento econômico dos municípios beneficiados e, em caso afirmativo, sob quais condições e com que magnitude. O FNE é o fundo com maior número de avaliações (24 estudos), seguido pelo FCO (10 estudos) e pelo FNO (8 estudos). A predominância dos Fundos Constitucionais na literatura decorre de sua maior longevidade em relação aos demais instrumentos da PNDR -- com operação contínua desde 1989 --, da disponibilidade de dados municipais anuais de PIB a partir de 2002, do elevado volume de recursos envolvidos e da capilaridade dos programas em todos os municípios das macrorregiões-alvo. A quase totalidade dos estudos adota o PIB \textit{per capita} municipal como variável de resultado, tratando-o como aproximação da renda média local, dada a escassez de alternativas em frequência anual e cobertura municipal. Os resultados obtidos pela literatura variam consideravelmente conforme a especificação econométrica adotada, o período amostral, a escala geográfica de análise e o controle de características locais não observáveis. As subseções seguintes organizam os achados por abordagem metodológica, permitindo identificar padrões gerais e divergências nas estimativas de impacto. O Quadro~\ref{tab:estudos_fc_pib} apresenta a síntese dos 21 estudos que avaliam os efeitos dos Fundos Constitucionais sobre o PIB \textit{per capita}, organizados por método de análise.

\input{tabelas/tabela_macros}
\input{tabelas/survey_artigos_fc_pib}

\subsubsection{Modelos de Efeitos Fixos}

Modelos de painel com efeitos fixos controlam características não observáveis constantes no tempo (como geografia e instituições locais) e choques temporais comuns a todas as unidades (como crises macroeconômicas), permitindo isolar de forma mais precisa o impacto da política regional. Essa abordagem foi a primeira empregada para avaliar os Fundos Constitucionais sobre o PIB municipal, aproveitando a disponibilidade de dados em painel desde 2002.

\citeonline{Resende2014a} analisa a influência do FNE sobre o crescimento do PIB \textit{per capita} dos municípios nordestinos no período 2004-2010, utilizando médias em três subperíodos (2004-2006, 2006-2008 e 2008-2010) para minimizar viés de ciclos de negócios. Os resultados apontam impacto positivo do FNE-total em nível municipal (0,021) e microrregional (0,032), mas não significativo em nível mesorregional, o que Resende (2014a) atribui à maior heterogeneidade em regiões geográficas maiores. As estimativas sugerem que o resultado positivo estaria concentrado nos empréstimos ao setor agropecuário, porém não se mostraram robustas à inclusão de controles temporais. Para o FNO, aplicando a mesma especificação econométrica à região Norte, \citeonline{Resende2014a} obtém relação inversa entre o fundo e a expansão do produto municipal, com ausência de robustez nas análises setoriais quando se controla por períodos.

\citeonline{Oliveira2017a} aplicam especificações de efeitos fixos juntamente com outros métodos (PSM e GPS) para avaliar o FCO no estado de Goiás, obtendo evidências de impacto positivo sobre a dinâmica econômica municipal, com magnitude variando conforme o método empregado.

Mais recentemente, \citeonline{Filho2024} adotam painel de efeitos fixos para avaliar conjuntamente os três fundos constitucionais, não obtendo evidências consistentes de associação positiva com o crescimento do PIB \textit{per capita} para o período 2016-2019, resultado que contrasta com os estudos anteriores.

Esses estudos ressaltam limitações metodológicas importantes. A principal dificuldade consiste na identificação causal, dado que a alocação dos fundos é endógena à demanda por crédito subsidiado, o que tende a privilegiar municípios com maior capacidade de absorção e infraestrutura socioeconômica mais desenvolvida. Adicionalmente, identificam-se dificuldades em mensurar os custos fiscais efetivos dos programas, a possibilidade de existência de peso morto (investimentos que seriam realizados mesmo sem subsídio) e potenciais deslocamentos de fatores produtivos entre municípios vizinhos. A ausência de robustez à inclusão de controles temporais sugere que parte do impacto estimado pode capturar tendências macroeconômicas comuns aos municípios tratados e não tratados, em vez do efeito causal dos fundos. Em síntese, os modelos de efeitos fixos indicam correlação positiva dos fundos constitucionais com a expansão econômica local, com impactos mais pronunciados em nível municipal e maior magnitude para o FCO, mas a fragilidade da identificação causal e a ausência de robustez em diversas especificações limitam a interpretação causal dos resultados.

\subsubsection{Modelos de Painel Espacial}

Uma vertente da literatura admite a possibilidade de interdependência econômica entre unidades espaciais geograficamente próximas, o que pode ocasionar efeitos indiretos de transbordamento da política para outras regiões, expandindo a área potencialmente beneficiada. Por outro lado, o efeito poderia consistir em vazamentos de recursos das regiões-alvo da política em direção a regiões de maior dinamismo econômico. Portanto, entender a natureza da dinâmica espacial é importante para um planejamento mais eficiente da política. \citeonline{Resende2014c} utilizam modelo de painel SDM (\textit{Spatial Durbin Model}) e estimam efeito significativo do FCO sobre o PIB \textit{per capita}, com evidência de que uma elevação de 1 ponto percentual (p.p.) na relação FCO/PIB está associada a um incremento de cerca de 0,09 p.p. na taxa média anual de crescimento do PIB \textit{per capita} municipal nos dois anos seguintes. Na análise setorial, o resultado seria atribuído à modalidade de crédito FCO-Empresarial, em que o efeito seria na ordem de 0,13 p.p. de incremento. Porém, o resultado é não significativo para o FNE e negativo para o FNO. Com respeito aos efeitos indiretos (transbordamento espacial), obteve-se resultado relevante, mas não positivo, para o FCO, o que pode estar ligado à migração de fatores de produção.

\citeonline{Oliveira2015} usam modelo espacial para avaliar efeitos do FCO nos municípios de Goiás e concluem que efeitos diretos são menores quando se admite a ocorrência de efeitos espaciais, sugerindo que modelos não espaciais tendem a superestimar o impacto real dos fundos. Entretanto, as evidências obtidas por \citeonline{Resende2014c} e \citeonline{Oliveira2015} apontam em direção contrária, com os efeitos diretos do FCO maiores em modelos espaciais do que em modelos não espaciais. Essa divergência pode estar associada a diferenças na cobertura geográfica (Goiás versus toda a região Centro-Oeste), no tipo de matriz de pesos espaciais empregada (\textit{queen} de contiguidade versus distância inversa) e no período de análise considerado. Os resultados de \citeonline{Oliveira2015} indicam que cada 1\% de aumento no valor dos financiamentos está associado a aproximadamente 0,1\% de crescimento no PIB \textit{per capita} do município, usando especificação SAR com matriz espacial \textit{queen} do tipo contiguidade e dados de 2004 a 2011.

\citeonline{ResendeSilvaFilho2017} distinguem-se dos trabalhos anteriores por incorporar na análise a heterogeneidade definida pela própria PNDR por meio das tipologias regionais. Em modelo não espacial de efeitos fixos, o efeito estimado do FNE é significativo para municípios pertencentes às tipologias dinâmica e baixa renda. Quando se admite dependência espacial, há evidências de transbordamento espacial dos investimentos do FNE quando o município diretamente beneficiado tem tipologia dinâmica. As estimativas apontam que para cada 1 p.p. de aumento na relação FNE/PIB em municípios dinâmicos há um incremento de 0,07 p.p. na taxa média anual de crescimento do PIB \textit{per capita} do próprio município nos três anos seguintes e de 0,33 p.p. na taxa média dos municípios vizinhos. Para os autores, esses resultados ressaltam a relevância de se considerar a dimensão espacial nas análises e também o caráter heterogêneo dos efeitos do FNE nos territórios. \citeonline{ResendeSilvaFilho2018} adotam a mesma abordagem econométrica, mas incluindo também as regiões Norte e Centro-Oeste, a fim de avaliar conjuntamente os fundos FNE, FNO e FCO. Os resultados são positivos para efeitos diretos e indiretos em municípios de tipologia dinâmica quando se usam valores acumulados na variável dos fundos (dois primeiros anos de cada subperíodo). Os resultados apontam que um aumento de 1 p.p. na relação FC/PIB em municípios dinâmicos está associado a 0,03 p.p. de incremento no crescimento médio do PIB \textit{per capita} do próprio município pelos três anos seguintes e 0,12 p.p. de incremento nos municípios vizinhos. \citeonline{Filho2024} não obtêm evidências de associação espacial positiva entre os financiamentos do FNE, FNO e FCO nos anos de 2016 a 2018 e o crescimento médio do PIB \textit{per capita} dos municípios de 2016 a 2019, porém, o resultado é positivo sobre o PIB \textit{per capita} de 2019. Os autores concluem que os efeitos dos fundos são em sua maior parte esporádicos e pontuais.

Em resumo, os estudos convergem para a relevância de se incorporar a dimensão espacial nas análises de efeito dos fundos constitucionais, justificada pela ocorrência de correlação espacial significativa nas variáveis envolvidas. Porém, há menor consenso quando se trata da magnitude dos efeitos. Em geral, o resultado estaria condicionado a outras características econômicas locais. O efeito direto mostrou-se significativo em municípios de tipologia dinâmica e microrregiões de baixa renda, com incrementos de 0,03 a 0,14 p.p. no crescimento anual do PIB \textit{per capita} municipal por até dois anos, para cada 1 p.p. de aumento na relação Fundo/PIB. Para o efeito indireto, as estimativas variam entre 0,1 e 0,4 p.p. de incremento no crescimento para cada 1 p.p. de aumento em Fundo/PIB, também concentrado em municípios de tipologia dinâmica e microrregiões de baixa renda. Em análise mais desagregada, evidências apontam para associação entre o FCO e taxa de crescimento, atribuída à linha de crédito empresarial. No FNE, os resultados estariam ligados ao setor industrial e de serviços. Já para o FNO não há nenhuma evidência positiva e significativa. No aspecto metodológico, nota-se a necessidade de estudos mais robustos, que combinem intervalo de tempo mais longo, especificações espaciais alternativas, correção de endogeneidade e maior número de controles (inclusive variáveis de confusão da própria PNDR).

\subsubsection{Modelos Não Lineares e de Eficiência}

\citeonline{Linharesetal2014} usam modelo \textit{threshold} e obtêm evidências de que o efeito do FNE é pouco significativo para municípios com PIB \textit{per capita} inicial abaixo de R\$ 2.143 ou acima de R\$ 7.406. Entre esses limites, os autores identificam dois grupos distintos. O primeiro é formado por municípios de PIB \textit{per capita} inicial entre R\$ 2.143 e R\$ 3.866. Nesse grupo, estima-se que 1 p.p. de aumento médio em FNE/PIB no período 2002-2006 esteja associado a um incremento de 0,08 p.p. no crescimento médio do PIB \textit{per capita} entre 2002 e 2008. No segundo grupo, de municípios com PIB \textit{per capita} inicial entre R\$ 3.866 e R\$ 7.406, o efeito estimado foi maior, na ordem de 0,11 p.p. Os resultados sugerem que o FNE tem maior efetividade em municípios com níveis intermediários de renda, reafirmando os achados de \citeonline{ResendeSilvaFilho2017} e \citeonline{ResendeSilvaFilho2018}. Em municípios de baixa renda, porém, os resultados não foram significativos, o que indica a necessidade de aprimoramento dos mecanismos da política para melhor adequação às características e necessidades dessas regiões.

\citeonline{Oliveira2015} usam modelagem quantílica para estimar a significância e heterogeneidade de efeitos do FNE ao longo da distribuição de crescimento dos municípios, empregando estimação em dois estágios e a localização de agências bancárias do BNB como instrumento para a alocação de recursos do fundo. Os efeitos foram positivos e significativos para o FNE total e para as linhas de crédito associadas à agricultura e pecuária. No entanto, os largos intervalos de confiança não permitiram constatar presença de heterogeneidade ao longo da curva de crescimento. \citeonline{Oliveira2018} analisa a relação entre a participação do FNE e a eficiência média dos municípios nordestinos. Usando modelo de análise envoltória de dados e de fronteira estocástica, o autor identifica elevada eficiência nos municípios pertencentes à região do MATOPIBA, associada à agricultura extensiva voltada à exportação. Adicionalmente, municípios do Semiárido indicaram menor eficiência no crescimento do produto médio, enquanto municípios menos densos e mais distantes das respectivas capitais e com empresas menores tendem a fazer melhor uso dos financiamentos.

\subsubsection{Modelos Dinâmicos}

\citeonline{Cambota2019} adotam modelo de painel dinâmico e instrumentos sintéticos para corrigir o viés de endogeneidade na variável PIB defasada e na variável de interesse FNE/PIB, usando dados de 2003 a 2014. Os autores confirmam a importância desse tipo de controle para a obtenção de estimativas mais realistas do impacto da política. As estimativas apontam associação positiva entre o FNE e o crescimento econômico dos municípios, na ordem de 2,96 p.p. para cada 1 p.p. de aumento em FNE/PIB no ano anterior. Cabe notar que essa magnitude é consideravelmente superior à dos demais estudos revisados, cujos coeficientes situam-se tipicamente entre 0,03 e 0,13 p.p. A discrepância pode ser atribuída à diferente especificação do modelo (painel dinâmico com instrumentos sintéticos), ao período mais longo de análise (2003-2014) ou a possível sobreestimação decorrente da proliferação de instrumentos no procedimento GMM, aspecto reconhecido na literatura como potencial fonte de viés.

\subsubsection{Modelos de Equilíbrio Geral}

Em \citeonline{Nascimento2017}, um cenário hipotético de retirada do FNE e realocação dos recursos em gastos correntes do governo resultaria a longo prazo em uma redução líquida de 0,16\% no PIB nacional e em maior concentração regional da produção. O estudo também sugere que os fundos têm papel mais relevante no crescimento de longo prazo, via elevação do estoque de capital e da capacidade produtiva local, frente à alternativa apontada. No curto prazo, ambos os cenários seriam equivalentes em termos de ativação da demanda. \citeonline{Ribeiroetal2020} indicam que os investimentos do FNE realizados em 2014 e 2015 teriam elevado o PIB da região Nordeste em 3,51\% até o ano 2025. \citeonline{Oliveira2021} constroem dois cenários alternativos envolvendo o término da política de subsídios, e ambos indicam como consequência a redução na atividade econômica na região Nordeste, puxada principalmente pela redução do produto da atividade agropecuária.

\subsubsection{Modelos Quase Experimentais}

\citeonline{Resende2014c} usa o método de primeira diferença para avaliar o efeito de empréstimos industriais do FNE sobre o crescimento do PIB \textit{per capita} municipal. Usando taxas de crescimento do PIB \textit{per capita} nos períodos 2000-2003 e 2000-2006, o autor não encontra efeito significativo das operações envolvendo o FNE-Indústria durante o ano 2000, seja em nível municipal ou microrregional. Argumenta-se que o resultado pode estar associado à elevada participação do setor agropecuário nas operações do FNE, setor em que há predomínio da informalidade, o que dificulta sua captura pelas medições do PIB.

\citeonline{Oliveira2020} adotam desenho quase-experimental que combina regressão com descontinuidade geográfica (RDD) e modelo de diferenças em diferenças. Os autores constroem um painel de dados decenal (1980, 1990, 2000 e 2010) composto por municípios dos estados de Minas Gerais e Espírito Santo que ingressaram nos critérios de elegibilidade ao FNE em períodos recentes (dados pré e pós-ingresso) e por municípios vizinhos imediatos que nunca tenham sido tratados pelo FNE. Os resultados obtidos levantam dúvidas sobre a efetividade do FNE, sugerindo que o Fundo não teve influência sobre as variáveis de produto e renda nos municípios avaliados no ano 2010. O estudo ressalta a importância de melhor vinculação do crédito ao aumento da produtividade, mediante condições diferenciadas e combinação com outros instrumentos.

Por outro lado, \citeonline{Carneiroetal2024a} obtêm evidências de efeito positivo e significativo do FNE sobre o PIB \textit{per capita} municipal. Os autores usam dados de 2002 a 2019 junto a estimadores DID dois estágios e DID escalonado, combinando três instrumentos de política regional: FNE, FDNE e incentivo fiscal da SUDENE. O estudo aponta efeito médio de tratamento de 8,8\% a 11\% no PIB \textit{per capita} do município beneficiado em relação ao grupo de controle, sendo atribuído principalmente ao setor serviços e administração pública. O estudo de eventos aponta que o efeito de tratamento é crescente e inicia no quinto ano após o primeiro tratamento. Quanto à sinergia com os demais instrumentos, o resultado é não significativo. \citeonline{MonteIrffiBastosCarneiro2025} exploram heterogeneidades de gênero, tipo de tomador e magnitude dos empréstimos (efeito dosagem) do FNE, aplicando o método \textit{Generalized Propensity Score} (GPS). O estudo encontra correlação positiva entre a participação do FNE-Mulheres no valor total do programa e a variação do PIB \textit{per capita} municipal, o que sinaliza uma eficácia maior do programa nesse grupo populacional. Efeito semelhante é obtido em relação à variável participação de tomadores pessoa jurídica no total de financiamentos. Com respeito ao efeito dosagem, os resultados apontam efeitos positivos no PIB \textit{per capita} somente na última e mais elevada dose de tratamento.

Em conjunto, as evidências sobre o efeito dos Fundos Constitucionais no PIB \textit{per capita} revelam um quadro heterogêneo e condicionado à especificação metodológica adotada. Os modelos de efeitos fixos e espaciais indicam associação positiva para o FCO e resultados ambíguos para FNE e FNO, enquanto os modelos quase experimentais apontam tanto para ausência de efeito quanto para impactos positivos significativos. Os modelos não lineares e de eficiência sugerem que a eficácia dos fundos é condicionada ao nível inicial de renda do município, com efeitos mais pronunciados em faixas intermediárias de desenvolvimento. As simulações de equilíbrio geral reafirmam a importância dos fundos para a atividade econômica regional em cenários contrafactuais de retirada da política. A constatação de que os efeitos são mais pronunciados em municípios com níveis intermediários de renda, e nulos ou inconclusivos em municípios de baixa renda, pode ser interpretada à luz da literatura sobre capacidade de absorção e armadilhas de pobreza: municípios de baixa renda podem carecer do capital humano e da infraestrutura básica necessários para converter o crédito subsidiado em crescimento econômico sustentado; municípios de renda intermediária possuiriam o mínimo de capacidade necessária para ativar o efeito multiplicador dos investimentos; e municípios de alta renda estariam próximos do efeito marginal decrescente do capital.

\subsection{Avaliações de Impacto dos Fundos Constitucionais sobre o Mercado de Trabalho}

Dezesseis estudos aprovados nesta revisão avaliam o impacto dos Fundos Constitucionais sobre variáveis do mercado de trabalho, com foco predominante no emprego formal, massa salarial e salário médio. O Quadro~\ref{tab:estudos_fc_emprego} apresenta a síntese desses estudos, organizados por método de análise.

\input{tabelas/survey_artigos_fc_emprego}

\citeonline{Silva2009} foi o primeiro estudo a investigar o impacto dos empréstimos dos FCs sobre a contratação de trabalhadores nas empresas beneficiadas. O estudo faz a interseção dos registros de empresas beneficiadas com os dados da RAIS (Relação Anual de Informações Sociais), obtendo a quantidade de vínculos de empregos formais de 449 empresas (NE=211, N=174, CO=74) nos anos de 2000 e 2003. Aplicando o método \textit{Propensity Score Matching (PSM)}, os autores concluem que nas empresas beneficiadas pelo FNE a quantidade de empregos cresceu em média 61,2 pontos percentuais a mais do que nas empresas não beneficiadas. No entanto, os valores não foram significativos para o salário médio. E ainda, para o FNO e o FCO, não foi obtido efeito positivo em nenhum dos indicadores avaliados.

\citeonline{Soares2017} também aplicam o método PSM, mas fazem uso de um conjunto mais rico de informações sobre o FNE, o que permitiu ampliar consideravelmente o número de empresas e o horizonte de tempo investigado. Além disso, os autores fazem uma análise ano a ano dos efeitos do fundo, o que permite estimativas mais eficientes, com quantidades maiores de empresas nos primeiros anos pós-tratamento. Com dados do período de 2000 a 2006, no horizonte de 3 anos pós-tratamento há 1.104 empresas e em 5 anos há 314. Os resultados obtidos reforçam as evidências obtidas por \citeonline{Silva2009}. No primeiro ano pós-tratamento o crescimento do estoque de empregos é em média 7,96 pontos percentuais maior nas empresas beneficiadas, relativamente às demais. Além disso, a trajetória é crescente nos anos subsequentes, alcançando 33,59 pontos percentuais no terceiro ano e 132,23 no quinto ano. Se confirmados, esses resultados atribuem grande relevância para a política regional, indicando que o FNE tem papel importante na geração de empregos formais nas empresas beneficiadas. Os autores também destacam que o efeito de maior monta estaria associado às empresas financiadas entre 1999 e 2001, as únicas no horizonte de 5 anos pós-tratamento, estando mais concentradas no setor industrial. O mesmo estudo ainda analisa possíveis efeitos do FNE sobre a massa salarial, obtendo resultados significativos e com dinâmica semelhante ao do efeito no emprego. Como consequência, a variável salário médio não apresentou variação significativa como decorrência do FNE, indicando que as novas contratações teriam ocorrido ao mesmo nível dos salários vigentes.

\citeonline{Oliveira2018} apresentam como inovação principal o uso de variável de tratamento do tipo contínua, permitindo analisar possíveis heterogeneidades no efeito do fundo decorrentes da intensidade do tratamento (valor do empréstimo), em contraste à abordagem dicotômica de tratamento presente nos estudos anteriores. Os autores usam o método \textit{Generalized Propensity Score Matching (GPS)} e dados de financiamentos a empresas do estado de Goiás realizados pelo FCO por meio da linha Empresarial (associada à indústria, comércio e serviços) no período de 2004 a 2011. Os resultados indicam relação positiva e significativa com o crescimento do emprego apenas em um intervalo específico da distribuição de tratamento e limitada ao período 2004-2008 (pré-crise), cenário semelhante ao obtido na abordagem PSM com variável de tratamento dicotômica. Com respeito ao salário médio, não se obteve nenhum resultado significativo. \citeonline{Oliveira2018} expandem a abordagem de \citeonline{Oliveira2018} para os demais fundos constitucionais. O dado mais relevante obtido indica efeito decrescente dos fundos sobre as variáveis de interesse, o que significa que quanto maior o valor do empréstimo, menor o efeito sobre o crescimento do emprego e do salário. Para os autores, isso se justifica do ponto de vista teórico pela propriedade de retornos marginais decrescentes da função de produção neoclássica.

\citeonline{Junior2024} contribuem com a literatura ao usar o método de função dose-resposta para captar efeitos heterogêneos e correção de viés de seleção por meio de variável instrumental. Para garantir comparabilidade, os autores optam por usar a mesma base de dados de \citeonline{Soares2017}. Na variável crescimento do emprego, o modelo IV aponta efeito médio de 357,54\%, acima de 133,94\% do modelo OLS. Para o crescimento da massa salarial, o valor médio é 220,96\% no modelo IV contra 188,55\% em OLS. Em ambas as variáveis a restrição de exogeneidade leva a subestimação do efeito médio e superestimação do ponto ótimo de tratamento, sinalizando para efeitos mais significativos em pontos mais baixos da distribuição de tratamento, quando a endogeneidade é corrigida.

\citeonline{Resende2014c} contribui para uma limitação importante dos métodos PSM e GSM: a incapacidade de controlar por efeitos não observáveis que possam influenciar tanto a probabilidade de tratamento quanto seu resultado (apenas características observáveis). Usando abordagem de primeira diferença e dados de 1999 a 2006, o autor identifica efeito positivo do FNE-Indústria sobre o crescimento do número de empregos formais nas empresas beneficiadas, com magnitude entre 10,83 e 16,11 pontos percentuais ao ano. No horizonte de três anos pós-tratamento, o valor é de 30 a 56 pontos percentuais a mais, de certa forma consistente com os resultados de \citeonline{Silva2009} e \citeonline{Soares2017}.

\citeonline{Oliveira2015} é o primeiro estudo a avaliar o efeito dos fundos constitucionais sobre variáveis agregadas do mercado de trabalho. O estudo adota os municípios de Goiás como unidade de análise e aplica o método de painel espacial autorregressivo (SAR), a fim de controlar também por características da vizinhança do município. Os resultados indicam efeitos positivos sobre emprego e salário, porém, sem diferenças relevantes em relação ao modelo de painel clássico com efeitos fixos, o que pode ser decorrência de baixa robustez dos resultados. \citeonline{Ribeiro2024} trazem duas contribuições principais: análise dos efeitos agregados do FNE por setor de atividade econômica e o uso de modelo de painel dinâmico para amenizar possível viés de endogeneidade. Os resultados apontam relação positiva entre os investimentos do FNE e o nível de emprego no município, com maior efeito sobre o setor primário (1,1\% de aumento no emprego para 1\% de aumento no investimento do FNE), seguido por terciário (0,53\%) e secundário (0,13\%). Esses resultados são observados para os investimentos correntes do FNE. Para o FNE defasado em um ano, o valor é significativo somente para o setor secundário e em menor magnitude (0,06\%). Sugere-se assim que o FNE tem efeito positivo e imediato sobre o emprego, especialmente no setor primário, setor mais intensivo em mão de obra.

Em maior nível de agregação, \citeonline{Ribeiroetal2020} aplicam um modelo dinâmico interregional de equilíbrio geral computável (CGE) para estimar os efeitos dos investimentos FNE sobre o emprego total da região Nordeste, dos estados e dos setores de atividade. As evidências apontam que os investimentos do FNE nos anos 2014 e 2015 estão associados a uma variação acumulada de 2,75\% na quantidade de empregos na região Nordeste até o ano de 2025, valor pouco menor que o estimado para a variação do PIB. Na análise por estados, Piauí, Ceará e Rio Grande do Norte também apresentam a maior variação acumulada do emprego, entre 3,19\% e 4,31\%, e Maranhão e Bahia têm as menores variações, de 0,45\% a 2,16\%. No aspecto setorial, também não há perdedores, com agropecuária (8\%) e construção (5,3\%) os setores com maior variação do emprego.

\citeonline{Oliveira2020} empregam o método diferenças-em-diferenças, que controla pela heterogeneidade não observável das empresas e possibilita uma estimativa do efeito causal do fundo. Além disso, têm como diferencial a análise dos efeitos microeconômicos do FNO por porte da empresa beneficiada, tipo de crédito concedido e horizonte de tempo. No geral, empresas beneficiadas apresentaram desempenho superior às não beneficiadas, em curto e médio prazo, em termos de geração de emprego e renda. No entanto, médias e grandes empresas beneficiadas por financiamento a capital de giro e custeio apresentam desempenho similar ao grupo de controle, no que diz respeito à geração de emprego e renda. Além disso, os autores apontam que o impacto positivo do crédito com finalidade capital de giro e custeio é em média maior no curto prazo do que no longo prazo. Por outro lado, em empresas beneficiadas pela finalidade investimento, os impactos são mais observáveis a longo prazo. Os efeitos assimétricos indicam que a finalidade investimento promove a acumulação de capital, o que aumenta a produtividade do trabalho a longo prazo, pois o impacto é positivo para o trabalho e a renda (salário médio).

\citeonline{Oliveira2021} utilizam modelo de descontinuidade geográfica (RDD) para estimar o impacto causal do FNE sobre o estoque de emprego, distinguindo-se dos estudos anteriores pela robustez metodológica e por analisar impacto sobre o emprego formal e informal, decorrência do uso de dados censitários. Os resultados indicam que o acesso ao FNE gerou impacto positivo e significativo sobre o estoque de empregos formais e informais no ano 2010, com magnitude 5,4\% (agropecuária), 2,8\% (indústria) e 3,5\% (serviços). Porém, impacto sobre a taxa de informalidade foi nulo.

Em síntese, as avaliações sobre o mercado de trabalho apresentam maior convergência do que as relativas ao PIB: há evidências consistentes de efeito positivo do FNE sobre a geração de emprego formal nas empresas beneficiadas, com magnitude crescente ao longo dos primeiros anos pós-tratamento. Os estudos também convergem em apontar efeitos mais pronunciados em micro e pequenas empresas e na modalidade de crédito para investimento. Por outro lado, a variável salário médio não apresenta variação significativa na maioria das especificações, sugerindo que os novos postos de trabalho são criados ao nível salarial vigente, sem ganhos adicionais de produtividade. Os resultados para FNO e FCO são mais limitados e inconclusivos, refletindo menor cobertura na literatura.

\subsection{Avaliações de Impacto dos Fundos de Desenvolvimento}
\label{subsec:fd}

Os Fundos de Desenvolvimento (FDNE, FDA e FDCO) constituem instrumentos mais recentes da PNDR, com primeiras contratações a partir de 2007 no caso do FDNE e de 2014 para o FDCO, o que explica a menor quantidade de avaliações empíricas em comparação aos Fundos Constitucionais. A escala dos empreendimentos financiados --- mínimo de R\$ 15 milhões, com concentração em infraestrutura e indústria --- sugere mecanismo de impacto distinto, cuja avaliação é relevante para informar a alocação de recursos entre os instrumentos da política. Dentre os estudos aprovados nesta revisão, 8 avaliam algum aspecto dos Fundos de Desenvolvimento, com concentração expressiva no FDNE. Não foram identificadas avaliações quantitativas dedicadas ao FDA nem ao FDCO, lacuna relevante dada a magnitude dos recursos envolvidos --- R\$ 5,5 bilhões e R\$ 1,8 bilhão contratados, respectivamente. As subseções seguintes organizam as evidências por abordagem metodológica. O Quadro~\ref{tab:estudos_fd} apresenta a síntese dos estudos que avaliam os Fundos de Desenvolvimento.

\input{tabelas/survey_artigos_fd}

\subsubsection{Modelos Quase Experimentais}

O método de diferenças em diferenças (DID) e suas extensões --- como o DID escalonado, adequado para intervenções implementadas em momentos distintos para diferentes unidades --- constituem a principal abordagem empregada para avaliar os Fundos de Desenvolvimento. Os três estudos discutidos a seguir representam as primeiras evidências empíricas de impacto causal desses fundos sobre indicadores municipais.

\citeonline{Braz2024} aplicam modelo DID escalonado aos municípios da área de atuação da SUDENE no período 2002--2021 e estimam que aqueles que receberam empreendimentos apoiados pelo FDNE experimentam elevação média de 24\% no PIB \textit{per capita} em relação ao grupo de controle. O estudo inova ao avaliar variáveis de resultado não estritamente econômicas: a renda domiciliar \textit{per capita} apresenta incremento de 4,6\%, e os indicadores de educação básica (IDEB) registram elevação de 19\% nos anos iniciais e 21\% nos anos finais. Não se obteve significância estatística para emprego formal, pobreza, mortalidade infantil e indicadores de saúde.

\citeonline{Carneiroetal2024a} avaliam simultaneamente os efeitos do FNE, do FDNE e dos incentivos fiscais da SUDENE sobre a economia dos municípios nordestinos, empregando estimadores DID em dois estágios e DID escalonado para o período 2002--2019. Para o FDNE, identificam impacto positivo e significativo sobre o PIB \textit{per capita} municipal, com coeficiente concentrado no setor de serviços e na arrecadação tributária. Os incentivos fiscais não apresentaram significância estatística, e a interação entre os instrumentos tampouco indicou sinergia, sugerindo que os efeitos operam de forma independente. Os resultados desse estudo para o FNE foram discutidos na Subseção~\ref{subsec:fc-pib}.

\citeonline{Irffietal2025} estimam o impacto da construção de parques eólicos sobre indicadores econômicos dos municípios nordestinos, empregando o método de controle sintético generalizado e estudos de eventos para o período 1999--2022. Os resultados apontam elevação de 19\% no PIB \textit{per capita}, de 72\% no valor adicionado bruto (VAB) industrial e de 10\% no VAB de serviços nos municípios com parques em operação ou construção. Identificam-se ainda incrementos de 23\% no emprego formal e de 31\% na massa salarial, além de redução temporária do VAB agropecuário durante a fase de construção, sem efeito permanente. Embora o estudo avalie todos os parques eólicos da região --- e não exclusivamente os financiados pelo FDNE ---, parcela expressiva dos empreendimentos de energia eólica no Nordeste contou com recursos desse fundo, e a proximidade dos valores obtidos com os de \citeonline{Braz2024} reforça a consistência das conclusões.

As três avaliações convergem em apontar impacto positivo do FDNE sobre o PIB \textit{per capita}, com magnitudes situadas entre 19\% e 24\%, consideravelmente superiores às estimadas para os Fundos Constitucionais --- cujos coeficientes situam-se tipicamente entre 0,03 e 0,13 p.p. de incremento no crescimento anual para cada 1 p.p. de aumento na relação Fundo/PIB. Essa diferença pode estar associada à escala dos empreendimentos financiados, com valor mínimo de R\$ 15 milhões, à concentração setorial em infraestrutura e indústria de transformação e ao horizonte mais longo de maturação desses investimentos. Cabe ressalvar que as três avaliações empregam variantes do método DID e analisam períodos parcialmente sobrepostos, o que limita a independência das estimativas.

\subsubsection{Modelos Espaciais}

Modelos de painel espacial permitem decompor o impacto de uma política em componente direto --- sobre o próprio município beneficiado --- e indireto --- sobre municípios vizinhos. Essa distinção é particularmente relevante para os Fundos de Desenvolvimento, cujos empreendimentos financiados (infraestrutura de energia, transportes, complexos industriais) podem gerar externalidades que se dispersam para além dos limites municipais.

\citeonline{FerreiraIrffiCarneiro2024} utilizam modelo SAC (\textit{Spatial Autoregressive Combined}) estático para avaliar o FDNE e os incentivos fiscais da SUDENE nos municípios nordestinos entre 2010 e 2021. A variável de resultado adotada é o Índice de Desenvolvimento Municipal (IDM) e suas componentes. Para o FDNE, o estudo identifica efeito positivo e significativo apenas sobre a componente IDM-renda (coeficiente 0,019), sem evidência sobre o IDM-geral, IDM-saúde ou IDM-educação. A adoção de variável de resultado distinta do PIB \textit{per capita} dificulta a comparação direta com os demais estudos desta subseção, embora contribua para ampliar o escopo das dimensões avaliadas.

\citeonline{gondim2025} analisam conjuntamente FNE, FDNE e incentivos fiscais da SUDENE por meio de modelo espacial generalizado com dados de 2010 a 2021. Os resultados para o FDNE revelam padrão distinto dos demais instrumentos: o efeito direto sobre o PIB \textit{per capita} é não significativo, enquanto o efeito indireto (transbordamento para municípios vizinhos) apresenta coeficiente positivo e significativo de 0,025, resultando em efeito total de 0,035. Para o VAB \textit{per capita}, o efeito direto é significativo (0,041), mas o indireto não. Esse resultado sugere que os investimentos do FDNE beneficiam não apenas o município-sede, mas também localidades vizinhas, padrão consistente com a natureza dos projetos financiados --- infraestrutura de energia e transportes cuja área de influência transcende limites municipais. Tal configuração contrasta com os Fundos Constitucionais, cujos efeitos diretos tendem a ser mais pronunciados.

\subsubsection{Avaliação Integrada e Síntese}

\citeonline{Souza2025} adotam painel dinâmico GMM com defasagens distribuídas para avaliar simultaneamente todos os instrumentos da PNDR --- incluindo FDNE, FDA e FDCO --- no período 2002--2021. O estudo constitui a única avaliação que contempla os três fundos de desenvolvimento em uma mesma especificação. Os resultados indicam heterogeneidade expressiva entre fundos, regiões e setores, com evidências inconclusivas para o FDCO, atribuídas ao reduzido número de projetos e ao curto período de operação.

Em conjunto, as evidências disponíveis apontam para impacto positivo expressivo do FDNE sobre o PIB \textit{per capita} municipal, com magnitudes entre 19\% e 24\% nas avaliações quase experimentais e presença de efeitos de transbordamento espacial nos modelos de painel. A base empírica, porém, apresenta limitações relevantes: concentra-se em estudos recentes (2024--2026) com períodos amostrais parcialmente sobrepostos e predominância do DID escalonado; não dispõe de avaliação quantitativa dedicada ao FDA nem ao FDCO; e restringe-se majoritariamente ao PIB \textit{per capita}, com evidências incipientes sobre emprego, renda domiciliar e pobreza. A diversificação metodológica e a extensão da cobertura para os três fundos de desenvolvimento constituem prioridades para a agenda de pesquisa.

\subsection{Avaliações de Impacto dos Incentivos Fiscais}
\label{subsec:if}

A avaliação empírica dos incentivos fiscais da SUDENE e da SUDAM constitui a vertente menos desenvolvida da literatura sobre a PNDR, a despeito do expressivo volume de renúncia tributária envolvido. Dentre os 34 estudos aprovados nesta revisão, 9 avaliam explicitamente esses instrumentos, refletindo a operação mais recente dos programas --- com concessões sistemáticas a partir de 2010 --- e as dificuldades de acesso a dados detalhados de renúncia fiscal em nível de empresa. Diferentemente dos Fundos Constitucionais, cujo mecanismo de transmissão opera via crédito subsidiado, os incentivos fiscais atuam pela redução de 75\% do Imposto de Renda de Pessoa Jurídica (IRPJ), diminuindo o custo tributário de operação em municípios da área de atuação da superintendência. Todos os estudos dedicados concentram-se nos benefícios concedidos pela SUDENE, não havendo avaliações específicas da SUDAM. Dois trabalhos avaliam adicionalmente o Prodepe --- programa estadual de Pernambuco --- em interação com os instrumentos federais. O Quadro~\ref{tab:estudos_if} apresenta a síntese dos estudos desta subseção.

\input{tabelas/survey_artigos_if}

No âmbito do mercado de trabalho, \citeonline{Garsous2017} conduzem a primeira avaliação de impacto dos incentivos fiscais da SUDENE, focalizando o setor turístico. Empregando modelo de diferenças em diferenças com \textit{matching} em municípios de Bahia, Rio de Janeiro, Espírito Santo e Minas Gerais no período 1998--2009, os autores estimam que o programa de benefícios tributários ao turismo elevou o emprego setorial em 30\% no acumulado entre 2002 e 2009. Os resultados mostram-se robustos à possibilidade de deslocamento de empregos para municípios vizinhos, embora os pesquisadores não avaliem a eficiência do programa em termos de custo fiscal por emprego gerado.

Em perspectiva mais abrangente, \citeonline{Braz2023} avaliam o efeito agregado da renúncia fiscal da SUDENE sobre o emprego formal e a renda nos municípios beneficiados, empregando estimador DID escalonado para o período 2002--2021. Os resultados indicam incremento de 3,2\% nos vínculos formais e de 1,2\% na renda nominal média, com trajetória crescente ao longo do tempo e duradoura por todo o período de vigência do benefício. Os autores destacam, porém, que os ganhos concentram-se em municípios de maior porte e dinamismo econômico, o que contrasta com o objetivo redistributivo da PNDR e levanta questões sobre a focalização territorial do instrumento.

\citeonline{Costa2024} adotam o método de controle sintético generalizado para estimar o impacto dos benefícios tributários da SUDENE sobre a criação de empregos em empresas da região Nordeste, com dados de 2006 a 2019. As evidências apontam elevação de 19,6\% nos vínculos formais no conjunto dos setores e de 18,3\% no setor manufatureiro, com persistência dos resultados por até cinco anos após o início da concessão. A magnitude consideravelmente superior à estimada por \citeonline{Braz2023} pode refletir o nível de agregação --- empresa versus município --- e a concentração no setor industrial, caracterizado por maior intensidade de mão de obra entre as atividades incentivadas.

Os três estudos convergem em apontar impacto positivo sobre o emprego formal, porém com magnitudes que variam de 3,2\% a 30\% conforme o setor analisado, o método empregado e a escala de agregação. Os ganhos mais expressivos concentram-se em segmentos específicos --- turismo e manufatura ---, enquanto o efeito agregado sobre o conjunto da economia é mais modesto. A preocupação distributiva levantada por \citeonline{Braz2023} merece atenção: se os benefícios se concentram em municípios já dinâmicos, o instrumento pode estar ampliando, em vez de reduzir, as disparidades intrarregionais.

Em abordagem complementar às avaliações de impacto, \citeonline{Carneiro2023} analisam a eficiência produtiva das empresas beneficiadas pela política de incentivos da SUDENE no período 2011--2019, combinando análise envoltória de dados (DEA) e fronteira estocástica (SFA). O estudo identifica predominância de empresas intensivas em mão de obra entre as beneficiadas e presença relevante de unidades eficientes no Semiárido. A ineficiência observada está associada mais ao setor de atividade incentivado do que às firmas individualmente, o que sugere possibilidade de aprimorar a seletividade setorial da política. Esses achados indicam que a eficiência no uso dos benefícios tributários não é uniforme e pode ser aprimorada mediante critérios diferenciados de concessão.

No que se refere ao produto e a indicadores de desenvolvimento, as evidências são mais escassas e derivam majoritariamente de estudos multi-instrumento discutidos nas subseções anteriores. \citeonline{Carneiroetal2024a}, ao avaliarem conjuntamente FNE, FDNE e incentivos fiscais, não encontram significância estatística para a renúncia fiscal da SUDENE sobre o PIB \textit{per capita} municipal nem sobre o valor adicionado setorial. Em contrapartida, \citeonline{FerreiraIrffiCarneiro2024} obtêm efeito direto significativo dos benefícios tributários sobre o IDM-geral (0,020), IDM-saúde (0,024) e IDM-renda (0,012), além de transbordamento espacial positivo sobre a componente renda nos municípios vizinhos (0,002). A divergência entre os dois estudos pode estar associada à diferença na variável de resultado e na especificação econométrica adotada.

Complementarmente, \citeonline{gondim2025} identificam efeito direto significativo da renúncia fiscal sobre o crescimento do VAB agropecuário \textit{per capita}, resultado que sugere influência do instrumento sobre setores primários não captada pelos estudos com foco no produto agregado. \citeonline{Souza2025}, ao avaliarem simultaneamente todos os instrumentos da PNDR em painel dinâmico, constituem a única referência que inclui os benefícios tributários da SUDAM, embora os resultados para esse instrumento permaneçam inconclusivos dada a escassez de observações.

Dois estudos avaliam o Programa de Desenvolvimento do Estado de Pernambuco (Prodepe), incentivo fiscal estadual que opera em complementaridade com o FNE, o FDNE e a renúncia fiscal da SUDENE. \citeonline{Oliveira2020a} empregam modelo DID com efeitos heterogêneos por tempo de exposição e dados de empresas pernambucanas no período 2000--2017, estimando incremento de 8,6\% no emprego das firmas beneficiadas. Simultaneamente, o estudo identifica redução de 10,3\% no salário médio, padrão que contrasta com a ausência de efeito salarial observada nos estudos sobre Fundos Constitucionais e sobre os incentivos federais. Não se obteve significância para a massa salarial em empresas que recebem exclusivamente o Prodepe, embora a combinação com outros instrumentos altere os resultados.

\citeonline{Alves2024} confirmam esse padrão empregando estimador DID escalonado ao mesmo recorte geográfico e temporal. As estimativas apontam elevação de 22,3\% no emprego e de 15,1\% na massa salarial, acompanhadas de redução de 8,2\% no salário médio. Os ganhos são mais pronunciados em empresas do Semiárido, de pequeno porte e do setor industrial, porém apresentam caráter temporário, com declínio progressivo ao longo dos anos de vigência do benefício. A redução salarial, consistente nos dois estudos, sugere que os incentivos estaduais atraem ou expandem atividades de menor remuneração, o que coloca em questão a qualidade dos postos de trabalho gerados.

Em síntese, as evidências sobre a renúncia fiscal revelam convergência no tocante ao emprego --- com impactos positivos variando de 3,2\% a 30\% conforme setor e método --- e divergência quanto ao produto e ao salário. Nenhum estudo identifica efeito significativo exclusivo dos incentivos fiscais sobre o PIB \textit{per capita}, embora resultados positivos sejam observados sobre indicadores multidimensionais de desenvolvimento. No mercado de trabalho, a política federal gera empregos sem alterar o salário médio, ao passo que o programa estadual eleva o emprego mas reduz a remuneração, suscitando preocupação sobre a qualidade dos postos criados. A literatura apresenta lacunas relevantes: ausência de avaliações dedicadas aos incentivos da SUDAM, inexistência de análises de custo-efetividade e predominância de métodos DID, com escassa aplicação de abordagens como RDD, PSM ou equilíbrio geral. A ampliação e diversificação metodológica desta vertente constituem prioridade para a agenda de pesquisa.

 no main.tex

A literatura empírica sobre os instrumentos da Política Nacional de Desenvolvimento Regional apresenta expressiva assimetria em termos de volume de evidências: os Fundos Constitucionais concentram a quase totalidade dos estudos, enquanto os Fundos de Desenvolvimento e os Incentivos Fiscais constituem vertentes consideravelmente menos desenvolvidas. Os resultados obtidos variam substancialmente conforme a especificação econométrica adotada, o período amostral, a escala geográfica de análise e o controle de características locais não observáveis. Esta seção apresenta os principais achados da literatura, organizados por instrumento e abordagem metodológica, buscando identificar para cada um deles convergências, divergências e condições que modulam a eficácia da política.

\subsection{Avaliações de Impacto dos Fundos Constitucionais sobre o PIB}
\label{subsec:fc-pib}

Os Fundos Constitucionais são os instrumentos da PNDR mais amplamente avaliados pela literatura empírica: dentre os 34 estudos aprovados nesta revisão, 28 investigam os efeitos desses fundos sobre variáveis de resultado relacionadas ao produto interno bruto ou ao mercado de trabalho. A questão central é se esses fundos --- que operam via crédito subsidiado de ampla capilaridade territorial --- efetivamente contribuem para o crescimento econômico dos municípios beneficiados e, em caso afirmativo, sob quais condições e com que magnitude. O FNE é o fundo com maior número de avaliações (24 estudos), seguido pelo FCO (10 estudos) e pelo FNO (8 estudos). A predominância dos Fundos Constitucionais na literatura decorre de sua maior longevidade em relação aos demais instrumentos da PNDR -- com operação contínua desde 1989 --, da disponibilidade de dados municipais anuais de PIB a partir de 2002, do elevado volume de recursos envolvidos e da capilaridade dos programas em todos os municípios das macrorregiões-alvo. A quase totalidade dos estudos adota o PIB \textit{per capita} municipal como variável de resultado, tratando-o como aproximação da renda média local, dada a escassez de alternativas em frequência anual e cobertura municipal. Os resultados obtidos pela literatura variam consideravelmente conforme a especificação econométrica adotada, o período amostral, a escala geográfica de análise e o controle de características locais não observáveis. As subseções seguintes organizam os achados por abordagem metodológica, permitindo identificar padrões gerais e divergências nas estimativas de impacto. O Quadro~\ref{tab:estudos_fc_pib} apresenta a síntese dos 21 estudos que avaliam os efeitos dos Fundos Constitucionais sobre o PIB \textit{per capita}, organizados por método de análise.

% Arquivo: latex/tabelas/tabela_macros.tex
% Macros para formatação de tabelas de resultados da revisão sistemática

% Macro principal para formatação de células com listas (label "--")
% Uso: \fc[largura]{conteúdo em itemize}
% Padrão: largura = 4.2cm (coluna Estimativa)
\newcommand{\fc}[2][4.2cm]{%
	\parbox[t]{#1}{\raggedright%
		\begin{itemize}[leftmargin=*,nosep,topsep=0pt,before=\vspace{-0.6\baselineskip},label={--},labelsep=0.4em]
			#2
		\end{itemize}%
		\vspace{0.3\baselineskip}%
	}%
}

% Macro para célula com label "--" (variáveis, amostragem)
% Uso: \fcvar[largura]{conteúdo em itemize}
% Padrão: largura = 2.2cm (colunas Amostragem e Variáveis)
\newcommand{\fcvar}[2][2.2cm]{%
	\parbox[t]{#1}{\raggedright%
		\begin{itemize}[leftmargin=*,nosep,topsep=0pt,before=\vspace{-0.6\baselineskip},label={--},labelsep=0.4em]
			#2
		\end{itemize}%
		\vspace{0.3\baselineskip}%
	}%
}

% Macro para mesclar células na coluna de método
% Parâmetros: #1 = texto, #2 = número de linhas a mesclar
\newcommand{\metodo}[2]{%
	\multirow{#2}{*}{\begin{minipage}[t]{1.4cm}\raggedright #1 \end{minipage}}%
}

% Arquivo: latex/tabelas/survey_artigos_fc_pib.tex
% Gerado automaticamente por scripts/generate_survey_tables.py
% Fonte: data/2-papers/survey_results.json

% IMPORTANTE: Este arquivo requer:
% - Pacotes: longtable, multirow, afterpage, enumitem, array
% - Macros: \fc e \metodo (definidas em tabelas/tabela_macros.tex)
% - Comando \citeonline (do pacote abntex2cite)
% - Contador: quadro (definido na classe abntex2)

\afterpage{\clearpage%
\tiny%
\setlength{\extrarowheight}{1pt}%
\setlength{\tabcolsep}{2pt}%
\linespread{0.9}\selectfont%
\renewcommand{\arraystretch}{1.2}%
\renewcommand{\tablename}{\quadroname}%
\renewcommand{\thetable}{\arabic{quadro}}%
\stepcounter{quadro}%
\begin{longtable}{|>{\raggedright\arraybackslash}p{2.2cm}
		|>{\raggedright\arraybackslash}p{2.2cm}
		|>{\raggedright\arraybackslash}p{2.3cm}
		|>{\raggedright\arraybackslash}p{3.0cm}
		|>{\raggedright\arraybackslash}p{4.2cm}|}

	\caption{Estudos de Impacto dos Fundos Constitucionais sobre o PIB \textit{per capita}} \label{tab:estudos_fc_pib} \\

\hline
	\begin{tabular}[c]{@{}>{\centering\arraybackslash}m{2.2cm}@{}}Método\end{tabular}
	& \begin{tabular}[c]{@{}>{\centering\arraybackslash}m{2.2cm}@{}}Artigo\end{tabular}
	& \begin{tabular}[c]{@{}>{\centering\arraybackslash}m{2.3cm}@{}}Amostragem\end{tabular}
	& \begin{tabular}[c]{@{}>{\centering\arraybackslash}m{3.0cm}@{}}Variáveis\end{tabular}
	& \begin{tabular}[c]{@{}>{\centering\arraybackslash}m{4.2cm}@{}}Estimativa\end{tabular} \\
	\hline
	\endfirsthead

	\multicolumn{5}{c}%
	{{\quadroname\ \thetable{} -- Continuação}} \\
	\hline
	\begin{tabular}[c]{@{}>{\centering\arraybackslash}m{2.2cm}@{}}Método\end{tabular}
	& \begin{tabular}[c]{@{}>{\centering\arraybackslash}m{2.2cm}@{}}Artigo\end{tabular}
	& \begin{tabular}[c]{@{}>{\centering\arraybackslash}m{2.3cm}@{}}Amostragem\end{tabular}
	& \begin{tabular}[c]{@{}>{\centering\arraybackslash}m{3.0cm}@{}}Variáveis\end{tabular}
	& \begin{tabular}[c]{@{}>{\centering\arraybackslash}m{4.2cm}@{}}Estimativa\end{tabular} \\
	\hline
	\endhead

	\hline
	\multicolumn{5}{r}{{Continua}} \\
	\endfoot

	\hline
	\multicolumn{5}{l}{
		\tiny
		\parbox{14.8cm}{%
			\vspace{0.1cm}
			Fonte: Elaborado pelo autor. Siglas: AC - Apresentação em congresso; TD - Trabalho para discussão; NS - Não significativo.
	}}
	\endlastfoot

Painel de Efeito Fixo
	& \citeonline{Resende2014a}
	& \fcvar[2.3cm]{\item NE \item 2004-2010}
	& \fcvar[3.0cm]{\item FC/PIB \item Cresc. PIB pc}
	& \fc[4.2cm]{
		\item FNE (munic) \hfill 0,021\textsuperscript{*}
		\item FNE (micro) \hfill 0,032\textsuperscript{*}
		\item FNE (meso) \hfill --0,288\textsuperscript{*}
		\item FNE-setorial \hfill NS\phantom{\textsuperscript{*}}
	}
	\\
	\cline{2-5}


	& \citeonline{Resende2014b}
	& \fcvar[2.3cm]{\item N \item 2004-2010}
	& \fcvar[3.0cm]{\item FC/PIB \item Cresc. PIB pc}
	& \fc[4.2cm]{
		\item FNO (munic) \hfill --0,399\textsuperscript{*}
		\item FNO (micro e meso) \hfill NS\phantom{\textsuperscript{*}}
		\item FNO setorial \hfill NS\phantom{\textsuperscript{*}}
	}
	\\
	\cline{2-5}


	& \citeonline{ResendeCravo2014}
	& \fcvar[2.3cm]{\item CO \item 2004-2010}
	& \fcvar[3.0cm]{\item FC/PIB \item Cresc. PIB pc}
	& \fc[4.2cm]{
		\item FCO (munic) \hfill 0,076\textsuperscript{*}
		\item FCO (micro e meso) \hfill NS\phantom{\textsuperscript{*}}
		\item FCO-empresarial (munic) \hfill 0,119\textsuperscript{*}
		\item FCO-rural (munic) \hfill --0,072\textsuperscript{*}
	}
	\\

	\hline
Modelo de \textit{Spillover} Espacial
	& \citeonline{Oliveira2005}
	& \fcvar[2.3cm]{\item N, CO \item 1991-2000}
	& \fcvar[3.0cm]{\item FC/PIB \item Cresc. PIB pc}
	& \fc[4.2cm]{
		\item FNO e FCO \hfill NS\phantom{\textsuperscript{*}}
	}
	\\
	\cline{2-5}


	& \citeonline{Cravo2015}
	& \fcvar[2.3cm]{\item N, NE, CO \item 2004-2010}
	& \fcvar[3.0cm]{\item FC/PIB \item Cresc. PIB pc}
	& \fc[4.2cm]{
		\item FNE (munic) \hfill NS\phantom{\textsuperscript{*}}
		\item FNO (munic) \hfill --0,313\textsuperscript{*}
		\item FCO (munic) \hfill 0,087\textsuperscript{*}
		\item FCO-empresarial (munic) \hfill +0,129\textsuperscript{*}
		\item \textit{Spillovers} espaciais \hfill NS\phantom{\textsuperscript{*}}
	}
	\\
	\cline{2-5}


	& \citeonline{Oliveira2016}
	& \fcvar[2.3cm]{\item GO \item 2004-2011}
	& \fcvar[3.0cm]{\item Log. FCO \item Log. PIB pc}
	& \fc[4.2cm]{
		\item FCO (painel espacial) \hfill 0,098\textsuperscript{*}
		\item FCO (painel clássico) \hfill 0,102\textsuperscript{*}
	}
	\\
	\cline{2-5}


	& \citeonline{Resende2017}
	& \fcvar[2.3cm]{\item NE \item 1999-2011}
	& \fcvar[3.0cm]{\item FC/PIB \item Cresc. PIB pc}
	& \fc[4.2cm]{
		\item FNE (munic. tipologia dinâmica)
		\item[\textbullet] Efeito direto \hfill 0,077\textsuperscript{*}
		\item[\textbullet] Efeito indireto \hfill 0,329\textsuperscript{*}
	}
	\\
	\cline{2-5}


	& \citeonline{Resende2018}
	& \fcvar[2.3cm]{\item N, NE, CO \item 1999-2011}
	& \fcvar[3.0cm]{\item FC/PIB \item Cresc. PIB pc}
	& \fc[4.2cm]{
		\item FC (munic. tip. dinâmica - acumulado)
		\item[\textbullet] Efeito direto \hfill 0,033\textsuperscript{*}
		\item[\textbullet] Efeito indireto \hfill 0,118\textsuperscript{*}
		\item FC (micro. tip. baixa renda - acumulado)
		\item[\textbullet] Efeito direto \hfill 0,146\textsuperscript{*}
		\item[\textbullet] Efeito indireto \hfill 0,466\textsuperscript{*}
	}
	\\
	\cline{2-5}


	& \citeonline{Silva2023}
	& \fcvar[2.3cm]{\item N, NE, CO \item 2016-2019}
	& \fcvar[3.0cm]{\item Log. FC \item Log. PIB pc}
	& \fc[4.2cm]{
		\item FNE, FNO e FCO \hfill NS\phantom{\textsuperscript{*}}
	}
	\\
	\cline{2-5}


	& \citeonline{Oliveira2026}
	& \fcvar[2.3cm]{\item NE \item 2010-2021}
	& \fcvar[3.0cm]{\item FC/PIB \item Cresc. PIB pc}
	& \fc[4.2cm]{
		\item FNE (efeito direto) \hfill 0,00001\textsuperscript{*}
		\item FNE (efeito indireto) \hfill NS\phantom{\textsuperscript{*}}
		\item FNE-serviços (efeito direto) \hfill 0,00001\textsuperscript{*}
		\item FNE-indústria (efeito direto) \hfill 0,00005\textsuperscript{*}
	}
	\\

	\hline
Modelo \textit{Threshold}
	& \citeonline{Linhares2014}
	& \fcvar[2.3cm]{\item NE \item 2000-2008}
	& \fcvar[3.0cm]{\item FC/PIB \item Cresc. PIB pc}
	& \fc[4.2cm]{
		\item FNE (PIB pc < 2.143) \hfill NS\phantom{\textsuperscript{*}}
		\item FNE (2.143 < PIB pc < 3.866) \hfill 0,078\textsuperscript{*}
		\item FNE (3.866 < PIB pc < 7.406) \hfill 0,109\textsuperscript{*}
		\item FNE (PIB pc > 7.406) \hfill NS\phantom{\textsuperscript{*}}
	}
	\\

	\hline
Primeira Diferença
	& \citeonline{Resende2014c}
	& \fcvar[2.3cm]{\item NE \item 1999-2006}
	& \fcvar[3.0cm]{\item FC/PIB \item Cresc. PIB pc}
	& \fc[4.2cm]{
		\item FNE-indústria (munic) \hfill NS\phantom{\textsuperscript{*}}
	}
	\\

	\hline
Regressão Quantílica e IV
	& \citeonline{Irffi2016}
	& \fcvar[2.3cm]{\item NE \item 2000-2010}
	& \fcvar[3.0cm]{\item FC/PIB \item Cresc. PIB pc}
	& \fc[4.2cm]{
		\item FNE \hfill (+)\textsuperscript{*}
		\item FNE-pecuária e agricultura \hfill (+)\textsuperscript{*}
	}
	\\

	\hline
Modelo de Eficiência
	& \citeonline{Carneiro2018}
	& \fcvar[2.3cm]{\item NE \item 2010-2015}
	& \fcvar[3.0cm]{\item FC/PIB \item Cresc. PIB pc}
	& \fc[4.2cm]{
		\item FNE: Maior eficiência em municípios integrados ao comércio internacional (Matopiba).
	}
	\\

	\hline
Painel Dinâmico
	& \citeonline{Cambota2019}
	& \fcvar[2.3cm]{\item NE \item 2003-2014}
	& \fcvar[3.0cm]{\item FC/PIB \item Cresc. PIB pc}
	& \fc[4.2cm]{
		\item FNE \hfill 2,96\textsuperscript{*}
		\item Efeito contemporâneo e endógeno significativo.
	}
	\\

	\hline
Equilíbrio Geral
	& \citeonline{Nascimento2017}
	& \fcvar[2.3cm]{\item BR \item 2000-2011}
	& \fcvar[3.0cm]{\item FC \item PIB}
	& \fc[4.2cm]{
		\item Remoção do FNE (alocação em despesa corrente) \hfill --0,16\%\phantom{\textsuperscript{*}}
		\item Maior concentração regional da produção
	}
	\\
	\cline{2-5}


	& \citeonline{Ribeiro2020}
	& \fcvar[2.3cm]{\item BR \item 2014-2025}
	& \fcvar[3.0cm]{\item FC/PIB \item Cresc. PIB pc}
	& \fc[4.2cm]{
		\item FNE (PIB NE) \hfill 3,51\%\textsuperscript{*}
	}
	\\
	\cline{2-5}


	& \citeonline{Goncalves2022}
	& \fcvar[2.3cm]{\item NE \item 2006-2015}
	& \fcvar[3.0cm]{\item Fundo \item Log. PIB pc}
	& \fc[4.2cm]{
		\item FNE (remoção): redução da atividade econômica no NE
	}
	\\

	\hline
RDD Geográfico
	& \citeonline{Oliveira2021}
	& \fcvar[2.3cm]{\item NE \item 1970-2010}
	& \fcvar[3.0cm]{\item Log. FNE \item Log. PIB pc}
	& \fc[4.2cm]{
		\item FNE \hfill NS\phantom{\textsuperscript{*}}
	}
	\\

	\hline
\textit{Propensity Score}
	& \citeonline{Monte2025}
	& \fcvar[2.3cm]{\item NE \item 2010-2020}
	& \fcvar[3.0cm]{\item FNE \item Log. PIB pc}
	& \fc[4.2cm]{
		\item FNE (PJ) \hfill 0,15pp\textsuperscript{*}
		\item FNE (mulheres) \hfill 0,10pp\textsuperscript{*}
		\item Efeito heterogêneo por magnitude
	}
	\\

	\hline
DID 2 Estágios, DID Escalonado
	& \citeonline{CarneiroVeloso2024}
	& \fcvar[2.3cm]{\item NE \item 2002-2019}
	& \fcvar[3.0cm]{\item FNE (\textit{dummy}) \item Log. PIB pc}
	& \fc[4.2cm]{
		\item FNE \hfill 0,08\textsuperscript{*} a 0,11\textsuperscript{*}
		\item A partir do 5º ano após o tratamento
		\item Principalmente serviços e administração pública\hfill \phantom{\textsuperscript{*}}
	}
	\\

\end{longtable}%
\normalsize%
\setlength{\extrarowheight}{0pt}%
\setlength{\tabcolsep}{6pt}%
\linespread{1.5}\selectfont%
\renewcommand{\tablename}{Tabela}%
\renewcommand{\thetable}{\arabic{table}}%
}


\subsubsection{Modelos de Efeitos Fixos}

Modelos de painel com efeitos fixos controlam características não observáveis constantes no tempo (como geografia e instituições locais) e choques temporais comuns a todas as unidades (como crises macroeconômicas), permitindo isolar de forma mais precisa o impacto da política regional. Essa abordagem foi a primeira empregada para avaliar os Fundos Constitucionais sobre o PIB municipal, aproveitando a disponibilidade de dados em painel desde 2002.

\citeonline{Resende2014a} analisa a influência do FNE sobre o crescimento do PIB \textit{per capita} dos municípios nordestinos no período 2004-2010, utilizando médias em três subperíodos (2004-2006, 2006-2008 e 2008-2010) para minimizar viés de ciclos de negócios. Os resultados apontam impacto positivo do FNE-total em nível municipal (0,021) e microrregional (0,032), mas não significativo em nível mesorregional, o que Resende (2014a) atribui à maior heterogeneidade em regiões geográficas maiores. As estimativas sugerem que o resultado positivo estaria concentrado nos empréstimos ao setor agropecuário, porém não se mostraram robustas à inclusão de controles temporais. Para o FNO, aplicando a mesma especificação econométrica à região Norte, \citeonline{Resende2014a} obtém relação inversa entre o fundo e a expansão do produto municipal, com ausência de robustez nas análises setoriais quando se controla por períodos.

\citeonline{Oliveira2017a} aplicam especificações de efeitos fixos juntamente com outros métodos (PSM e GPS) para avaliar o FCO no estado de Goiás, obtendo evidências de impacto positivo sobre a dinâmica econômica municipal, com magnitude variando conforme o método empregado.

Mais recentemente, \citeonline{Filho2024} adotam painel de efeitos fixos para avaliar conjuntamente os três fundos constitucionais, não obtendo evidências consistentes de associação positiva com o crescimento do PIB \textit{per capita} para o período 2016-2019, resultado que contrasta com os estudos anteriores.

Esses estudos ressaltam limitações metodológicas importantes. A principal dificuldade consiste na identificação causal, dado que a alocação dos fundos é endógena à demanda por crédito subsidiado, o que tende a privilegiar municípios com maior capacidade de absorção e infraestrutura socioeconômica mais desenvolvida. Adicionalmente, identificam-se dificuldades em mensurar os custos fiscais efetivos dos programas, a possibilidade de existência de peso morto (investimentos que seriam realizados mesmo sem subsídio) e potenciais deslocamentos de fatores produtivos entre municípios vizinhos. A ausência de robustez à inclusão de controles temporais sugere que parte do impacto estimado pode capturar tendências macroeconômicas comuns aos municípios tratados e não tratados, em vez do efeito causal dos fundos. Em síntese, os modelos de efeitos fixos indicam correlação positiva dos fundos constitucionais com a expansão econômica local, com impactos mais pronunciados em nível municipal e maior magnitude para o FCO, mas a fragilidade da identificação causal e a ausência de robustez em diversas especificações limitam a interpretação causal dos resultados.

\subsubsection{Modelos de Painel Espacial}

Uma vertente da literatura admite a possibilidade de interdependência econômica entre unidades espaciais geograficamente próximas, o que pode ocasionar efeitos indiretos de transbordamento da política para outras regiões, expandindo a área potencialmente beneficiada. Por outro lado, o efeito poderia consistir em vazamentos de recursos das regiões-alvo da política em direção a regiões de maior dinamismo econômico. Portanto, entender a natureza da dinâmica espacial é importante para um planejamento mais eficiente da política. \citeonline{Resende2014c} utilizam modelo de painel SDM (\textit{Spatial Durbin Model}) e estimam efeito significativo do FCO sobre o PIB \textit{per capita}, com evidência de que uma elevação de 1 ponto percentual (p.p.) na relação FCO/PIB está associada a um incremento de cerca de 0,09 p.p. na taxa média anual de crescimento do PIB \textit{per capita} municipal nos dois anos seguintes. Na análise setorial, o resultado seria atribuído à modalidade de crédito FCO-Empresarial, em que o efeito seria na ordem de 0,13 p.p. de incremento. Porém, o resultado é não significativo para o FNE e negativo para o FNO. Com respeito aos efeitos indiretos (transbordamento espacial), obteve-se resultado relevante, mas não positivo, para o FCO, o que pode estar ligado à migração de fatores de produção.

\citeonline{Oliveira2015} usam modelo espacial para avaliar efeitos do FCO nos municípios de Goiás e concluem que efeitos diretos são menores quando se admite a ocorrência de efeitos espaciais, sugerindo que modelos não espaciais tendem a superestimar o impacto real dos fundos. Entretanto, as evidências obtidas por \citeonline{Resende2014c} e \citeonline{Oliveira2015} apontam em direção contrária, com os efeitos diretos do FCO maiores em modelos espaciais do que em modelos não espaciais. Essa divergência pode estar associada a diferenças na cobertura geográfica (Goiás versus toda a região Centro-Oeste), no tipo de matriz de pesos espaciais empregada (\textit{queen} de contiguidade versus distância inversa) e no período de análise considerado. Os resultados de \citeonline{Oliveira2015} indicam que cada 1\% de aumento no valor dos financiamentos está associado a aproximadamente 0,1\% de crescimento no PIB \textit{per capita} do município, usando especificação SAR com matriz espacial \textit{queen} do tipo contiguidade e dados de 2004 a 2011.

\citeonline{ResendeSilvaFilho2017} distinguem-se dos trabalhos anteriores por incorporar na análise a heterogeneidade definida pela própria PNDR por meio das tipologias regionais. Em modelo não espacial de efeitos fixos, o efeito estimado do FNE é significativo para municípios pertencentes às tipologias dinâmica e baixa renda. Quando se admite dependência espacial, há evidências de transbordamento espacial dos investimentos do FNE quando o município diretamente beneficiado tem tipologia dinâmica. As estimativas apontam que para cada 1 p.p. de aumento na relação FNE/PIB em municípios dinâmicos há um incremento de 0,07 p.p. na taxa média anual de crescimento do PIB \textit{per capita} do próprio município nos três anos seguintes e de 0,33 p.p. na taxa média dos municípios vizinhos. Para os autores, esses resultados ressaltam a relevância de se considerar a dimensão espacial nas análises e também o caráter heterogêneo dos efeitos do FNE nos territórios. \citeonline{ResendeSilvaFilho2018} adotam a mesma abordagem econométrica, mas incluindo também as regiões Norte e Centro-Oeste, a fim de avaliar conjuntamente os fundos FNE, FNO e FCO. Os resultados são positivos para efeitos diretos e indiretos em municípios de tipologia dinâmica quando se usam valores acumulados na variável dos fundos (dois primeiros anos de cada subperíodo). Os resultados apontam que um aumento de 1 p.p. na relação FC/PIB em municípios dinâmicos está associado a 0,03 p.p. de incremento no crescimento médio do PIB \textit{per capita} do próprio município pelos três anos seguintes e 0,12 p.p. de incremento nos municípios vizinhos. \citeonline{Filho2024} não obtêm evidências de associação espacial positiva entre os financiamentos do FNE, FNO e FCO nos anos de 2016 a 2018 e o crescimento médio do PIB \textit{per capita} dos municípios de 2016 a 2019, porém, o resultado é positivo sobre o PIB \textit{per capita} de 2019. Os autores concluem que os efeitos dos fundos são em sua maior parte esporádicos e pontuais.

Em resumo, os estudos convergem para a relevância de se incorporar a dimensão espacial nas análises de efeito dos fundos constitucionais, justificada pela ocorrência de correlação espacial significativa nas variáveis envolvidas. Porém, há menor consenso quando se trata da magnitude dos efeitos. Em geral, o resultado estaria condicionado a outras características econômicas locais. O efeito direto mostrou-se significativo em municípios de tipologia dinâmica e microrregiões de baixa renda, com incrementos de 0,03 a 0,14 p.p. no crescimento anual do PIB \textit{per capita} municipal por até dois anos, para cada 1 p.p. de aumento na relação Fundo/PIB. Para o efeito indireto, as estimativas variam entre 0,1 e 0,4 p.p. de incremento no crescimento para cada 1 p.p. de aumento em Fundo/PIB, também concentrado em municípios de tipologia dinâmica e microrregiões de baixa renda. Em análise mais desagregada, evidências apontam para associação entre o FCO e taxa de crescimento, atribuída à linha de crédito empresarial. No FNE, os resultados estariam ligados ao setor industrial e de serviços. Já para o FNO não há nenhuma evidência positiva e significativa. No aspecto metodológico, nota-se a necessidade de estudos mais robustos, que combinem intervalo de tempo mais longo, especificações espaciais alternativas, correção de endogeneidade e maior número de controles (inclusive variáveis de confusão da própria PNDR).

\subsubsection{Modelos Não Lineares e de Eficiência}

\citeonline{Linharesetal2014} usam modelo \textit{threshold} e obtêm evidências de que o efeito do FNE é pouco significativo para municípios com PIB \textit{per capita} inicial abaixo de R\$ 2.143 ou acima de R\$ 7.406. Entre esses limites, os autores identificam dois grupos distintos. O primeiro é formado por municípios de PIB \textit{per capita} inicial entre R\$ 2.143 e R\$ 3.866. Nesse grupo, estima-se que 1 p.p. de aumento médio em FNE/PIB no período 2002-2006 esteja associado a um incremento de 0,08 p.p. no crescimento médio do PIB \textit{per capita} entre 2002 e 2008. No segundo grupo, de municípios com PIB \textit{per capita} inicial entre R\$ 3.866 e R\$ 7.406, o efeito estimado foi maior, na ordem de 0,11 p.p. Os resultados sugerem que o FNE tem maior efetividade em municípios com níveis intermediários de renda, reafirmando os achados de \citeonline{ResendeSilvaFilho2017} e \citeonline{ResendeSilvaFilho2018}. Em municípios de baixa renda, porém, os resultados não foram significativos, o que indica a necessidade de aprimoramento dos mecanismos da política para melhor adequação às características e necessidades dessas regiões.

\citeonline{Oliveira2015} usam modelagem quantílica para estimar a significância e heterogeneidade de efeitos do FNE ao longo da distribuição de crescimento dos municípios, empregando estimação em dois estágios e a localização de agências bancárias do BNB como instrumento para a alocação de recursos do fundo. Os efeitos foram positivos e significativos para o FNE total e para as linhas de crédito associadas à agricultura e pecuária. No entanto, os largos intervalos de confiança não permitiram constatar presença de heterogeneidade ao longo da curva de crescimento. \citeonline{Oliveira2018} analisa a relação entre a participação do FNE e a eficiência média dos municípios nordestinos. Usando modelo de análise envoltória de dados e de fronteira estocástica, o autor identifica elevada eficiência nos municípios pertencentes à região do MATOPIBA, associada à agricultura extensiva voltada à exportação. Adicionalmente, municípios do Semiárido indicaram menor eficiência no crescimento do produto médio, enquanto municípios menos densos e mais distantes das respectivas capitais e com empresas menores tendem a fazer melhor uso dos financiamentos.

\subsubsection{Modelos Dinâmicos}

\citeonline{Cambota2019} adotam modelo de painel dinâmico e instrumentos sintéticos para corrigir o viés de endogeneidade na variável PIB defasada e na variável de interesse FNE/PIB, usando dados de 2003 a 2014. Os autores confirmam a importância desse tipo de controle para a obtenção de estimativas mais realistas do impacto da política. As estimativas apontam associação positiva entre o FNE e o crescimento econômico dos municípios, na ordem de 2,96 p.p. para cada 1 p.p. de aumento em FNE/PIB no ano anterior. Cabe notar que essa magnitude é consideravelmente superior à dos demais estudos revisados, cujos coeficientes situam-se tipicamente entre 0,03 e 0,13 p.p. A discrepância pode ser atribuída à diferente especificação do modelo (painel dinâmico com instrumentos sintéticos), ao período mais longo de análise (2003-2014) ou a possível sobreestimação decorrente da proliferação de instrumentos no procedimento GMM, aspecto reconhecido na literatura como potencial fonte de viés.

\subsubsection{Modelos de Equilíbrio Geral}

Em \citeonline{Nascimento2017}, um cenário hipotético de retirada do FNE e realocação dos recursos em gastos correntes do governo resultaria a longo prazo em uma redução líquida de 0,16\% no PIB nacional e em maior concentração regional da produção. O estudo também sugere que os fundos têm papel mais relevante no crescimento de longo prazo, via elevação do estoque de capital e da capacidade produtiva local, frente à alternativa apontada. No curto prazo, ambos os cenários seriam equivalentes em termos de ativação da demanda. \citeonline{Ribeiroetal2020} indicam que os investimentos do FNE realizados em 2014 e 2015 teriam elevado o PIB da região Nordeste em 3,51\% até o ano 2025. \citeonline{Oliveira2021} constroem dois cenários alternativos envolvendo o término da política de subsídios, e ambos indicam como consequência a redução na atividade econômica na região Nordeste, puxada principalmente pela redução do produto da atividade agropecuária.

\subsubsection{Modelos Quase Experimentais}

\citeonline{Resende2014c} usa o método de primeira diferença para avaliar o efeito de empréstimos industriais do FNE sobre o crescimento do PIB \textit{per capita} municipal. Usando taxas de crescimento do PIB \textit{per capita} nos períodos 2000-2003 e 2000-2006, o autor não encontra efeito significativo das operações envolvendo o FNE-Indústria durante o ano 2000, seja em nível municipal ou microrregional. Argumenta-se que o resultado pode estar associado à elevada participação do setor agropecuário nas operações do FNE, setor em que há predomínio da informalidade, o que dificulta sua captura pelas medições do PIB.

\citeonline{Oliveira2020} adotam desenho quase-experimental que combina regressão com descontinuidade geográfica (RDD) e modelo de diferenças em diferenças. Os autores constroem um painel de dados decenal (1980, 1990, 2000 e 2010) composto por municípios dos estados de Minas Gerais e Espírito Santo que ingressaram nos critérios de elegibilidade ao FNE em períodos recentes (dados pré e pós-ingresso) e por municípios vizinhos imediatos que nunca tenham sido tratados pelo FNE. Os resultados obtidos levantam dúvidas sobre a efetividade do FNE, sugerindo que o Fundo não teve influência sobre as variáveis de produto e renda nos municípios avaliados no ano 2010. O estudo ressalta a importância de melhor vinculação do crédito ao aumento da produtividade, mediante condições diferenciadas e combinação com outros instrumentos.

Por outro lado, \citeonline{Carneiroetal2024a} obtêm evidências de efeito positivo e significativo do FNE sobre o PIB \textit{per capita} municipal. Os autores usam dados de 2002 a 2019 junto a estimadores DID dois estágios e DID escalonado, combinando três instrumentos de política regional: FNE, FDNE e incentivo fiscal da SUDENE. O estudo aponta efeito médio de tratamento de 8,8\% a 11\% no PIB \textit{per capita} do município beneficiado em relação ao grupo de controle, sendo atribuído principalmente ao setor serviços e administração pública. O estudo de eventos aponta que o efeito de tratamento é crescente e inicia no quinto ano após o primeiro tratamento. Quanto à sinergia com os demais instrumentos, o resultado é não significativo. \citeonline{MonteIrffiBastosCarneiro2025} exploram heterogeneidades de gênero, tipo de tomador e magnitude dos empréstimos (efeito dosagem) do FNE, aplicando o método \textit{Generalized Propensity Score} (GPS). O estudo encontra correlação positiva entre a participação do FNE-Mulheres no valor total do programa e a variação do PIB \textit{per capita} municipal, o que sinaliza uma eficácia maior do programa nesse grupo populacional. Efeito semelhante é obtido em relação à variável participação de tomadores pessoa jurídica no total de financiamentos. Com respeito ao efeito dosagem, os resultados apontam efeitos positivos no PIB \textit{per capita} somente na última e mais elevada dose de tratamento.

Em conjunto, as evidências sobre o efeito dos Fundos Constitucionais no PIB \textit{per capita} revelam um quadro heterogêneo e condicionado à especificação metodológica adotada. Os modelos de efeitos fixos e espaciais indicam associação positiva para o FCO e resultados ambíguos para FNE e FNO, enquanto os modelos quase experimentais apontam tanto para ausência de efeito quanto para impactos positivos significativos. Os modelos não lineares e de eficiência sugerem que a eficácia dos fundos é condicionada ao nível inicial de renda do município, com efeitos mais pronunciados em faixas intermediárias de desenvolvimento. As simulações de equilíbrio geral reafirmam a importância dos fundos para a atividade econômica regional em cenários contrafactuais de retirada da política. A constatação de que os efeitos são mais pronunciados em municípios com níveis intermediários de renda, e nulos ou inconclusivos em municípios de baixa renda, pode ser interpretada à luz da literatura sobre capacidade de absorção e armadilhas de pobreza: municípios de baixa renda podem carecer do capital humano e da infraestrutura básica necessários para converter o crédito subsidiado em crescimento econômico sustentado; municípios de renda intermediária possuiriam o mínimo de capacidade necessária para ativar o efeito multiplicador dos investimentos; e municípios de alta renda estariam próximos do efeito marginal decrescente do capital.

\subsection{Avaliações de Impacto dos Fundos Constitucionais sobre o Mercado de Trabalho}

Dezesseis estudos aprovados nesta revisão avaliam o impacto dos Fundos Constitucionais sobre variáveis do mercado de trabalho, com foco predominante no emprego formal, massa salarial e salário médio. O Quadro~\ref{tab:estudos_fc_emprego} apresenta a síntese desses estudos, organizados por método de análise.

% Arquivo: latex/tabelas/survey_artigos_fc_emprego.tex
% Gerado automaticamente por scripts/generate_survey_tables.py
% Fonte: data/2-papers/survey_results.json

% IMPORTANTE: Este arquivo requer:
% - Pacotes: longtable, multirow, afterpage, enumitem, array
% - Macros: \fc e \metodo (definidas em tabelas/tabela_macros.tex)
% - Comando \citeonline (do pacote abntex2cite)
% - Contador: quadro (definido na classe abntex2)

\afterpage{\clearpage%
\tiny%
\setlength{\extrarowheight}{1pt}%
\setlength{\tabcolsep}{2pt}%
\linespread{0.9}\selectfont%
\renewcommand{\arraystretch}{1.2}%
\renewcommand{\tablename}{\quadroname}%
\renewcommand{\thetable}{\arabic{quadro}}%
\stepcounter{quadro}%
\begin{longtable}{|>{\raggedright\arraybackslash}p{2.2cm}
		|>{\raggedright\arraybackslash}p{2.2cm}
		|>{\raggedright\arraybackslash}p{2.3cm}
		|>{\raggedright\arraybackslash}p{3.0cm}
		|>{\raggedright\arraybackslash}p{4.2cm}|}

	\caption{Estudos de Impacto dos Fundos Constitucionais sobre o Mercado de Trabalho} \label{tab:estudos_fc_emprego} \\

\hline
	\begin{tabular}[c]{@{}>{\centering\arraybackslash}m{2.2cm}@{}}Método\end{tabular}
	& \begin{tabular}[c]{@{}>{\centering\arraybackslash}m{2.2cm}@{}}Artigo\end{tabular}
	& \begin{tabular}[c]{@{}>{\centering\arraybackslash}m{2.3cm}@{}}Amostragem\end{tabular}
	& \begin{tabular}[c]{@{}>{\centering\arraybackslash}m{3.0cm}@{}}Variáveis\end{tabular}
	& \begin{tabular}[c]{@{}>{\centering\arraybackslash}m{4.2cm}@{}}Estimativa\end{tabular} \\
	\hline
	\endfirsthead

	\multicolumn{5}{c}%
	{{\quadroname\ \thetable{} -- Continuação}} \\
	\hline
	\begin{tabular}[c]{@{}>{\centering\arraybackslash}m{2.2cm}@{}}Método\end{tabular}
	& \begin{tabular}[c]{@{}>{\centering\arraybackslash}m{2.2cm}@{}}Artigo\end{tabular}
	& \begin{tabular}[c]{@{}>{\centering\arraybackslash}m{2.3cm}@{}}Amostragem\end{tabular}
	& \begin{tabular}[c]{@{}>{\centering\arraybackslash}m{3.0cm}@{}}Variáveis\end{tabular}
	& \begin{tabular}[c]{@{}>{\centering\arraybackslash}m{4.2cm}@{}}Estimativa\end{tabular} \\
	\hline
	\endhead

	\hline
	\multicolumn{5}{r}{{Continua}} \\
	\endfoot

	\hline
	\multicolumn{5}{l}{
		\tiny
		\parbox{14.8cm}{%
			\vspace{0.1cm}
			Fonte: Elaborado pelo autor. Siglas: NS - Não significativo; p.p. - Pontos percentuais.
	}}
	\endlastfoot

\textit{Propensity Score Matching}
	& \citeonline{SilvaResende2009}
	& \fcvar[2.3cm]{\item Empresas NE, N, CO \item 2000-2003}
	& \fcvar[3.0cm]{\item FNE, FNO, FCO (\textit{dummy}) \item Cresc. emprego \item Cresc. salário}
	& \fc[4.2cm]{
		\item FNE-emprego \hfill +61,2 p.p.\textsuperscript{*}
		\item FNE-salário \hfill NS\phantom{\textsuperscript{*}}
		\item FNO, FCO \hfill NS\phantom{\textsuperscript{*}}
	}
	\\
	\cline{2-5}


	& \citeonline{Soares2017}
	& \fcvar[2.3cm]{\item Empresas NE \item 2000-2006}
	& \fcvar[3.0cm]{\item FNE (\textit{dummy}) \item Cresc. emprego \item Cresc. massa salarial \item Cresc. salário}
	& \fc[4.2cm]{
		\item FNE-empr. (1 ano) \hfill +7,96 p.p.\textsuperscript{*}
		\item FNE-empr. (3 anos) \hfill +33,59 p.p.\textsuperscript{*}
		\item FNE-empr. (5 anos) \hfill +132,23 p.p.\textsuperscript{*}
		\item FNE-massa salarial \hfill (+)\textsuperscript{*}
		\item FNE-salário \hfill NS\phantom{\textsuperscript{*}}
	}
	\\

	\hline
\textit{Generalized Propensity Score}
	& \citeonline{OliveiraTerra2018}
	& \fcvar[2.3cm]{\item Empresas GO \item 2004-2011}
	& \fcvar[3.0cm]{\item FCO-Empresarial (contínuo) \item Cresc. emprego \item Cresc. salário}
	& \fc[4.2cm]{
		\item FCO-emprego \hfill (+)\textsuperscript{*}
		\item Efeito limitado a intervalo específico da dose
		\item FCO-salário \hfill NS\phantom{\textsuperscript{*}}
	}
	\\

	\hline
IV e Função Dose-Resposta
	& \citeonline{Cunha2024}
	& \fcvar[2.3cm]{\item Empresas NE \item 2000-2006}
	& \fcvar[3.0cm]{\item FNE (contínua) \item Cresc. emprego \item Cresc. massa salarial}
	& \fc[4.2cm]{
		\item Dose-resposta emprego: IV \hfill 357,54\%\textsuperscript{*}
		\item Dose-resposta emprego: OLS \hfill 133,94\%\textsuperscript{*}
		\item Dose-resposta massa: IV \hfill 220,96\%\textsuperscript{*}
		\item Dose-resposta massa: OLS \hfill 188,55\%\textsuperscript{*}
	}
	\\

	\hline
Primeira Diferença
	& \citeonline{Resende2014c}
	& \fcvar[2.3cm]{\item Empresas NE \item 1999-2006}
	& \fcvar[3.0cm]{\item FNE-Indústria (\textit{dummy}) \item Cresc. emprego}
	& \fc[4.2cm]{
		\item FNE-Ind. (anual) \hfill +10,83 a 16,11 p.p.\textsuperscript{*}
		\item FNE-Ind. (3 anos) \hfill +30 a 56 p.p.\textsuperscript{*}
	}
	\\

	\hline
Painel Espacial (SAR)
	& \citeonline{Oliveira2015}
	& \fcvar[2.3cm]{\item Municípios GO \item 2004-2011}
	& \fcvar[3.0cm]{\item Log. FCO \item Log. emprego \item Salário médio}
	& \fc[4.2cm]{
		\item FCO-emprego \hfill (+)\textsuperscript{*}
		\item FCO-salário \hfill (+)\textsuperscript{*}
		\item Sem diferença relevante vs. painel clássico
	}
	\\

	\hline
Painel Dinâmico
	& \citeonline{Ribeiro2024}
	& \fcvar[2.3cm]{\item Municípios NE \item 2010-2015}
	& \fcvar[3.0cm]{\item Log. FNE \item Log. emprego por setor}
	& \fc[4.2cm]{
		\item FNE-primário \hfill +1,10\%\textsuperscript{*}
		\item FNE-terciário \hfill +0,53\%\textsuperscript{*}
		\item FNE-secundário \hfill +0,13\%\textsuperscript{*}
		\item FNE defasado (secundário) \hfill +0,06\%\textsuperscript{*}
	}
	\\

	\hline
Equilíbrio Geral
	& \citeonline{Ribeiro2020}
	& \fcvar[2.3cm]{\item Região NE \item 2014-2025}
	& \fcvar[3.0cm]{\item FNE \item Emprego total}
	& \fc[4.2cm]{
		\item FNE (emprego NE acumulado) \hfill +2,75\%\textsuperscript{*}
		\item NE-agropecuária \hfill +8,0\%\textsuperscript{*}
		\item NE-construção \hfill +5,3\%\textsuperscript{*}
	}
	\\

	\hline
Diferenças em Diferenças
	& \citeonline{Oliveira2020}
	& \fcvar[2.3cm]{\item Empresas N (FNO) \item 2000-2010}
	& \fcvar[3.0cm]{\item FNO (\textit{dummy}) \item Emprego \item Massa salarial \item Salário médio}
	& \fc[4.2cm]{
		\item Emprego e renda \hfill (+)\textsuperscript{*}
		\item Curto e médio prazo (+), exceto grandes empresas
		\item Investimento: efeito a longo prazo
		\item Giro/custeio: efeito a curto prazo
	}
	\\

	\hline
RDD Geográfico
	& \citeonline{Oliveira2021}
	& \fcvar[2.3cm]{\item Municípios NE \item 1970-2010}
	& \fcvar[3.0cm]{\item FNE \item Emprego formal e informal (2010)}
	& \fc[4.2cm]{
		\item Agropecuária \hfill +5,4\%\textsuperscript{*}
		\item Indústria \hfill +2,8\%\textsuperscript{*}
		\item Serviços \hfill +3,5\%\textsuperscript{*}
		\item Taxa de informalidade \hfill NS\phantom{\textsuperscript{*}}
	}
	\\

\end{longtable}%
\normalsize%
\setlength{\extrarowheight}{0pt}%
\setlength{\tabcolsep}{6pt}%
\linespread{1.5}\selectfont%
\renewcommand{\tablename}{Tabela}%
\renewcommand{\thetable}{\arabic{table}}%
}


\citeonline{Silva2009} foi o primeiro estudo a investigar o impacto dos empréstimos dos FCs sobre a contratação de trabalhadores nas empresas beneficiadas. O estudo faz a interseção dos registros de empresas beneficiadas com os dados da RAIS (Relação Anual de Informações Sociais), obtendo a quantidade de vínculos de empregos formais de 449 empresas (NE=211, N=174, CO=74) nos anos de 2000 e 2003. Aplicando o método \textit{Propensity Score Matching (PSM)}, os autores concluem que nas empresas beneficiadas pelo FNE a quantidade de empregos cresceu em média 61,2 pontos percentuais a mais do que nas empresas não beneficiadas. No entanto, os valores não foram significativos para o salário médio. E ainda, para o FNO e o FCO, não foi obtido efeito positivo em nenhum dos indicadores avaliados.

\citeonline{Soares2017} também aplicam o método PSM, mas fazem uso de um conjunto mais rico de informações sobre o FNE, o que permitiu ampliar consideravelmente o número de empresas e o horizonte de tempo investigado. Além disso, os autores fazem uma análise ano a ano dos efeitos do fundo, o que permite estimativas mais eficientes, com quantidades maiores de empresas nos primeiros anos pós-tratamento. Com dados do período de 2000 a 2006, no horizonte de 3 anos pós-tratamento há 1.104 empresas e em 5 anos há 314. Os resultados obtidos reforçam as evidências obtidas por \citeonline{Silva2009}. No primeiro ano pós-tratamento o crescimento do estoque de empregos é em média 7,96 pontos percentuais maior nas empresas beneficiadas, relativamente às demais. Além disso, a trajetória é crescente nos anos subsequentes, alcançando 33,59 pontos percentuais no terceiro ano e 132,23 no quinto ano. Se confirmados, esses resultados atribuem grande relevância para a política regional, indicando que o FNE tem papel importante na geração de empregos formais nas empresas beneficiadas. Os autores também destacam que o efeito de maior monta estaria associado às empresas financiadas entre 1999 e 2001, as únicas no horizonte de 5 anos pós-tratamento, estando mais concentradas no setor industrial. O mesmo estudo ainda analisa possíveis efeitos do FNE sobre a massa salarial, obtendo resultados significativos e com dinâmica semelhante ao do efeito no emprego. Como consequência, a variável salário médio não apresentou variação significativa como decorrência do FNE, indicando que as novas contratações teriam ocorrido ao mesmo nível dos salários vigentes.

\citeonline{Oliveira2018} apresentam como inovação principal o uso de variável de tratamento do tipo contínua, permitindo analisar possíveis heterogeneidades no efeito do fundo decorrentes da intensidade do tratamento (valor do empréstimo), em contraste à abordagem dicotômica de tratamento presente nos estudos anteriores. Os autores usam o método \textit{Generalized Propensity Score Matching (GPS)} e dados de financiamentos a empresas do estado de Goiás realizados pelo FCO por meio da linha Empresarial (associada à indústria, comércio e serviços) no período de 2004 a 2011. Os resultados indicam relação positiva e significativa com o crescimento do emprego apenas em um intervalo específico da distribuição de tratamento e limitada ao período 2004-2008 (pré-crise), cenário semelhante ao obtido na abordagem PSM com variável de tratamento dicotômica. Com respeito ao salário médio, não se obteve nenhum resultado significativo. \citeonline{Oliveira2018} expandem a abordagem de \citeonline{Oliveira2018} para os demais fundos constitucionais. O dado mais relevante obtido indica efeito decrescente dos fundos sobre as variáveis de interesse, o que significa que quanto maior o valor do empréstimo, menor o efeito sobre o crescimento do emprego e do salário. Para os autores, isso se justifica do ponto de vista teórico pela propriedade de retornos marginais decrescentes da função de produção neoclássica.

\citeonline{Junior2024} contribuem com a literatura ao usar o método de função dose-resposta para captar efeitos heterogêneos e correção de viés de seleção por meio de variável instrumental. Para garantir comparabilidade, os autores optam por usar a mesma base de dados de \citeonline{Soares2017}. Na variável crescimento do emprego, o modelo IV aponta efeito médio de 357,54\%, acima de 133,94\% do modelo OLS. Para o crescimento da massa salarial, o valor médio é 220,96\% no modelo IV contra 188,55\% em OLS. Em ambas as variáveis a restrição de exogeneidade leva a subestimação do efeito médio e superestimação do ponto ótimo de tratamento, sinalizando para efeitos mais significativos em pontos mais baixos da distribuição de tratamento, quando a endogeneidade é corrigida.

\citeonline{Resende2014c} contribui para uma limitação importante dos métodos PSM e GSM: a incapacidade de controlar por efeitos não observáveis que possam influenciar tanto a probabilidade de tratamento quanto seu resultado (apenas características observáveis). Usando abordagem de primeira diferença e dados de 1999 a 2006, o autor identifica efeito positivo do FNE-Indústria sobre o crescimento do número de empregos formais nas empresas beneficiadas, com magnitude entre 10,83 e 16,11 pontos percentuais ao ano. No horizonte de três anos pós-tratamento, o valor é de 30 a 56 pontos percentuais a mais, de certa forma consistente com os resultados de \citeonline{Silva2009} e \citeonline{Soares2017}.

\citeonline{Oliveira2015} é o primeiro estudo a avaliar o efeito dos fundos constitucionais sobre variáveis agregadas do mercado de trabalho. O estudo adota os municípios de Goiás como unidade de análise e aplica o método de painel espacial autorregressivo (SAR), a fim de controlar também por características da vizinhança do município. Os resultados indicam efeitos positivos sobre emprego e salário, porém, sem diferenças relevantes em relação ao modelo de painel clássico com efeitos fixos, o que pode ser decorrência de baixa robustez dos resultados. \citeonline{Ribeiro2024} trazem duas contribuições principais: análise dos efeitos agregados do FNE por setor de atividade econômica e o uso de modelo de painel dinâmico para amenizar possível viés de endogeneidade. Os resultados apontam relação positiva entre os investimentos do FNE e o nível de emprego no município, com maior efeito sobre o setor primário (1,1\% de aumento no emprego para 1\% de aumento no investimento do FNE), seguido por terciário (0,53\%) e secundário (0,13\%). Esses resultados são observados para os investimentos correntes do FNE. Para o FNE defasado em um ano, o valor é significativo somente para o setor secundário e em menor magnitude (0,06\%). Sugere-se assim que o FNE tem efeito positivo e imediato sobre o emprego, especialmente no setor primário, setor mais intensivo em mão de obra.

Em maior nível de agregação, \citeonline{Ribeiroetal2020} aplicam um modelo dinâmico interregional de equilíbrio geral computável (CGE) para estimar os efeitos dos investimentos FNE sobre o emprego total da região Nordeste, dos estados e dos setores de atividade. As evidências apontam que os investimentos do FNE nos anos 2014 e 2015 estão associados a uma variação acumulada de 2,75\% na quantidade de empregos na região Nordeste até o ano de 2025, valor pouco menor que o estimado para a variação do PIB. Na análise por estados, Piauí, Ceará e Rio Grande do Norte também apresentam a maior variação acumulada do emprego, entre 3,19\% e 4,31\%, e Maranhão e Bahia têm as menores variações, de 0,45\% a 2,16\%. No aspecto setorial, também não há perdedores, com agropecuária (8\%) e construção (5,3\%) os setores com maior variação do emprego.

\citeonline{Oliveira2020} empregam o método diferenças-em-diferenças, que controla pela heterogeneidade não observável das empresas e possibilita uma estimativa do efeito causal do fundo. Além disso, têm como diferencial a análise dos efeitos microeconômicos do FNO por porte da empresa beneficiada, tipo de crédito concedido e horizonte de tempo. No geral, empresas beneficiadas apresentaram desempenho superior às não beneficiadas, em curto e médio prazo, em termos de geração de emprego e renda. No entanto, médias e grandes empresas beneficiadas por financiamento a capital de giro e custeio apresentam desempenho similar ao grupo de controle, no que diz respeito à geração de emprego e renda. Além disso, os autores apontam que o impacto positivo do crédito com finalidade capital de giro e custeio é em média maior no curto prazo do que no longo prazo. Por outro lado, em empresas beneficiadas pela finalidade investimento, os impactos são mais observáveis a longo prazo. Os efeitos assimétricos indicam que a finalidade investimento promove a acumulação de capital, o que aumenta a produtividade do trabalho a longo prazo, pois o impacto é positivo para o trabalho e a renda (salário médio).

\citeonline{Oliveira2021} utilizam modelo de descontinuidade geográfica (RDD) para estimar o impacto causal do FNE sobre o estoque de emprego, distinguindo-se dos estudos anteriores pela robustez metodológica e por analisar impacto sobre o emprego formal e informal, decorrência do uso de dados censitários. Os resultados indicam que o acesso ao FNE gerou impacto positivo e significativo sobre o estoque de empregos formais e informais no ano 2010, com magnitude 5,4\% (agropecuária), 2,8\% (indústria) e 3,5\% (serviços). Porém, impacto sobre a taxa de informalidade foi nulo.

Em síntese, as avaliações sobre o mercado de trabalho apresentam maior convergência do que as relativas ao PIB: há evidências consistentes de efeito positivo do FNE sobre a geração de emprego formal nas empresas beneficiadas, com magnitude crescente ao longo dos primeiros anos pós-tratamento. Os estudos também convergem em apontar efeitos mais pronunciados em micro e pequenas empresas e na modalidade de crédito para investimento. Por outro lado, a variável salário médio não apresenta variação significativa na maioria das especificações, sugerindo que os novos postos de trabalho são criados ao nível salarial vigente, sem ganhos adicionais de produtividade. Os resultados para FNO e FCO são mais limitados e inconclusivos, refletindo menor cobertura na literatura.

\subsection{Avaliações de Impacto dos Fundos de Desenvolvimento}
\label{subsec:fd}

Os Fundos de Desenvolvimento (FDNE, FDA e FDCO) constituem instrumentos mais recentes da PNDR, com primeiras contratações a partir de 2007 no caso do FDNE e de 2014 para o FDCO, o que explica a menor quantidade de avaliações empíricas em comparação aos Fundos Constitucionais. A escala dos empreendimentos financiados --- mínimo de R\$ 15 milhões, com concentração em infraestrutura e indústria --- sugere mecanismo de impacto distinto, cuja avaliação é relevante para informar a alocação de recursos entre os instrumentos da política. Dentre os estudos aprovados nesta revisão, 8 avaliam algum aspecto dos Fundos de Desenvolvimento, com concentração expressiva no FDNE. Não foram identificadas avaliações quantitativas dedicadas ao FDA nem ao FDCO, lacuna relevante dada a magnitude dos recursos envolvidos --- R\$ 5,5 bilhões e R\$ 1,8 bilhão contratados, respectivamente. As subseções seguintes organizam as evidências por abordagem metodológica. O Quadro~\ref{tab:estudos_fd} apresenta a síntese dos estudos que avaliam os Fundos de Desenvolvimento.

% Arquivo: latex/tabelas/survey_artigos_fd.tex
% Quadro de estudos sobre impacto dos Fundos de Desenvolvimento
% Fonte: Revisão sistemática - dados extraídos do pipeline pndr_survey

% IMPORTANTE: Este arquivo requer:
% - Pacotes: longtable, multirow, afterpage, enumitem, array
% - Macros: \fc e \metodo (definidas em tabelas/tabela_macros.tex)
% - Comando \citeonline (do pacote abntex2cite)
% - Contador: quadro (definido no main.tex)

\afterpage{\clearpage%
\tiny%
\setlength{\extrarowheight}{1pt}%
\setlength{\tabcolsep}{2pt}%
\linespread{0.9}\selectfont%
\renewcommand{\arraystretch}{1.2}%
\renewcommand{\tablename}{\quadroname}%
\renewcommand{\thetable}{\arabic{quadro}}%
\stepcounter{quadro}%
\begin{longtable}{|>{\raggedright\arraybackslash}p{1.6cm}
		|>{\raggedright\arraybackslash}p{1.7cm}
		|>{\raggedright\arraybackslash}p{2cm}
		|>{\raggedright\arraybackslash}p{1.5cm}
		|>{\raggedright\arraybackslash}p{2.6cm}
		|>{\raggedright\arraybackslash}p{4.2cm}|}

	\caption{Estudos de Impacto dos Fundos de Desenvolvimento} \label{tab:estudos_fd} \\

\hline
	\begin{tabular}[c]{@{}>{\centering\arraybackslash}m{1.6cm}@{}}Método\end{tabular}
	& \begin{tabular}[c]{@{}>{\centering\arraybackslash}m{1.7cm}@{}}Artigo\end{tabular}
	& \begin{tabular}[c]{@{}>{\centering\arraybackslash}m{2cm}@{}}Amostragem\end{tabular}
	& \begin{tabular}[c]{@{}>{\centering\arraybackslash}m{1.5cm}@{}}Variável Independente\end{tabular}
	& \begin{tabular}[c]{@{}>{\centering\arraybackslash}m{2.6cm}@{}}Variável Dependente\end{tabular}
	& \begin{tabular}[c]{@{}>{\centering\arraybackslash}m{4.2cm}@{}}Resultado\end{tabular} \\
	\hline
	\endfirsthead

	\multicolumn{6}{c}%
	{{\quadroname\ \thetable{} -- Continuação}} \\
	\hline
	\begin{tabular}[c]{@{}>{\centering\arraybackslash}m{1.6cm}@{}}Método\end{tabular}
	& \begin{tabular}[c]{@{}>{\centering\arraybackslash}m{1.7cm}@{}}Artigo\end{tabular}
	& \begin{tabular}[c]{@{}>{\centering\arraybackslash}m{2cm}@{}}Amostragem\end{tabular}
	& \begin{tabular}[c]{@{}>{\centering\arraybackslash}m{1.5cm}@{}}Variável Independente\end{tabular}
	& \begin{tabular}[c]{@{}>{\centering\arraybackslash}m{2.6cm}@{}}Variável Dependente\end{tabular}
	& \begin{tabular}[c]{@{}>{\centering\arraybackslash}m{4.2cm}@{}}Resultado\end{tabular} \\
	\hline
	\endhead

	\hline
	\multicolumn{6}{r}{{Continua}} \\
	\endfoot

	\hline
	\multicolumn{6}{l}{
		\tiny
		\parbox{14.8cm}{%
			\vspace{0.1cm}
			Fonte: Elaborado pelo autor. Siglas: AC - Apresentação em congresso; NS - Não significativo.
	}}
	\endlastfoot

	\hline
Painel Espacial & \citeonline{FerreiraIrffiCarneiro2024}
	& \fc[2cm]{\item NE \item 2010-2021}
	& \fc[1.5cm]{\item FDNE}
	& \fc[2cm]{\item IDM}
	& \fc[4.5cm]{
		\item IDM renda \hfill 0,019\textsuperscript{*}
		\item IDM, IDM saúde e IDM educação \hfill NS\phantom{\textsuperscript{*}}
	} \\

	\cline{2-6}

	& \citeonline{gondim2025}
	& \fc[2cm]{\item NE \item 2010-2021}
	& \fc[1.5cm]{\item FDNE (\textit{dummy})}
	& \fc[2cm]{\item $\Delta$ ln(PIB pc)}
	& \fc[4.5cm]{
		\item Efeito direto \hfill NS\phantom{\textsuperscript{*}}
		\item Efeito indireto \hfill 0,0252\textsuperscript{*}
		\item Efeito total \hfill 0,0351\textsuperscript{*}
	} \\
	\cline{5-6}
	& & &
	& \fc[2cm]{\item $\Delta$ ln(VAB pc)}
	& \fc[4.5cm]{
		\item Efeito direto \hfill 0,0409\textsuperscript{*}
		\item Efeito indireto \hfill NS\phantom{\textsuperscript{*}}
		\item Efeito total \hfill 0,6014\textsuperscript{*}
	} \\

	\hline
DID Dois Estágios & \citeonline{Carneiroetal2024a}
	& \fc[2cm]{\item NE \item 2002-2019}
	& \fc[1.5cm]{\item FDNE}
	& \fc[2cm]{\item log PIB\textit{pc}, VAB setorial}
	& \fc[4.5cm]{
		\item PIB\textit{pc} \hfill 0,063\textsuperscript{*}
		\item VAB\textit{pc} serviços \hfill 0,111\textsuperscript{*}
		\item Arrecadação \hfill 0,207\textsuperscript{*}
		\item VAB ind. e VAB agro. \hfill NS\phantom{\textsuperscript{*}}
	} \\

	\hline
DID Escalonado
	& \citeonline{BrazBastosIrffi2024}
	& \fc[2cm]{\item NE \item 2002-2021}
	& \fc[1.5cm]{\item FDNE}
	& \fc[2cm]{\item PIB\textit{pc}, emprego, renda, IDEB}
	& \fc[4.5cm]{
		\item PIB\textit{pc} \hfill 24,07\%\textsuperscript{*}
		\item Renda \hfill 4,63\%\textsuperscript{*}
		\item IDEB anos iniciais \hfill 18,67\%\textsuperscript{*}
		\item IDEB anos finais \hfill 21,49\%\textsuperscript{*}
		\item Emprego, pobreza, mortalidade infantil e saúde \hfill NS\phantom{\textsuperscript{*}}
	} \\
	\cline{2-6}

	& \citeonline{Irffietal2025}
	& \fc[2cm]{\item NE \item 1999-2022}
	& \fc[1.5cm]{\item FDNE-parques eólicos}
	& \fc[2cm]{\item PIB\textit{pc}, VAB, emprego, massa salarial}
	& \fc[4.5cm]{
		\item PIB\textit{pc} \hfill 18,97\%\textsuperscript{*}
		\item VAB indústria \hfill 71,90\%\textsuperscript{*}
		\item VAB serviço \hfill 10,24\%\textsuperscript{*}
		\item Emprego \hfill 23,41\%\textsuperscript{*}
		\item Redução temporária do VAB-agro durante construção
	} \\

	\hline
Painel Dinâmico GMM
	& \citeonline{Souza2025}
	& \fc[2cm]{\item N, NE, CO \item 2002-2021}
	& \fc[1.5cm]{\item FDNE, FDA, FDCO (+ FC + IF)}
	& \fc[2cm]{\item PIB\textit{pc}, emprego, renda}
	& \fc[4.5cm]{
		\item Heterogeneidade entre fundos, regiões e setores
		\item FDCO inconclusivo (poucos projetos)
	} \\

\end{longtable}%
\normalsize%
\setlength{\extrarowheight}{0pt}%
\setlength{\tabcolsep}{6pt}%
\linespread{1.5}\selectfont%
\renewcommand{\tablename}{Tabela}%
\renewcommand{\thetable}{\arabic{table}}%
}


\subsubsection{Modelos Quase Experimentais}

O método de diferenças em diferenças (DID) e suas extensões --- como o DID escalonado, adequado para intervenções implementadas em momentos distintos para diferentes unidades --- constituem a principal abordagem empregada para avaliar os Fundos de Desenvolvimento. Os três estudos discutidos a seguir representam as primeiras evidências empíricas de impacto causal desses fundos sobre indicadores municipais.

\citeonline{Braz2024} aplicam modelo DID escalonado aos municípios da área de atuação da SUDENE no período 2002--2021 e estimam que aqueles que receberam empreendimentos apoiados pelo FDNE experimentam elevação média de 24\% no PIB \textit{per capita} em relação ao grupo de controle. O estudo inova ao avaliar variáveis de resultado não estritamente econômicas: a renda domiciliar \textit{per capita} apresenta incremento de 4,6\%, e os indicadores de educação básica (IDEB) registram elevação de 19\% nos anos iniciais e 21\% nos anos finais. Não se obteve significância estatística para emprego formal, pobreza, mortalidade infantil e indicadores de saúde.

\citeonline{Carneiroetal2024a} avaliam simultaneamente os efeitos do FNE, do FDNE e dos incentivos fiscais da SUDENE sobre a economia dos municípios nordestinos, empregando estimadores DID em dois estágios e DID escalonado para o período 2002--2019. Para o FDNE, identificam impacto positivo e significativo sobre o PIB \textit{per capita} municipal, com coeficiente concentrado no setor de serviços e na arrecadação tributária. Os incentivos fiscais não apresentaram significância estatística, e a interação entre os instrumentos tampouco indicou sinergia, sugerindo que os efeitos operam de forma independente. Os resultados desse estudo para o FNE foram discutidos na Subseção~\ref{subsec:fc-pib}.

\citeonline{Irffietal2025} estimam o impacto da construção de parques eólicos sobre indicadores econômicos dos municípios nordestinos, empregando o método de controle sintético generalizado e estudos de eventos para o período 1999--2022. Os resultados apontam elevação de 19\% no PIB \textit{per capita}, de 72\% no valor adicionado bruto (VAB) industrial e de 10\% no VAB de serviços nos municípios com parques em operação ou construção. Identificam-se ainda incrementos de 23\% no emprego formal e de 31\% na massa salarial, além de redução temporária do VAB agropecuário durante a fase de construção, sem efeito permanente. Embora o estudo avalie todos os parques eólicos da região --- e não exclusivamente os financiados pelo FDNE ---, parcela expressiva dos empreendimentos de energia eólica no Nordeste contou com recursos desse fundo, e a proximidade dos valores obtidos com os de \citeonline{Braz2024} reforça a consistência das conclusões.

As três avaliações convergem em apontar impacto positivo do FDNE sobre o PIB \textit{per capita}, com magnitudes situadas entre 19\% e 24\%, consideravelmente superiores às estimadas para os Fundos Constitucionais --- cujos coeficientes situam-se tipicamente entre 0,03 e 0,13 p.p. de incremento no crescimento anual para cada 1 p.p. de aumento na relação Fundo/PIB. Essa diferença pode estar associada à escala dos empreendimentos financiados, com valor mínimo de R\$ 15 milhões, à concentração setorial em infraestrutura e indústria de transformação e ao horizonte mais longo de maturação desses investimentos. Cabe ressalvar que as três avaliações empregam variantes do método DID e analisam períodos parcialmente sobrepostos, o que limita a independência das estimativas.

\subsubsection{Modelos Espaciais}

Modelos de painel espacial permitem decompor o impacto de uma política em componente direto --- sobre o próprio município beneficiado --- e indireto --- sobre municípios vizinhos. Essa distinção é particularmente relevante para os Fundos de Desenvolvimento, cujos empreendimentos financiados (infraestrutura de energia, transportes, complexos industriais) podem gerar externalidades que se dispersam para além dos limites municipais.

\citeonline{FerreiraIrffiCarneiro2024} utilizam modelo SAC (\textit{Spatial Autoregressive Combined}) estático para avaliar o FDNE e os incentivos fiscais da SUDENE nos municípios nordestinos entre 2010 e 2021. A variável de resultado adotada é o Índice de Desenvolvimento Municipal (IDM) e suas componentes. Para o FDNE, o estudo identifica efeito positivo e significativo apenas sobre a componente IDM-renda (coeficiente 0,019), sem evidência sobre o IDM-geral, IDM-saúde ou IDM-educação. A adoção de variável de resultado distinta do PIB \textit{per capita} dificulta a comparação direta com os demais estudos desta subseção, embora contribua para ampliar o escopo das dimensões avaliadas.

\citeonline{gondim2025} analisam conjuntamente FNE, FDNE e incentivos fiscais da SUDENE por meio de modelo espacial generalizado com dados de 2010 a 2021. Os resultados para o FDNE revelam padrão distinto dos demais instrumentos: o efeito direto sobre o PIB \textit{per capita} é não significativo, enquanto o efeito indireto (transbordamento para municípios vizinhos) apresenta coeficiente positivo e significativo de 0,025, resultando em efeito total de 0,035. Para o VAB \textit{per capita}, o efeito direto é significativo (0,041), mas o indireto não. Esse resultado sugere que os investimentos do FDNE beneficiam não apenas o município-sede, mas também localidades vizinhas, padrão consistente com a natureza dos projetos financiados --- infraestrutura de energia e transportes cuja área de influência transcende limites municipais. Tal configuração contrasta com os Fundos Constitucionais, cujos efeitos diretos tendem a ser mais pronunciados.

\subsubsection{Avaliação Integrada e Síntese}

\citeonline{Souza2025} adotam painel dinâmico GMM com defasagens distribuídas para avaliar simultaneamente todos os instrumentos da PNDR --- incluindo FDNE, FDA e FDCO --- no período 2002--2021. O estudo constitui a única avaliação que contempla os três fundos de desenvolvimento em uma mesma especificação. Os resultados indicam heterogeneidade expressiva entre fundos, regiões e setores, com evidências inconclusivas para o FDCO, atribuídas ao reduzido número de projetos e ao curto período de operação.

Em conjunto, as evidências disponíveis apontam para impacto positivo expressivo do FDNE sobre o PIB \textit{per capita} municipal, com magnitudes entre 19\% e 24\% nas avaliações quase experimentais e presença de efeitos de transbordamento espacial nos modelos de painel. A base empírica, porém, apresenta limitações relevantes: concentra-se em estudos recentes (2024--2026) com períodos amostrais parcialmente sobrepostos e predominância do DID escalonado; não dispõe de avaliação quantitativa dedicada ao FDA nem ao FDCO; e restringe-se majoritariamente ao PIB \textit{per capita}, com evidências incipientes sobre emprego, renda domiciliar e pobreza. A diversificação metodológica e a extensão da cobertura para os três fundos de desenvolvimento constituem prioridades para a agenda de pesquisa.

\subsection{Avaliações de Impacto dos Incentivos Fiscais}
\label{subsec:if}

A avaliação empírica dos incentivos fiscais da SUDENE e da SUDAM constitui a vertente menos desenvolvida da literatura sobre a PNDR, a despeito do expressivo volume de renúncia tributária envolvido. Dentre os 34 estudos aprovados nesta revisão, 9 avaliam explicitamente esses instrumentos, refletindo a operação mais recente dos programas --- com concessões sistemáticas a partir de 2010 --- e as dificuldades de acesso a dados detalhados de renúncia fiscal em nível de empresa. Diferentemente dos Fundos Constitucionais, cujo mecanismo de transmissão opera via crédito subsidiado, os incentivos fiscais atuam pela redução de 75\% do Imposto de Renda de Pessoa Jurídica (IRPJ), diminuindo o custo tributário de operação em municípios da área de atuação da superintendência. Todos os estudos dedicados concentram-se nos benefícios concedidos pela SUDENE, não havendo avaliações específicas da SUDAM. Dois trabalhos avaliam adicionalmente o Prodepe --- programa estadual de Pernambuco --- em interação com os instrumentos federais. O Quadro~\ref{tab:estudos_if} apresenta a síntese dos estudos desta subseção.

% Arquivo: latex/tabelas/survey_artigos_if.tex
% Quadro de estudos sobre impacto dos Incentivos Fiscais
% Fonte: Revisão sistemática - dados extraídos do pipeline pndr_survey
% Adaptado da tese: removido CostaCarneiroIrffi2023 (não aprovado),
%   adicionados Oliveira2020a, Alves2024, Gondim2025, Souza2025.
%   Chaves BibTeX corrigidas para consistência com references.bib.

% IMPORTANTE: Este arquivo requer:
% - Pacotes: longtable, multirow, afterpage, enumitem, array
% - Macros: \fc e \metodo (definidas em tabelas/tabela_macros.tex)
% - Comando \citeonline (do pacote abntex2cite)
% - Contador: quadro (definido no main.tex)

\afterpage{\clearpage%
\tiny%
\setlength{\extrarowheight}{1pt}%
\setlength{\tabcolsep}{2pt}%
\linespread{0.9}\selectfont%
\renewcommand{\arraystretch}{1.2}%
\renewcommand{\tablename}{\quadroname}%
\renewcommand{\thetable}{\arabic{quadro}}%
\stepcounter{quadro}%
\begin{longtable}{|>{\raggedright\arraybackslash}p{1.4cm}
		|>{\raggedright\arraybackslash}p{1.5cm}
		|>{\raggedright\arraybackslash}p{1.5cm}
		|>{\raggedright\arraybackslash}p{2cm}
		|>{\raggedright\arraybackslash}p{1.5cm}
		|>{\raggedright\arraybackslash}p{2cm}
		|>{\raggedright\arraybackslash}p{4.2cm}|}

	\caption{Estudos de Impacto dos Incentivos Fiscais} \label{tab:estudos_if} \\

	\hline
	\begin{tabular}[c]{@{}>{\centering\arraybackslash}m{1.4cm}@{}}Método\end{tabular}
	& \begin{tabular}[c]{@{}>{\centering\arraybackslash}m{1.5cm}@{}}Artigo\end{tabular}
	& \begin{tabular}[c]{@{}>{\centering\arraybackslash}m{1.5cm}@{}}Publicação\end{tabular}
	& \begin{tabular}[c]{@{}>{\centering\arraybackslash}m{2cm}@{}}Amostragem\end{tabular}
	& \begin{tabular}[c]{@{}>{\centering\arraybackslash}m{1.5cm}@{}}Variável Independente\end{tabular}
	& \begin{tabular}[c]{@{}>{\centering\arraybackslash}m{2cm}@{}}Variável Dependente\end{tabular}
	& \begin{tabular}[c]{@{}>{\centering\arraybackslash}m{4.2cm}@{}}Resultado\end{tabular} \\
	\hline
	\endfirsthead

	\multicolumn{7}{c}%
	{{\quadroname\ \thetable{} -- Continuação}} \\
	\hline
	\begin{tabular}[c]{@{}>{\centering\arraybackslash}m{1.4cm}@{}}Método\end{tabular}
	& \begin{tabular}[c]{@{}>{\centering\arraybackslash}m{1.5cm}@{}}Artigo\end{tabular}
	& \begin{tabular}[c]{@{}>{\centering\arraybackslash}m{1.5cm}@{}}Publicação\end{tabular}
	& \begin{tabular}[c]{@{}>{\centering\arraybackslash}m{2cm}@{}}Amostragem\end{tabular}
	& \begin{tabular}[c]{@{}>{\centering\arraybackslash}m{1.5cm}@{}}Variável Independente\end{tabular}
	& \begin{tabular}[c]{@{}>{\centering\arraybackslash}m{2cm}@{}}Variável Dependente\end{tabular}
	& \begin{tabular}[c]{@{}>{\centering\arraybackslash}m{4.2cm}@{}}Resultado\end{tabular} \\
	\hline
	\endhead

	\hline
	\multicolumn{7}{r}{{Continua}} \\
	\endfoot

	\hline
	\multicolumn{7}{l}{
		\tiny
		\parbox{14.8cm}{%
			\vspace{0.1cm}
			Fonte: Elaborado pelo autor. Siglas: AC - Apresentação em congresso; NS - Não significativo.
	}}
	\endlastfoot

	\hline
DID & \citeonline{Garsous2017}
	& \textit{World Development}
	& \fc[2cm]{\item BA, RJ, ES, MG \item 1998-2009}
	& \fc[1.5cm]{\item IF-SUDENE (turismo)}
	& \fc[2cm]{\item Emprego no turismo (log)}
	& \fc[4.2cm]{
		\item Emprego turismo \hfill 20\%\textsuperscript{*}, 30\%\textsuperscript{*}
	} \\

	\cline{2-7}
	& \citeonline{Oliveira2020a}
	& AC
	& \fc[2cm]{\item PE \item 2000-2017}
	& \fc[1.5cm]{\item Prodepe + IF-SUDENE}
	& \fc[2cm]{\item Emprego, salário médio, massa salarial}
	& \fc[4.2cm]{
		\item Emprego \hfill 8,6\%\textsuperscript{*}
		\item Salário médio \hfill $-$10,3\%\textsuperscript{*}
		\item Massa salarial \hfill NS\phantom{\textsuperscript{*}}
	} \\

	\hline
DID Escalonado & \citeonline{Braz2023}
	& AC
	& \fc[2cm]{\item SUDENE \item 2002-2021}
	& \fc[1.5cm]{\item IF-SUDENE}
	& \fc[2cm]{\item Emprego, renda}
	& \fc[4.2cm]{
		\item Emprego \hfill 3,2\%\textsuperscript{*}
		\item Renda \hfill 1,2\%\textsuperscript{*}
		\item Efeito crescente e duradouro
	} \\

	\cline{2-7}
	& \citeonline{Alves2024}
	& AC
	& \fc[2cm]{\item PE \item 2000-2017}
	& \fc[1.5cm]{\item Prodepe + IF-SUDENE}
	& \fc[2cm]{\item Emprego, salário médio, massa salarial}
	& \fc[4.2cm]{
		\item Emprego \hfill 22,3\%\textsuperscript{*}
		\item Massa salarial \hfill 15,1\%\textsuperscript{*}
		\item Salário médio \hfill $-$8,2\%\textsuperscript{*}
		\item Efeito temporário e declinante
	} \\

	\hline
DEA + SFA & \citeonline{Carneiro2023a}
	& Rev. Cadernos de Finanças Públicas
	& \fc[2cm]{\item NE \item 2011-2019}
	& \fc[1.5cm]{\item IF-SUDENE}
	& \fc[2cm]{\item Eficiência produtiva (massa salarial)}
	& \fc[4.2cm]{
		\item Empresas eficientes no Semiárido
		\item Intensivas em mão de obra
		\item Ineficiência associada ao setor, não à empresa
	} \\

	\hline
DID Dois Estágios & \citeonline{Carneiroetal2024a}
	& AC
	& \fc[2cm]{\item NE, MG, ES \item 2002-2019}
	& \fc[1.5cm]{\item IF-SUDENE (+ FNE, FDNE)}
	& \fc[2cm]{\item log PIB\textit{pc}, VAB setorial}
	& \fc[4.2cm]{
		\item IF-SUDENE \hfill NS\phantom{\textsuperscript{*}}
	} \\

	\hline
Controle Sintético & \citeonline{Costa2024}
	& \textit{Applied Economics Letters}
	& \fc[2cm]{\item NE \item 2006-2019}
	& \fc[1.5cm]{\item IF-SUDENE}
	& \fc[2cm]{\item Vínculos (log)}
	& \fc[4.2cm]{
		\item Todos os setores \hfill 19,6\%\textsuperscript{*}
		\item Indústria \hfill 18,3\%\textsuperscript{*}
		\item Efeitos por até 5 anos
		\item Demais setores \hfill NS\phantom{\textsuperscript{*}}
	} \\

	\hline
Painel Espacial & \citeonline{FerreiraIrffiCarneiro2024}
	& AC
	& \fc[2cm]{\item NE \item 2010-2021}
	& \fc[1.5cm]{\item IF-SUDENE (+ FDNE)}
	& \fc[2cm]{\item IDM e subcomponentes}
	& \fc[4.2cm]{
		\item IDM \hfill 0,020\textsuperscript{*}
		\item IDM saúde \hfill 0,024\textsuperscript{*}
		\item IDM renda \hfill 0,012\textsuperscript{*}
		\item IDM educação \hfill NS\phantom{\textsuperscript{*}}
		\item Transbordamento renda \hfill 0,002\textsuperscript{*}
	} \\

	\cline{2-7}
	& \citeonline{Gondim2025}
	& Rev. Cadernos de Finanças Públicas
	& \fc[2cm]{\item NE \item 2010-2021}
	& \fc[1.5cm]{\item IF-SUDENE (+ FNE, FDNE)}
	& \fc[2cm]{\item $\Delta$ ln(PIB\textit{pc}), $\Delta$ ln(VAB\textit{pc})}
	& \fc[4.2cm]{
		\item VAB agropecuário\textit{pc} \hfill Significativo\textsuperscript{*}
		\item PIB\textit{pc} \hfill NS\phantom{\textsuperscript{*}}
	} \\

	\hline
Painel Dinâmico GMM & \citeonline{Souza2025}
	& AC
	& \fc[2cm]{\item N, NE, CO \item 2002-2021}
	& \fc[1.5cm]{\item IF-SUDENE, IF-SUDAM (+ FC, FD)}
	& \fc[2cm]{\item PIB\textit{pc}, emprego, renda}
	& \fc[4.2cm]{
		\item IF-SUDAM inconclusivo (poucas observações)
		\item Heterogeneidade entre instrumentos e regiões
	} \\

\end{longtable}%
\normalsize%
\setlength{\extrarowheight}{0pt}%
\setlength{\tabcolsep}{6pt}%
\linespread{1.5}\selectfont%
\renewcommand{\tablename}{Tabela}%
\renewcommand{\thetable}{\arabic{table}}%
}


No âmbito do mercado de trabalho, \citeonline{Garsous2017} conduzem a primeira avaliação de impacto dos incentivos fiscais da SUDENE, focalizando o setor turístico. Empregando modelo de diferenças em diferenças com \textit{matching} em municípios de Bahia, Rio de Janeiro, Espírito Santo e Minas Gerais no período 1998--2009, os autores estimam que o programa de benefícios tributários ao turismo elevou o emprego setorial em 30\% no acumulado entre 2002 e 2009. Os resultados mostram-se robustos à possibilidade de deslocamento de empregos para municípios vizinhos, embora os pesquisadores não avaliem a eficiência do programa em termos de custo fiscal por emprego gerado.

Em perspectiva mais abrangente, \citeonline{Braz2023} avaliam o efeito agregado da renúncia fiscal da SUDENE sobre o emprego formal e a renda nos municípios beneficiados, empregando estimador DID escalonado para o período 2002--2021. Os resultados indicam incremento de 3,2\% nos vínculos formais e de 1,2\% na renda nominal média, com trajetória crescente ao longo do tempo e duradoura por todo o período de vigência do benefício. Os autores destacam, porém, que os ganhos concentram-se em municípios de maior porte e dinamismo econômico, o que contrasta com o objetivo redistributivo da PNDR e levanta questões sobre a focalização territorial do instrumento.

\citeonline{Costa2024} adotam o método de controle sintético generalizado para estimar o impacto dos benefícios tributários da SUDENE sobre a criação de empregos em empresas da região Nordeste, com dados de 2006 a 2019. As evidências apontam elevação de 19,6\% nos vínculos formais no conjunto dos setores e de 18,3\% no setor manufatureiro, com persistência dos resultados por até cinco anos após o início da concessão. A magnitude consideravelmente superior à estimada por \citeonline{Braz2023} pode refletir o nível de agregação --- empresa versus município --- e a concentração no setor industrial, caracterizado por maior intensidade de mão de obra entre as atividades incentivadas.

Os três estudos convergem em apontar impacto positivo sobre o emprego formal, porém com magnitudes que variam de 3,2\% a 30\% conforme o setor analisado, o método empregado e a escala de agregação. Os ganhos mais expressivos concentram-se em segmentos específicos --- turismo e manufatura ---, enquanto o efeito agregado sobre o conjunto da economia é mais modesto. A preocupação distributiva levantada por \citeonline{Braz2023} merece atenção: se os benefícios se concentram em municípios já dinâmicos, o instrumento pode estar ampliando, em vez de reduzir, as disparidades intrarregionais.

Em abordagem complementar às avaliações de impacto, \citeonline{Carneiro2023} analisam a eficiência produtiva das empresas beneficiadas pela política de incentivos da SUDENE no período 2011--2019, combinando análise envoltória de dados (DEA) e fronteira estocástica (SFA). O estudo identifica predominância de empresas intensivas em mão de obra entre as beneficiadas e presença relevante de unidades eficientes no Semiárido. A ineficiência observada está associada mais ao setor de atividade incentivado do que às firmas individualmente, o que sugere possibilidade de aprimorar a seletividade setorial da política. Esses achados indicam que a eficiência no uso dos benefícios tributários não é uniforme e pode ser aprimorada mediante critérios diferenciados de concessão.

No que se refere ao produto e a indicadores de desenvolvimento, as evidências são mais escassas e derivam majoritariamente de estudos multi-instrumento discutidos nas subseções anteriores. \citeonline{Carneiroetal2024a}, ao avaliarem conjuntamente FNE, FDNE e incentivos fiscais, não encontram significância estatística para a renúncia fiscal da SUDENE sobre o PIB \textit{per capita} municipal nem sobre o valor adicionado setorial. Em contrapartida, \citeonline{FerreiraIrffiCarneiro2024} obtêm efeito direto significativo dos benefícios tributários sobre o IDM-geral (0,020), IDM-saúde (0,024) e IDM-renda (0,012), além de transbordamento espacial positivo sobre a componente renda nos municípios vizinhos (0,002). A divergência entre os dois estudos pode estar associada à diferença na variável de resultado e na especificação econométrica adotada.

Complementarmente, \citeonline{gondim2025} identificam efeito direto significativo da renúncia fiscal sobre o crescimento do VAB agropecuário \textit{per capita}, resultado que sugere influência do instrumento sobre setores primários não captada pelos estudos com foco no produto agregado. \citeonline{Souza2025}, ao avaliarem simultaneamente todos os instrumentos da PNDR em painel dinâmico, constituem a única referência que inclui os benefícios tributários da SUDAM, embora os resultados para esse instrumento permaneçam inconclusivos dada a escassez de observações.

Dois estudos avaliam o Programa de Desenvolvimento do Estado de Pernambuco (Prodepe), incentivo fiscal estadual que opera em complementaridade com o FNE, o FDNE e a renúncia fiscal da SUDENE. \citeonline{Oliveira2020a} empregam modelo DID com efeitos heterogêneos por tempo de exposição e dados de empresas pernambucanas no período 2000--2017, estimando incremento de 8,6\% no emprego das firmas beneficiadas. Simultaneamente, o estudo identifica redução de 10,3\% no salário médio, padrão que contrasta com a ausência de efeito salarial observada nos estudos sobre Fundos Constitucionais e sobre os incentivos federais. Não se obteve significância para a massa salarial em empresas que recebem exclusivamente o Prodepe, embora a combinação com outros instrumentos altere os resultados.

\citeonline{Alves2024} confirmam esse padrão empregando estimador DID escalonado ao mesmo recorte geográfico e temporal. As estimativas apontam elevação de 22,3\% no emprego e de 15,1\% na massa salarial, acompanhadas de redução de 8,2\% no salário médio. Os ganhos são mais pronunciados em empresas do Semiárido, de pequeno porte e do setor industrial, porém apresentam caráter temporário, com declínio progressivo ao longo dos anos de vigência do benefício. A redução salarial, consistente nos dois estudos, sugere que os incentivos estaduais atraem ou expandem atividades de menor remuneração, o que coloca em questão a qualidade dos postos de trabalho gerados.

Em síntese, as evidências sobre a renúncia fiscal revelam convergência no tocante ao emprego --- com impactos positivos variando de 3,2\% a 30\% conforme setor e método --- e divergência quanto ao produto e ao salário. Nenhum estudo identifica efeito significativo exclusivo dos incentivos fiscais sobre o PIB \textit{per capita}, embora resultados positivos sejam observados sobre indicadores multidimensionais de desenvolvimento. No mercado de trabalho, a política federal gera empregos sem alterar o salário médio, ao passo que o programa estadual eleva o emprego mas reduz a remuneração, suscitando preocupação sobre a qualidade dos postos criados. A literatura apresenta lacunas relevantes: ausência de avaliações dedicadas aos incentivos da SUDAM, inexistência de análises de custo-efetividade e predominância de métodos DID, com escassa aplicação de abordagens como RDD, PSM ou equilíbrio geral. A ampliação e diversificação metodológica desta vertente constituem prioridade para a agenda de pesquisa.

 no main.tex

A literatura empírica sobre os instrumentos da Política Nacional de Desenvolvimento Regional apresenta expressiva assimetria em termos de volume de evidências: os Fundos Constitucionais concentram a quase totalidade dos estudos, enquanto os Fundos de Desenvolvimento e os Incentivos Fiscais constituem vertentes consideravelmente menos desenvolvidas. Os resultados obtidos variam substancialmente conforme a especificação econométrica adotada, o período amostral, a escala geográfica de análise e o controle de características locais não observáveis. Esta seção apresenta os principais achados da literatura, organizados por instrumento e abordagem metodológica, buscando identificar para cada um deles convergências, divergências e condições que modulam a eficácia da política.

\subsection{Avaliações de Impacto dos Fundos Constitucionais sobre o PIB}
\label{subsec:fc-pib}

Os Fundos Constitucionais são os instrumentos da PNDR mais amplamente avaliados pela literatura empírica: dentre os 34 estudos aprovados nesta revisão, 28 investigam os efeitos desses fundos sobre variáveis de resultado relacionadas ao produto interno bruto ou ao mercado de trabalho. A questão central é se esses fundos --- que operam via crédito subsidiado de ampla capilaridade territorial --- efetivamente contribuem para o crescimento econômico dos municípios beneficiados e, em caso afirmativo, sob quais condições e com que magnitude. O FNE é o fundo com maior número de avaliações (24 estudos), seguido pelo FCO (10 estudos) e pelo FNO (8 estudos). A predominância dos Fundos Constitucionais na literatura decorre de sua maior longevidade em relação aos demais instrumentos da PNDR -- com operação contínua desde 1989 --, da disponibilidade de dados municipais anuais de PIB a partir de 2002, do elevado volume de recursos envolvidos e da capilaridade dos programas em todos os municípios das macrorregiões-alvo. A quase totalidade dos estudos adota o PIB \textit{per capita} municipal como variável de resultado, tratando-o como aproximação da renda média local, dada a escassez de alternativas em frequência anual e cobertura municipal. Os resultados obtidos pela literatura variam consideravelmente conforme a especificação econométrica adotada, o período amostral, a escala geográfica de análise e o controle de características locais não observáveis. As subseções seguintes organizam os achados por abordagem metodológica, permitindo identificar padrões gerais e divergências nas estimativas de impacto. O Quadro~\ref{tab:estudos_fc_pib} apresenta a síntese dos 21 estudos que avaliam os efeitos dos Fundos Constitucionais sobre o PIB \textit{per capita}, organizados por método de análise.

% Arquivo: latex/tabelas/tabela_macros.tex
% Macros para formatação de tabelas de resultados da revisão sistemática

% Macro principal para formatação de células com listas (label "--")
% Uso: \fc[largura]{conteúdo em itemize}
% Padrão: largura = 4.2cm (coluna Estimativa)
\newcommand{\fc}[2][4.2cm]{%
	\parbox[t]{#1}{\raggedright%
		\begin{itemize}[leftmargin=*,nosep,topsep=0pt,before=\vspace{-0.6\baselineskip},label={--},labelsep=0.4em]
			#2
		\end{itemize}%
		\vspace{0.3\baselineskip}%
	}%
}

% Macro para célula com label "--" (variáveis, amostragem)
% Uso: \fcvar[largura]{conteúdo em itemize}
% Padrão: largura = 2.2cm (colunas Amostragem e Variáveis)
\newcommand{\fcvar}[2][2.2cm]{%
	\parbox[t]{#1}{\raggedright%
		\begin{itemize}[leftmargin=*,nosep,topsep=0pt,before=\vspace{-0.6\baselineskip},label={--},labelsep=0.4em]
			#2
		\end{itemize}%
		\vspace{0.3\baselineskip}%
	}%
}

% Macro para mesclar células na coluna de método
% Parâmetros: #1 = texto, #2 = número de linhas a mesclar
\newcommand{\metodo}[2]{%
	\multirow{#2}{*}{\begin{minipage}[t]{1.4cm}\raggedright #1 \end{minipage}}%
}

% Arquivo: latex/tabelas/survey_artigos_fc_pib.tex
% Gerado automaticamente por scripts/generate_survey_tables.py
% Fonte: data/2-papers/survey_results.json

% IMPORTANTE: Este arquivo requer:
% - Pacotes: longtable, multirow, afterpage, enumitem, array
% - Macros: \fc e \metodo (definidas em tabelas/tabela_macros.tex)
% - Comando \citeonline (do pacote abntex2cite)
% - Contador: quadro (definido na classe abntex2)

\afterpage{\clearpage%
\tiny%
\setlength{\extrarowheight}{1pt}%
\setlength{\tabcolsep}{2pt}%
\linespread{0.9}\selectfont%
\renewcommand{\arraystretch}{1.2}%
\renewcommand{\tablename}{\quadroname}%
\renewcommand{\thetable}{\arabic{quadro}}%
\stepcounter{quadro}%
\begin{longtable}{|>{\raggedright\arraybackslash}p{2.2cm}
		|>{\raggedright\arraybackslash}p{2.2cm}
		|>{\raggedright\arraybackslash}p{2.3cm}
		|>{\raggedright\arraybackslash}p{3.0cm}
		|>{\raggedright\arraybackslash}p{4.2cm}|}

	\caption{Estudos de Impacto dos Fundos Constitucionais sobre o PIB \textit{per capita}} \label{tab:estudos_fc_pib} \\

\hline
	\begin{tabular}[c]{@{}>{\centering\arraybackslash}m{2.2cm}@{}}Método\end{tabular}
	& \begin{tabular}[c]{@{}>{\centering\arraybackslash}m{2.2cm}@{}}Artigo\end{tabular}
	& \begin{tabular}[c]{@{}>{\centering\arraybackslash}m{2.3cm}@{}}Amostragem\end{tabular}
	& \begin{tabular}[c]{@{}>{\centering\arraybackslash}m{3.0cm}@{}}Variáveis\end{tabular}
	& \begin{tabular}[c]{@{}>{\centering\arraybackslash}m{4.2cm}@{}}Estimativa\end{tabular} \\
	\hline
	\endfirsthead

	\multicolumn{5}{c}%
	{{\quadroname\ \thetable{} -- Continuação}} \\
	\hline
	\begin{tabular}[c]{@{}>{\centering\arraybackslash}m{2.2cm}@{}}Método\end{tabular}
	& \begin{tabular}[c]{@{}>{\centering\arraybackslash}m{2.2cm}@{}}Artigo\end{tabular}
	& \begin{tabular}[c]{@{}>{\centering\arraybackslash}m{2.3cm}@{}}Amostragem\end{tabular}
	& \begin{tabular}[c]{@{}>{\centering\arraybackslash}m{3.0cm}@{}}Variáveis\end{tabular}
	& \begin{tabular}[c]{@{}>{\centering\arraybackslash}m{4.2cm}@{}}Estimativa\end{tabular} \\
	\hline
	\endhead

	\hline
	\multicolumn{5}{r}{{Continua}} \\
	\endfoot

	\hline
	\multicolumn{5}{l}{
		\tiny
		\parbox{14.8cm}{%
			\vspace{0.1cm}
			Fonte: Elaborado pelo autor. Siglas: AC - Apresentação em congresso; TD - Trabalho para discussão; NS - Não significativo.
	}}
	\endlastfoot

Painel de Efeito Fixo
	& \citeonline{Resende2014a}
	& \fcvar[2.3cm]{\item NE \item 2004-2010}
	& \fcvar[3.0cm]{\item FC/PIB \item Cresc. PIB pc}
	& \fc[4.2cm]{
		\item FNE (munic) \hfill 0,021\textsuperscript{*}
		\item FNE (micro) \hfill 0,032\textsuperscript{*}
		\item FNE (meso) \hfill --0,288\textsuperscript{*}
		\item FNE-setorial \hfill NS\phantom{\textsuperscript{*}}
	}
	\\
	\cline{2-5}


	& \citeonline{Resende2014b}
	& \fcvar[2.3cm]{\item N \item 2004-2010}
	& \fcvar[3.0cm]{\item FC/PIB \item Cresc. PIB pc}
	& \fc[4.2cm]{
		\item FNO (munic) \hfill --0,399\textsuperscript{*}
		\item FNO (micro e meso) \hfill NS\phantom{\textsuperscript{*}}
		\item FNO setorial \hfill NS\phantom{\textsuperscript{*}}
	}
	\\
	\cline{2-5}


	& \citeonline{ResendeCravo2014}
	& \fcvar[2.3cm]{\item CO \item 2004-2010}
	& \fcvar[3.0cm]{\item FC/PIB \item Cresc. PIB pc}
	& \fc[4.2cm]{
		\item FCO (munic) \hfill 0,076\textsuperscript{*}
		\item FCO (micro e meso) \hfill NS\phantom{\textsuperscript{*}}
		\item FCO-empresarial (munic) \hfill 0,119\textsuperscript{*}
		\item FCO-rural (munic) \hfill --0,072\textsuperscript{*}
	}
	\\

	\hline
Modelo de \textit{Spillover} Espacial
	& \citeonline{Oliveira2005}
	& \fcvar[2.3cm]{\item N, CO \item 1991-2000}
	& \fcvar[3.0cm]{\item FC/PIB \item Cresc. PIB pc}
	& \fc[4.2cm]{
		\item FNO e FCO \hfill NS\phantom{\textsuperscript{*}}
	}
	\\
	\cline{2-5}


	& \citeonline{Cravo2015}
	& \fcvar[2.3cm]{\item N, NE, CO \item 2004-2010}
	& \fcvar[3.0cm]{\item FC/PIB \item Cresc. PIB pc}
	& \fc[4.2cm]{
		\item FNE (munic) \hfill NS\phantom{\textsuperscript{*}}
		\item FNO (munic) \hfill --0,313\textsuperscript{*}
		\item FCO (munic) \hfill 0,087\textsuperscript{*}
		\item FCO-empresarial (munic) \hfill +0,129\textsuperscript{*}
		\item \textit{Spillovers} espaciais \hfill NS\phantom{\textsuperscript{*}}
	}
	\\
	\cline{2-5}


	& \citeonline{Oliveira2016}
	& \fcvar[2.3cm]{\item GO \item 2004-2011}
	& \fcvar[3.0cm]{\item Log. FCO \item Log. PIB pc}
	& \fc[4.2cm]{
		\item FCO (painel espacial) \hfill 0,098\textsuperscript{*}
		\item FCO (painel clássico) \hfill 0,102\textsuperscript{*}
	}
	\\
	\cline{2-5}


	& \citeonline{Resende2017}
	& \fcvar[2.3cm]{\item NE \item 1999-2011}
	& \fcvar[3.0cm]{\item FC/PIB \item Cresc. PIB pc}
	& \fc[4.2cm]{
		\item FNE (munic. tipologia dinâmica)
		\item[\textbullet] Efeito direto \hfill 0,077\textsuperscript{*}
		\item[\textbullet] Efeito indireto \hfill 0,329\textsuperscript{*}
	}
	\\
	\cline{2-5}


	& \citeonline{Resende2018}
	& \fcvar[2.3cm]{\item N, NE, CO \item 1999-2011}
	& \fcvar[3.0cm]{\item FC/PIB \item Cresc. PIB pc}
	& \fc[4.2cm]{
		\item FC (munic. tip. dinâmica - acumulado)
		\item[\textbullet] Efeito direto \hfill 0,033\textsuperscript{*}
		\item[\textbullet] Efeito indireto \hfill 0,118\textsuperscript{*}
		\item FC (micro. tip. baixa renda - acumulado)
		\item[\textbullet] Efeito direto \hfill 0,146\textsuperscript{*}
		\item[\textbullet] Efeito indireto \hfill 0,466\textsuperscript{*}
	}
	\\
	\cline{2-5}


	& \citeonline{Silva2023}
	& \fcvar[2.3cm]{\item N, NE, CO \item 2016-2019}
	& \fcvar[3.0cm]{\item Log. FC \item Log. PIB pc}
	& \fc[4.2cm]{
		\item FNE, FNO e FCO \hfill NS\phantom{\textsuperscript{*}}
	}
	\\
	\cline{2-5}


	& \citeonline{Oliveira2026}
	& \fcvar[2.3cm]{\item NE \item 2010-2021}
	& \fcvar[3.0cm]{\item FC/PIB \item Cresc. PIB pc}
	& \fc[4.2cm]{
		\item FNE (efeito direto) \hfill 0,00001\textsuperscript{*}
		\item FNE (efeito indireto) \hfill NS\phantom{\textsuperscript{*}}
		\item FNE-serviços (efeito direto) \hfill 0,00001\textsuperscript{*}
		\item FNE-indústria (efeito direto) \hfill 0,00005\textsuperscript{*}
	}
	\\

	\hline
Modelo \textit{Threshold}
	& \citeonline{Linhares2014}
	& \fcvar[2.3cm]{\item NE \item 2000-2008}
	& \fcvar[3.0cm]{\item FC/PIB \item Cresc. PIB pc}
	& \fc[4.2cm]{
		\item FNE (PIB pc < 2.143) \hfill NS\phantom{\textsuperscript{*}}
		\item FNE (2.143 < PIB pc < 3.866) \hfill 0,078\textsuperscript{*}
		\item FNE (3.866 < PIB pc < 7.406) \hfill 0,109\textsuperscript{*}
		\item FNE (PIB pc > 7.406) \hfill NS\phantom{\textsuperscript{*}}
	}
	\\

	\hline
Primeira Diferença
	& \citeonline{Resende2014c}
	& \fcvar[2.3cm]{\item NE \item 1999-2006}
	& \fcvar[3.0cm]{\item FC/PIB \item Cresc. PIB pc}
	& \fc[4.2cm]{
		\item FNE-indústria (munic) \hfill NS\phantom{\textsuperscript{*}}
	}
	\\

	\hline
Regressão Quantílica e IV
	& \citeonline{Irffi2016}
	& \fcvar[2.3cm]{\item NE \item 2000-2010}
	& \fcvar[3.0cm]{\item FC/PIB \item Cresc. PIB pc}
	& \fc[4.2cm]{
		\item FNE \hfill (+)\textsuperscript{*}
		\item FNE-pecuária e agricultura \hfill (+)\textsuperscript{*}
	}
	\\

	\hline
Modelo de Eficiência
	& \citeonline{Carneiro2018}
	& \fcvar[2.3cm]{\item NE \item 2010-2015}
	& \fcvar[3.0cm]{\item FC/PIB \item Cresc. PIB pc}
	& \fc[4.2cm]{
		\item FNE: Maior eficiência em municípios integrados ao comércio internacional (Matopiba).
	}
	\\

	\hline
Painel Dinâmico
	& \citeonline{Cambota2019}
	& \fcvar[2.3cm]{\item NE \item 2003-2014}
	& \fcvar[3.0cm]{\item FC/PIB \item Cresc. PIB pc}
	& \fc[4.2cm]{
		\item FNE \hfill 2,96\textsuperscript{*}
		\item Efeito contemporâneo e endógeno significativo.
	}
	\\

	\hline
Equilíbrio Geral
	& \citeonline{Nascimento2017}
	& \fcvar[2.3cm]{\item BR \item 2000-2011}
	& \fcvar[3.0cm]{\item FC \item PIB}
	& \fc[4.2cm]{
		\item Remoção do FNE (alocação em despesa corrente) \hfill --0,16\%\phantom{\textsuperscript{*}}
		\item Maior concentração regional da produção
	}
	\\
	\cline{2-5}


	& \citeonline{Ribeiro2020}
	& \fcvar[2.3cm]{\item BR \item 2014-2025}
	& \fcvar[3.0cm]{\item FC/PIB \item Cresc. PIB pc}
	& \fc[4.2cm]{
		\item FNE (PIB NE) \hfill 3,51\%\textsuperscript{*}
	}
	\\
	\cline{2-5}


	& \citeonline{Goncalves2022}
	& \fcvar[2.3cm]{\item NE \item 2006-2015}
	& \fcvar[3.0cm]{\item Fundo \item Log. PIB pc}
	& \fc[4.2cm]{
		\item FNE (remoção): redução da atividade econômica no NE
	}
	\\

	\hline
RDD Geográfico
	& \citeonline{Oliveira2021}
	& \fcvar[2.3cm]{\item NE \item 1970-2010}
	& \fcvar[3.0cm]{\item Log. FNE \item Log. PIB pc}
	& \fc[4.2cm]{
		\item FNE \hfill NS\phantom{\textsuperscript{*}}
	}
	\\

	\hline
\textit{Propensity Score}
	& \citeonline{Monte2025}
	& \fcvar[2.3cm]{\item NE \item 2010-2020}
	& \fcvar[3.0cm]{\item FNE \item Log. PIB pc}
	& \fc[4.2cm]{
		\item FNE (PJ) \hfill 0,15pp\textsuperscript{*}
		\item FNE (mulheres) \hfill 0,10pp\textsuperscript{*}
		\item Efeito heterogêneo por magnitude
	}
	\\

	\hline
DID 2 Estágios, DID Escalonado
	& \citeonline{CarneiroVeloso2024}
	& \fcvar[2.3cm]{\item NE \item 2002-2019}
	& \fcvar[3.0cm]{\item FNE (\textit{dummy}) \item Log. PIB pc}
	& \fc[4.2cm]{
		\item FNE \hfill 0,08\textsuperscript{*} a 0,11\textsuperscript{*}
		\item A partir do 5º ano após o tratamento
		\item Principalmente serviços e administração pública\hfill \phantom{\textsuperscript{*}}
	}
	\\

\end{longtable}%
\normalsize%
\setlength{\extrarowheight}{0pt}%
\setlength{\tabcolsep}{6pt}%
\linespread{1.5}\selectfont%
\renewcommand{\tablename}{Tabela}%
\renewcommand{\thetable}{\arabic{table}}%
}


\subsubsection{Modelos de Efeitos Fixos}

Modelos de painel com efeitos fixos controlam características não observáveis constantes no tempo (como geografia e instituições locais) e choques temporais comuns a todas as unidades (como crises macroeconômicas), permitindo isolar de forma mais precisa o impacto da política regional. Essa abordagem foi a primeira empregada para avaliar os Fundos Constitucionais sobre o PIB municipal, aproveitando a disponibilidade de dados em painel desde 2002.

\citeonline{Resende2014a} analisa a influência do FNE sobre o crescimento do PIB \textit{per capita} dos municípios nordestinos no período 2004-2010, utilizando médias em três subperíodos (2004-2006, 2006-2008 e 2008-2010) para minimizar viés de ciclos de negócios. Os resultados apontam impacto positivo do FNE-total em nível municipal (0,021) e microrregional (0,032), mas não significativo em nível mesorregional, o que Resende (2014a) atribui à maior heterogeneidade em regiões geográficas maiores. As estimativas sugerem que o resultado positivo estaria concentrado nos empréstimos ao setor agropecuário, porém não se mostraram robustas à inclusão de controles temporais. Para o FNO, aplicando a mesma especificação econométrica à região Norte, \citeonline{Resende2014a} obtém relação inversa entre o fundo e a expansão do produto municipal, com ausência de robustez nas análises setoriais quando se controla por períodos.

\citeonline{Oliveira2017a} aplicam especificações de efeitos fixos juntamente com outros métodos (PSM e GPS) para avaliar o FCO no estado de Goiás, obtendo evidências de impacto positivo sobre a dinâmica econômica municipal, com magnitude variando conforme o método empregado.

Mais recentemente, \citeonline{Filho2024} adotam painel de efeitos fixos para avaliar conjuntamente os três fundos constitucionais, não obtendo evidências consistentes de associação positiva com o crescimento do PIB \textit{per capita} para o período 2016-2019, resultado que contrasta com os estudos anteriores.

Esses estudos ressaltam limitações metodológicas importantes. A principal dificuldade consiste na identificação causal, dado que a alocação dos fundos é endógena à demanda por crédito subsidiado, o que tende a privilegiar municípios com maior capacidade de absorção e infraestrutura socioeconômica mais desenvolvida. Adicionalmente, identificam-se dificuldades em mensurar os custos fiscais efetivos dos programas, a possibilidade de existência de peso morto (investimentos que seriam realizados mesmo sem subsídio) e potenciais deslocamentos de fatores produtivos entre municípios vizinhos. A ausência de robustez à inclusão de controles temporais sugere que parte do impacto estimado pode capturar tendências macroeconômicas comuns aos municípios tratados e não tratados, em vez do efeito causal dos fundos. Em síntese, os modelos de efeitos fixos indicam correlação positiva dos fundos constitucionais com a expansão econômica local, com impactos mais pronunciados em nível municipal e maior magnitude para o FCO, mas a fragilidade da identificação causal e a ausência de robustez em diversas especificações limitam a interpretação causal dos resultados.

\subsubsection{Modelos de Painel Espacial}

Uma vertente da literatura admite a possibilidade de interdependência econômica entre unidades espaciais geograficamente próximas, o que pode ocasionar efeitos indiretos de transbordamento da política para outras regiões, expandindo a área potencialmente beneficiada. Por outro lado, o efeito poderia consistir em vazamentos de recursos das regiões-alvo da política em direção a regiões de maior dinamismo econômico. Portanto, entender a natureza da dinâmica espacial é importante para um planejamento mais eficiente da política. \citeonline{Resende2014c} utilizam modelo de painel SDM (\textit{Spatial Durbin Model}) e estimam efeito significativo do FCO sobre o PIB \textit{per capita}, com evidência de que uma elevação de 1 ponto percentual (p.p.) na relação FCO/PIB está associada a um incremento de cerca de 0,09 p.p. na taxa média anual de crescimento do PIB \textit{per capita} municipal nos dois anos seguintes. Na análise setorial, o resultado seria atribuído à modalidade de crédito FCO-Empresarial, em que o efeito seria na ordem de 0,13 p.p. de incremento. Porém, o resultado é não significativo para o FNE e negativo para o FNO. Com respeito aos efeitos indiretos (transbordamento espacial), obteve-se resultado relevante, mas não positivo, para o FCO, o que pode estar ligado à migração de fatores de produção.

\citeonline{Oliveira2015} usam modelo espacial para avaliar efeitos do FCO nos municípios de Goiás e concluem que efeitos diretos são menores quando se admite a ocorrência de efeitos espaciais, sugerindo que modelos não espaciais tendem a superestimar o impacto real dos fundos. Entretanto, as evidências obtidas por \citeonline{Resende2014c} e \citeonline{Oliveira2015} apontam em direção contrária, com os efeitos diretos do FCO maiores em modelos espaciais do que em modelos não espaciais. Essa divergência pode estar associada a diferenças na cobertura geográfica (Goiás versus toda a região Centro-Oeste), no tipo de matriz de pesos espaciais empregada (\textit{queen} de contiguidade versus distância inversa) e no período de análise considerado. Os resultados de \citeonline{Oliveira2015} indicam que cada 1\% de aumento no valor dos financiamentos está associado a aproximadamente 0,1\% de crescimento no PIB \textit{per capita} do município, usando especificação SAR com matriz espacial \textit{queen} do tipo contiguidade e dados de 2004 a 2011.

\citeonline{ResendeSilvaFilho2017} distinguem-se dos trabalhos anteriores por incorporar na análise a heterogeneidade definida pela própria PNDR por meio das tipologias regionais. Em modelo não espacial de efeitos fixos, o efeito estimado do FNE é significativo para municípios pertencentes às tipologias dinâmica e baixa renda. Quando se admite dependência espacial, há evidências de transbordamento espacial dos investimentos do FNE quando o município diretamente beneficiado tem tipologia dinâmica. As estimativas apontam que para cada 1 p.p. de aumento na relação FNE/PIB em municípios dinâmicos há um incremento de 0,07 p.p. na taxa média anual de crescimento do PIB \textit{per capita} do próprio município nos três anos seguintes e de 0,33 p.p. na taxa média dos municípios vizinhos. Para os autores, esses resultados ressaltam a relevância de se considerar a dimensão espacial nas análises e também o caráter heterogêneo dos efeitos do FNE nos territórios. \citeonline{ResendeSilvaFilho2018} adotam a mesma abordagem econométrica, mas incluindo também as regiões Norte e Centro-Oeste, a fim de avaliar conjuntamente os fundos FNE, FNO e FCO. Os resultados são positivos para efeitos diretos e indiretos em municípios de tipologia dinâmica quando se usam valores acumulados na variável dos fundos (dois primeiros anos de cada subperíodo). Os resultados apontam que um aumento de 1 p.p. na relação FC/PIB em municípios dinâmicos está associado a 0,03 p.p. de incremento no crescimento médio do PIB \textit{per capita} do próprio município pelos três anos seguintes e 0,12 p.p. de incremento nos municípios vizinhos. \citeonline{Filho2024} não obtêm evidências de associação espacial positiva entre os financiamentos do FNE, FNO e FCO nos anos de 2016 a 2018 e o crescimento médio do PIB \textit{per capita} dos municípios de 2016 a 2019, porém, o resultado é positivo sobre o PIB \textit{per capita} de 2019. Os autores concluem que os efeitos dos fundos são em sua maior parte esporádicos e pontuais.

Em resumo, os estudos convergem para a relevância de se incorporar a dimensão espacial nas análises de efeito dos fundos constitucionais, justificada pela ocorrência de correlação espacial significativa nas variáveis envolvidas. Porém, há menor consenso quando se trata da magnitude dos efeitos. Em geral, o resultado estaria condicionado a outras características econômicas locais. O efeito direto mostrou-se significativo em municípios de tipologia dinâmica e microrregiões de baixa renda, com incrementos de 0,03 a 0,14 p.p. no crescimento anual do PIB \textit{per capita} municipal por até dois anos, para cada 1 p.p. de aumento na relação Fundo/PIB. Para o efeito indireto, as estimativas variam entre 0,1 e 0,4 p.p. de incremento no crescimento para cada 1 p.p. de aumento em Fundo/PIB, também concentrado em municípios de tipologia dinâmica e microrregiões de baixa renda. Em análise mais desagregada, evidências apontam para associação entre o FCO e taxa de crescimento, atribuída à linha de crédito empresarial. No FNE, os resultados estariam ligados ao setor industrial e de serviços. Já para o FNO não há nenhuma evidência positiva e significativa. No aspecto metodológico, nota-se a necessidade de estudos mais robustos, que combinem intervalo de tempo mais longo, especificações espaciais alternativas, correção de endogeneidade e maior número de controles (inclusive variáveis de confusão da própria PNDR).

\subsubsection{Modelos Não Lineares e de Eficiência}

\citeonline{Linharesetal2014} usam modelo \textit{threshold} e obtêm evidências de que o efeito do FNE é pouco significativo para municípios com PIB \textit{per capita} inicial abaixo de R\$ 2.143 ou acima de R\$ 7.406. Entre esses limites, os autores identificam dois grupos distintos. O primeiro é formado por municípios de PIB \textit{per capita} inicial entre R\$ 2.143 e R\$ 3.866. Nesse grupo, estima-se que 1 p.p. de aumento médio em FNE/PIB no período 2002-2006 esteja associado a um incremento de 0,08 p.p. no crescimento médio do PIB \textit{per capita} entre 2002 e 2008. No segundo grupo, de municípios com PIB \textit{per capita} inicial entre R\$ 3.866 e R\$ 7.406, o efeito estimado foi maior, na ordem de 0,11 p.p. Os resultados sugerem que o FNE tem maior efetividade em municípios com níveis intermediários de renda, reafirmando os achados de \citeonline{ResendeSilvaFilho2017} e \citeonline{ResendeSilvaFilho2018}. Em municípios de baixa renda, porém, os resultados não foram significativos, o que indica a necessidade de aprimoramento dos mecanismos da política para melhor adequação às características e necessidades dessas regiões.

\citeonline{Oliveira2015} usam modelagem quantílica para estimar a significância e heterogeneidade de efeitos do FNE ao longo da distribuição de crescimento dos municípios, empregando estimação em dois estágios e a localização de agências bancárias do BNB como instrumento para a alocação de recursos do fundo. Os efeitos foram positivos e significativos para o FNE total e para as linhas de crédito associadas à agricultura e pecuária. No entanto, os largos intervalos de confiança não permitiram constatar presença de heterogeneidade ao longo da curva de crescimento. \citeonline{Oliveira2018} analisa a relação entre a participação do FNE e a eficiência média dos municípios nordestinos. Usando modelo de análise envoltória de dados e de fronteira estocástica, o autor identifica elevada eficiência nos municípios pertencentes à região do MATOPIBA, associada à agricultura extensiva voltada à exportação. Adicionalmente, municípios do Semiárido indicaram menor eficiência no crescimento do produto médio, enquanto municípios menos densos e mais distantes das respectivas capitais e com empresas menores tendem a fazer melhor uso dos financiamentos.

\subsubsection{Modelos Dinâmicos}

\citeonline{Cambota2019} adotam modelo de painel dinâmico e instrumentos sintéticos para corrigir o viés de endogeneidade na variável PIB defasada e na variável de interesse FNE/PIB, usando dados de 2003 a 2014. Os autores confirmam a importância desse tipo de controle para a obtenção de estimativas mais realistas do impacto da política. As estimativas apontam associação positiva entre o FNE e o crescimento econômico dos municípios, na ordem de 2,96 p.p. para cada 1 p.p. de aumento em FNE/PIB no ano anterior. Cabe notar que essa magnitude é consideravelmente superior à dos demais estudos revisados, cujos coeficientes situam-se tipicamente entre 0,03 e 0,13 p.p. A discrepância pode ser atribuída à diferente especificação do modelo (painel dinâmico com instrumentos sintéticos), ao período mais longo de análise (2003-2014) ou a possível sobreestimação decorrente da proliferação de instrumentos no procedimento GMM, aspecto reconhecido na literatura como potencial fonte de viés.

\subsubsection{Modelos de Equilíbrio Geral}

Em \citeonline{Nascimento2017}, um cenário hipotético de retirada do FNE e realocação dos recursos em gastos correntes do governo resultaria a longo prazo em uma redução líquida de 0,16\% no PIB nacional e em maior concentração regional da produção. O estudo também sugere que os fundos têm papel mais relevante no crescimento de longo prazo, via elevação do estoque de capital e da capacidade produtiva local, frente à alternativa apontada. No curto prazo, ambos os cenários seriam equivalentes em termos de ativação da demanda. \citeonline{Ribeiroetal2020} indicam que os investimentos do FNE realizados em 2014 e 2015 teriam elevado o PIB da região Nordeste em 3,51\% até o ano 2025. \citeonline{Oliveira2021} constroem dois cenários alternativos envolvendo o término da política de subsídios, e ambos indicam como consequência a redução na atividade econômica na região Nordeste, puxada principalmente pela redução do produto da atividade agropecuária.

\subsubsection{Modelos Quase Experimentais}

\citeonline{Resende2014c} usa o método de primeira diferença para avaliar o efeito de empréstimos industriais do FNE sobre o crescimento do PIB \textit{per capita} municipal. Usando taxas de crescimento do PIB \textit{per capita} nos períodos 2000-2003 e 2000-2006, o autor não encontra efeito significativo das operações envolvendo o FNE-Indústria durante o ano 2000, seja em nível municipal ou microrregional. Argumenta-se que o resultado pode estar associado à elevada participação do setor agropecuário nas operações do FNE, setor em que há predomínio da informalidade, o que dificulta sua captura pelas medições do PIB.

\citeonline{Oliveira2020} adotam desenho quase-experimental que combina regressão com descontinuidade geográfica (RDD) e modelo de diferenças em diferenças. Os autores constroem um painel de dados decenal (1980, 1990, 2000 e 2010) composto por municípios dos estados de Minas Gerais e Espírito Santo que ingressaram nos critérios de elegibilidade ao FNE em períodos recentes (dados pré e pós-ingresso) e por municípios vizinhos imediatos que nunca tenham sido tratados pelo FNE. Os resultados obtidos levantam dúvidas sobre a efetividade do FNE, sugerindo que o Fundo não teve influência sobre as variáveis de produto e renda nos municípios avaliados no ano 2010. O estudo ressalta a importância de melhor vinculação do crédito ao aumento da produtividade, mediante condições diferenciadas e combinação com outros instrumentos.

Por outro lado, \citeonline{Carneiroetal2024a} obtêm evidências de efeito positivo e significativo do FNE sobre o PIB \textit{per capita} municipal. Os autores usam dados de 2002 a 2019 junto a estimadores DID dois estágios e DID escalonado, combinando três instrumentos de política regional: FNE, FDNE e incentivo fiscal da SUDENE. O estudo aponta efeito médio de tratamento de 8,8\% a 11\% no PIB \textit{per capita} do município beneficiado em relação ao grupo de controle, sendo atribuído principalmente ao setor serviços e administração pública. O estudo de eventos aponta que o efeito de tratamento é crescente e inicia no quinto ano após o primeiro tratamento. Quanto à sinergia com os demais instrumentos, o resultado é não significativo. \citeonline{MonteIrffiBastosCarneiro2025} exploram heterogeneidades de gênero, tipo de tomador e magnitude dos empréstimos (efeito dosagem) do FNE, aplicando o método \textit{Generalized Propensity Score} (GPS). O estudo encontra correlação positiva entre a participação do FNE-Mulheres no valor total do programa e a variação do PIB \textit{per capita} municipal, o que sinaliza uma eficácia maior do programa nesse grupo populacional. Efeito semelhante é obtido em relação à variável participação de tomadores pessoa jurídica no total de financiamentos. Com respeito ao efeito dosagem, os resultados apontam efeitos positivos no PIB \textit{per capita} somente na última e mais elevada dose de tratamento.

Em conjunto, as evidências sobre o efeito dos Fundos Constitucionais no PIB \textit{per capita} revelam um quadro heterogêneo e condicionado à especificação metodológica adotada. Os modelos de efeitos fixos e espaciais indicam associação positiva para o FCO e resultados ambíguos para FNE e FNO, enquanto os modelos quase experimentais apontam tanto para ausência de efeito quanto para impactos positivos significativos. Os modelos não lineares e de eficiência sugerem que a eficácia dos fundos é condicionada ao nível inicial de renda do município, com efeitos mais pronunciados em faixas intermediárias de desenvolvimento. As simulações de equilíbrio geral reafirmam a importância dos fundos para a atividade econômica regional em cenários contrafactuais de retirada da política. A constatação de que os efeitos são mais pronunciados em municípios com níveis intermediários de renda, e nulos ou inconclusivos em municípios de baixa renda, pode ser interpretada à luz da literatura sobre capacidade de absorção e armadilhas de pobreza: municípios de baixa renda podem carecer do capital humano e da infraestrutura básica necessários para converter o crédito subsidiado em crescimento econômico sustentado; municípios de renda intermediária possuiriam o mínimo de capacidade necessária para ativar o efeito multiplicador dos investimentos; e municípios de alta renda estariam próximos do efeito marginal decrescente do capital.

\subsection{Avaliações de Impacto dos Fundos Constitucionais sobre o Mercado de Trabalho}

Dezesseis estudos aprovados nesta revisão avaliam o impacto dos Fundos Constitucionais sobre variáveis do mercado de trabalho, com foco predominante no emprego formal, massa salarial e salário médio. O Quadro~\ref{tab:estudos_fc_emprego} apresenta a síntese desses estudos, organizados por método de análise.

% Arquivo: latex/tabelas/survey_artigos_fc_emprego.tex
% Gerado automaticamente por scripts/generate_survey_tables.py
% Fonte: data/2-papers/survey_results.json

% IMPORTANTE: Este arquivo requer:
% - Pacotes: longtable, multirow, afterpage, enumitem, array
% - Macros: \fc e \metodo (definidas em tabelas/tabela_macros.tex)
% - Comando \citeonline (do pacote abntex2cite)
% - Contador: quadro (definido na classe abntex2)

\afterpage{\clearpage%
\tiny%
\setlength{\extrarowheight}{1pt}%
\setlength{\tabcolsep}{2pt}%
\linespread{0.9}\selectfont%
\renewcommand{\arraystretch}{1.2}%
\renewcommand{\tablename}{\quadroname}%
\renewcommand{\thetable}{\arabic{quadro}}%
\stepcounter{quadro}%
\begin{longtable}{|>{\raggedright\arraybackslash}p{2.2cm}
		|>{\raggedright\arraybackslash}p{2.2cm}
		|>{\raggedright\arraybackslash}p{2.3cm}
		|>{\raggedright\arraybackslash}p{3.0cm}
		|>{\raggedright\arraybackslash}p{4.2cm}|}

	\caption{Estudos de Impacto dos Fundos Constitucionais sobre o Mercado de Trabalho} \label{tab:estudos_fc_emprego} \\

\hline
	\begin{tabular}[c]{@{}>{\centering\arraybackslash}m{2.2cm}@{}}Método\end{tabular}
	& \begin{tabular}[c]{@{}>{\centering\arraybackslash}m{2.2cm}@{}}Artigo\end{tabular}
	& \begin{tabular}[c]{@{}>{\centering\arraybackslash}m{2.3cm}@{}}Amostragem\end{tabular}
	& \begin{tabular}[c]{@{}>{\centering\arraybackslash}m{3.0cm}@{}}Variáveis\end{tabular}
	& \begin{tabular}[c]{@{}>{\centering\arraybackslash}m{4.2cm}@{}}Estimativa\end{tabular} \\
	\hline
	\endfirsthead

	\multicolumn{5}{c}%
	{{\quadroname\ \thetable{} -- Continuação}} \\
	\hline
	\begin{tabular}[c]{@{}>{\centering\arraybackslash}m{2.2cm}@{}}Método\end{tabular}
	& \begin{tabular}[c]{@{}>{\centering\arraybackslash}m{2.2cm}@{}}Artigo\end{tabular}
	& \begin{tabular}[c]{@{}>{\centering\arraybackslash}m{2.3cm}@{}}Amostragem\end{tabular}
	& \begin{tabular}[c]{@{}>{\centering\arraybackslash}m{3.0cm}@{}}Variáveis\end{tabular}
	& \begin{tabular}[c]{@{}>{\centering\arraybackslash}m{4.2cm}@{}}Estimativa\end{tabular} \\
	\hline
	\endhead

	\hline
	\multicolumn{5}{r}{{Continua}} \\
	\endfoot

	\hline
	\multicolumn{5}{l}{
		\tiny
		\parbox{14.8cm}{%
			\vspace{0.1cm}
			Fonte: Elaborado pelo autor. Siglas: NS - Não significativo; p.p. - Pontos percentuais.
	}}
	\endlastfoot

\textit{Propensity Score Matching}
	& \citeonline{SilvaResende2009}
	& \fcvar[2.3cm]{\item Empresas NE, N, CO \item 2000-2003}
	& \fcvar[3.0cm]{\item FNE, FNO, FCO (\textit{dummy}) \item Cresc. emprego \item Cresc. salário}
	& \fc[4.2cm]{
		\item FNE-emprego \hfill +61,2 p.p.\textsuperscript{*}
		\item FNE-salário \hfill NS\phantom{\textsuperscript{*}}
		\item FNO, FCO \hfill NS\phantom{\textsuperscript{*}}
	}
	\\
	\cline{2-5}


	& \citeonline{Soares2017}
	& \fcvar[2.3cm]{\item Empresas NE \item 2000-2006}
	& \fcvar[3.0cm]{\item FNE (\textit{dummy}) \item Cresc. emprego \item Cresc. massa salarial \item Cresc. salário}
	& \fc[4.2cm]{
		\item FNE-empr. (1 ano) \hfill +7,96 p.p.\textsuperscript{*}
		\item FNE-empr. (3 anos) \hfill +33,59 p.p.\textsuperscript{*}
		\item FNE-empr. (5 anos) \hfill +132,23 p.p.\textsuperscript{*}
		\item FNE-massa salarial \hfill (+)\textsuperscript{*}
		\item FNE-salário \hfill NS\phantom{\textsuperscript{*}}
	}
	\\

	\hline
\textit{Generalized Propensity Score}
	& \citeonline{OliveiraTerra2018}
	& \fcvar[2.3cm]{\item Empresas GO \item 2004-2011}
	& \fcvar[3.0cm]{\item FCO-Empresarial (contínuo) \item Cresc. emprego \item Cresc. salário}
	& \fc[4.2cm]{
		\item FCO-emprego \hfill (+)\textsuperscript{*}
		\item Efeito limitado a intervalo específico da dose
		\item FCO-salário \hfill NS\phantom{\textsuperscript{*}}
	}
	\\

	\hline
IV e Função Dose-Resposta
	& \citeonline{Cunha2024}
	& \fcvar[2.3cm]{\item Empresas NE \item 2000-2006}
	& \fcvar[3.0cm]{\item FNE (contínua) \item Cresc. emprego \item Cresc. massa salarial}
	& \fc[4.2cm]{
		\item Dose-resposta emprego: IV \hfill 357,54\%\textsuperscript{*}
		\item Dose-resposta emprego: OLS \hfill 133,94\%\textsuperscript{*}
		\item Dose-resposta massa: IV \hfill 220,96\%\textsuperscript{*}
		\item Dose-resposta massa: OLS \hfill 188,55\%\textsuperscript{*}
	}
	\\

	\hline
Primeira Diferença
	& \citeonline{Resende2014c}
	& \fcvar[2.3cm]{\item Empresas NE \item 1999-2006}
	& \fcvar[3.0cm]{\item FNE-Indústria (\textit{dummy}) \item Cresc. emprego}
	& \fc[4.2cm]{
		\item FNE-Ind. (anual) \hfill +10,83 a 16,11 p.p.\textsuperscript{*}
		\item FNE-Ind. (3 anos) \hfill +30 a 56 p.p.\textsuperscript{*}
	}
	\\

	\hline
Painel Espacial (SAR)
	& \citeonline{Oliveira2015}
	& \fcvar[2.3cm]{\item Municípios GO \item 2004-2011}
	& \fcvar[3.0cm]{\item Log. FCO \item Log. emprego \item Salário médio}
	& \fc[4.2cm]{
		\item FCO-emprego \hfill (+)\textsuperscript{*}
		\item FCO-salário \hfill (+)\textsuperscript{*}
		\item Sem diferença relevante vs. painel clássico
	}
	\\

	\hline
Painel Dinâmico
	& \citeonline{Ribeiro2024}
	& \fcvar[2.3cm]{\item Municípios NE \item 2010-2015}
	& \fcvar[3.0cm]{\item Log. FNE \item Log. emprego por setor}
	& \fc[4.2cm]{
		\item FNE-primário \hfill +1,10\%\textsuperscript{*}
		\item FNE-terciário \hfill +0,53\%\textsuperscript{*}
		\item FNE-secundário \hfill +0,13\%\textsuperscript{*}
		\item FNE defasado (secundário) \hfill +0,06\%\textsuperscript{*}
	}
	\\

	\hline
Equilíbrio Geral
	& \citeonline{Ribeiro2020}
	& \fcvar[2.3cm]{\item Região NE \item 2014-2025}
	& \fcvar[3.0cm]{\item FNE \item Emprego total}
	& \fc[4.2cm]{
		\item FNE (emprego NE acumulado) \hfill +2,75\%\textsuperscript{*}
		\item NE-agropecuária \hfill +8,0\%\textsuperscript{*}
		\item NE-construção \hfill +5,3\%\textsuperscript{*}
	}
	\\

	\hline
Diferenças em Diferenças
	& \citeonline{Oliveira2020}
	& \fcvar[2.3cm]{\item Empresas N (FNO) \item 2000-2010}
	& \fcvar[3.0cm]{\item FNO (\textit{dummy}) \item Emprego \item Massa salarial \item Salário médio}
	& \fc[4.2cm]{
		\item Emprego e renda \hfill (+)\textsuperscript{*}
		\item Curto e médio prazo (+), exceto grandes empresas
		\item Investimento: efeito a longo prazo
		\item Giro/custeio: efeito a curto prazo
	}
	\\

	\hline
RDD Geográfico
	& \citeonline{Oliveira2021}
	& \fcvar[2.3cm]{\item Municípios NE \item 1970-2010}
	& \fcvar[3.0cm]{\item FNE \item Emprego formal e informal (2010)}
	& \fc[4.2cm]{
		\item Agropecuária \hfill +5,4\%\textsuperscript{*}
		\item Indústria \hfill +2,8\%\textsuperscript{*}
		\item Serviços \hfill +3,5\%\textsuperscript{*}
		\item Taxa de informalidade \hfill NS\phantom{\textsuperscript{*}}
	}
	\\

\end{longtable}%
\normalsize%
\setlength{\extrarowheight}{0pt}%
\setlength{\tabcolsep}{6pt}%
\linespread{1.5}\selectfont%
\renewcommand{\tablename}{Tabela}%
\renewcommand{\thetable}{\arabic{table}}%
}


\citeonline{Silva2009} foi o primeiro estudo a investigar o impacto dos empréstimos dos FCs sobre a contratação de trabalhadores nas empresas beneficiadas. O estudo faz a interseção dos registros de empresas beneficiadas com os dados da RAIS (Relação Anual de Informações Sociais), obtendo a quantidade de vínculos de empregos formais de 449 empresas (NE=211, N=174, CO=74) nos anos de 2000 e 2003. Aplicando o método \textit{Propensity Score Matching (PSM)}, os autores concluem que nas empresas beneficiadas pelo FNE a quantidade de empregos cresceu em média 61,2 pontos percentuais a mais do que nas empresas não beneficiadas. No entanto, os valores não foram significativos para o salário médio. E ainda, para o FNO e o FCO, não foi obtido efeito positivo em nenhum dos indicadores avaliados.

\citeonline{Soares2017} também aplicam o método PSM, mas fazem uso de um conjunto mais rico de informações sobre o FNE, o que permitiu ampliar consideravelmente o número de empresas e o horizonte de tempo investigado. Além disso, os autores fazem uma análise ano a ano dos efeitos do fundo, o que permite estimativas mais eficientes, com quantidades maiores de empresas nos primeiros anos pós-tratamento. Com dados do período de 2000 a 2006, no horizonte de 3 anos pós-tratamento há 1.104 empresas e em 5 anos há 314. Os resultados obtidos reforçam as evidências obtidas por \citeonline{Silva2009}. No primeiro ano pós-tratamento o crescimento do estoque de empregos é em média 7,96 pontos percentuais maior nas empresas beneficiadas, relativamente às demais. Além disso, a trajetória é crescente nos anos subsequentes, alcançando 33,59 pontos percentuais no terceiro ano e 132,23 no quinto ano. Se confirmados, esses resultados atribuem grande relevância para a política regional, indicando que o FNE tem papel importante na geração de empregos formais nas empresas beneficiadas. Os autores também destacam que o efeito de maior monta estaria associado às empresas financiadas entre 1999 e 2001, as únicas no horizonte de 5 anos pós-tratamento, estando mais concentradas no setor industrial. O mesmo estudo ainda analisa possíveis efeitos do FNE sobre a massa salarial, obtendo resultados significativos e com dinâmica semelhante ao do efeito no emprego. Como consequência, a variável salário médio não apresentou variação significativa como decorrência do FNE, indicando que as novas contratações teriam ocorrido ao mesmo nível dos salários vigentes.

\citeonline{Oliveira2018} apresentam como inovação principal o uso de variável de tratamento do tipo contínua, permitindo analisar possíveis heterogeneidades no efeito do fundo decorrentes da intensidade do tratamento (valor do empréstimo), em contraste à abordagem dicotômica de tratamento presente nos estudos anteriores. Os autores usam o método \textit{Generalized Propensity Score Matching (GPS)} e dados de financiamentos a empresas do estado de Goiás realizados pelo FCO por meio da linha Empresarial (associada à indústria, comércio e serviços) no período de 2004 a 2011. Os resultados indicam relação positiva e significativa com o crescimento do emprego apenas em um intervalo específico da distribuição de tratamento e limitada ao período 2004-2008 (pré-crise), cenário semelhante ao obtido na abordagem PSM com variável de tratamento dicotômica. Com respeito ao salário médio, não se obteve nenhum resultado significativo. \citeonline{Oliveira2018} expandem a abordagem de \citeonline{Oliveira2018} para os demais fundos constitucionais. O dado mais relevante obtido indica efeito decrescente dos fundos sobre as variáveis de interesse, o que significa que quanto maior o valor do empréstimo, menor o efeito sobre o crescimento do emprego e do salário. Para os autores, isso se justifica do ponto de vista teórico pela propriedade de retornos marginais decrescentes da função de produção neoclássica.

\citeonline{Junior2024} contribuem com a literatura ao usar o método de função dose-resposta para captar efeitos heterogêneos e correção de viés de seleção por meio de variável instrumental. Para garantir comparabilidade, os autores optam por usar a mesma base de dados de \citeonline{Soares2017}. Na variável crescimento do emprego, o modelo IV aponta efeito médio de 357,54\%, acima de 133,94\% do modelo OLS. Para o crescimento da massa salarial, o valor médio é 220,96\% no modelo IV contra 188,55\% em OLS. Em ambas as variáveis a restrição de exogeneidade leva a subestimação do efeito médio e superestimação do ponto ótimo de tratamento, sinalizando para efeitos mais significativos em pontos mais baixos da distribuição de tratamento, quando a endogeneidade é corrigida.

\citeonline{Resende2014c} contribui para uma limitação importante dos métodos PSM e GSM: a incapacidade de controlar por efeitos não observáveis que possam influenciar tanto a probabilidade de tratamento quanto seu resultado (apenas características observáveis). Usando abordagem de primeira diferença e dados de 1999 a 2006, o autor identifica efeito positivo do FNE-Indústria sobre o crescimento do número de empregos formais nas empresas beneficiadas, com magnitude entre 10,83 e 16,11 pontos percentuais ao ano. No horizonte de três anos pós-tratamento, o valor é de 30 a 56 pontos percentuais a mais, de certa forma consistente com os resultados de \citeonline{Silva2009} e \citeonline{Soares2017}.

\citeonline{Oliveira2015} é o primeiro estudo a avaliar o efeito dos fundos constitucionais sobre variáveis agregadas do mercado de trabalho. O estudo adota os municípios de Goiás como unidade de análise e aplica o método de painel espacial autorregressivo (SAR), a fim de controlar também por características da vizinhança do município. Os resultados indicam efeitos positivos sobre emprego e salário, porém, sem diferenças relevantes em relação ao modelo de painel clássico com efeitos fixos, o que pode ser decorrência de baixa robustez dos resultados. \citeonline{Ribeiro2024} trazem duas contribuições principais: análise dos efeitos agregados do FNE por setor de atividade econômica e o uso de modelo de painel dinâmico para amenizar possível viés de endogeneidade. Os resultados apontam relação positiva entre os investimentos do FNE e o nível de emprego no município, com maior efeito sobre o setor primário (1,1\% de aumento no emprego para 1\% de aumento no investimento do FNE), seguido por terciário (0,53\%) e secundário (0,13\%). Esses resultados são observados para os investimentos correntes do FNE. Para o FNE defasado em um ano, o valor é significativo somente para o setor secundário e em menor magnitude (0,06\%). Sugere-se assim que o FNE tem efeito positivo e imediato sobre o emprego, especialmente no setor primário, setor mais intensivo em mão de obra.

Em maior nível de agregação, \citeonline{Ribeiroetal2020} aplicam um modelo dinâmico interregional de equilíbrio geral computável (CGE) para estimar os efeitos dos investimentos FNE sobre o emprego total da região Nordeste, dos estados e dos setores de atividade. As evidências apontam que os investimentos do FNE nos anos 2014 e 2015 estão associados a uma variação acumulada de 2,75\% na quantidade de empregos na região Nordeste até o ano de 2025, valor pouco menor que o estimado para a variação do PIB. Na análise por estados, Piauí, Ceará e Rio Grande do Norte também apresentam a maior variação acumulada do emprego, entre 3,19\% e 4,31\%, e Maranhão e Bahia têm as menores variações, de 0,45\% a 2,16\%. No aspecto setorial, também não há perdedores, com agropecuária (8\%) e construção (5,3\%) os setores com maior variação do emprego.

\citeonline{Oliveira2020} empregam o método diferenças-em-diferenças, que controla pela heterogeneidade não observável das empresas e possibilita uma estimativa do efeito causal do fundo. Além disso, têm como diferencial a análise dos efeitos microeconômicos do FNO por porte da empresa beneficiada, tipo de crédito concedido e horizonte de tempo. No geral, empresas beneficiadas apresentaram desempenho superior às não beneficiadas, em curto e médio prazo, em termos de geração de emprego e renda. No entanto, médias e grandes empresas beneficiadas por financiamento a capital de giro e custeio apresentam desempenho similar ao grupo de controle, no que diz respeito à geração de emprego e renda. Além disso, os autores apontam que o impacto positivo do crédito com finalidade capital de giro e custeio é em média maior no curto prazo do que no longo prazo. Por outro lado, em empresas beneficiadas pela finalidade investimento, os impactos são mais observáveis a longo prazo. Os efeitos assimétricos indicam que a finalidade investimento promove a acumulação de capital, o que aumenta a produtividade do trabalho a longo prazo, pois o impacto é positivo para o trabalho e a renda (salário médio).

\citeonline{Oliveira2021} utilizam modelo de descontinuidade geográfica (RDD) para estimar o impacto causal do FNE sobre o estoque de emprego, distinguindo-se dos estudos anteriores pela robustez metodológica e por analisar impacto sobre o emprego formal e informal, decorrência do uso de dados censitários. Os resultados indicam que o acesso ao FNE gerou impacto positivo e significativo sobre o estoque de empregos formais e informais no ano 2010, com magnitude 5,4\% (agropecuária), 2,8\% (indústria) e 3,5\% (serviços). Porém, impacto sobre a taxa de informalidade foi nulo.

Em síntese, as avaliações sobre o mercado de trabalho apresentam maior convergência do que as relativas ao PIB: há evidências consistentes de efeito positivo do FNE sobre a geração de emprego formal nas empresas beneficiadas, com magnitude crescente ao longo dos primeiros anos pós-tratamento. Os estudos também convergem em apontar efeitos mais pronunciados em micro e pequenas empresas e na modalidade de crédito para investimento. Por outro lado, a variável salário médio não apresenta variação significativa na maioria das especificações, sugerindo que os novos postos de trabalho são criados ao nível salarial vigente, sem ganhos adicionais de produtividade. Os resultados para FNO e FCO são mais limitados e inconclusivos, refletindo menor cobertura na literatura.

\subsection{Avaliações de Impacto dos Fundos de Desenvolvimento}
\label{subsec:fd}

Os Fundos de Desenvolvimento (FDNE, FDA e FDCO) constituem instrumentos mais recentes da PNDR, com primeiras contratações a partir de 2007 no caso do FDNE e de 2014 para o FDCO, o que explica a menor quantidade de avaliações empíricas em comparação aos Fundos Constitucionais. A escala dos empreendimentos financiados --- mínimo de R\$ 15 milhões, com concentração em infraestrutura e indústria --- sugere mecanismo de impacto distinto, cuja avaliação é relevante para informar a alocação de recursos entre os instrumentos da política. Dentre os estudos aprovados nesta revisão, 8 avaliam algum aspecto dos Fundos de Desenvolvimento, com concentração expressiva no FDNE. Não foram identificadas avaliações quantitativas dedicadas ao FDA nem ao FDCO, lacuna relevante dada a magnitude dos recursos envolvidos --- R\$ 5,5 bilhões e R\$ 1,8 bilhão contratados, respectivamente. As subseções seguintes organizam as evidências por abordagem metodológica. O Quadro~\ref{tab:estudos_fd} apresenta a síntese dos estudos que avaliam os Fundos de Desenvolvimento.

% Arquivo: latex/tabelas/survey_artigos_fd.tex
% Quadro de estudos sobre impacto dos Fundos de Desenvolvimento
% Fonte: Revisão sistemática - dados extraídos do pipeline pndr_survey

% IMPORTANTE: Este arquivo requer:
% - Pacotes: longtable, multirow, afterpage, enumitem, array
% - Macros: \fc e \metodo (definidas em tabelas/tabela_macros.tex)
% - Comando \citeonline (do pacote abntex2cite)
% - Contador: quadro (definido no main.tex)

\afterpage{\clearpage%
\tiny%
\setlength{\extrarowheight}{1pt}%
\setlength{\tabcolsep}{2pt}%
\linespread{0.9}\selectfont%
\renewcommand{\arraystretch}{1.2}%
\renewcommand{\tablename}{\quadroname}%
\renewcommand{\thetable}{\arabic{quadro}}%
\stepcounter{quadro}%
\begin{longtable}{|>{\raggedright\arraybackslash}p{1.6cm}
		|>{\raggedright\arraybackslash}p{1.7cm}
		|>{\raggedright\arraybackslash}p{2cm}
		|>{\raggedright\arraybackslash}p{1.5cm}
		|>{\raggedright\arraybackslash}p{2.6cm}
		|>{\raggedright\arraybackslash}p{4.2cm}|}

	\caption{Estudos de Impacto dos Fundos de Desenvolvimento} \label{tab:estudos_fd} \\

\hline
	\begin{tabular}[c]{@{}>{\centering\arraybackslash}m{1.6cm}@{}}Método\end{tabular}
	& \begin{tabular}[c]{@{}>{\centering\arraybackslash}m{1.7cm}@{}}Artigo\end{tabular}
	& \begin{tabular}[c]{@{}>{\centering\arraybackslash}m{2cm}@{}}Amostragem\end{tabular}
	& \begin{tabular}[c]{@{}>{\centering\arraybackslash}m{1.5cm}@{}}Variável Independente\end{tabular}
	& \begin{tabular}[c]{@{}>{\centering\arraybackslash}m{2.6cm}@{}}Variável Dependente\end{tabular}
	& \begin{tabular}[c]{@{}>{\centering\arraybackslash}m{4.2cm}@{}}Resultado\end{tabular} \\
	\hline
	\endfirsthead

	\multicolumn{6}{c}%
	{{\quadroname\ \thetable{} -- Continuação}} \\
	\hline
	\begin{tabular}[c]{@{}>{\centering\arraybackslash}m{1.6cm}@{}}Método\end{tabular}
	& \begin{tabular}[c]{@{}>{\centering\arraybackslash}m{1.7cm}@{}}Artigo\end{tabular}
	& \begin{tabular}[c]{@{}>{\centering\arraybackslash}m{2cm}@{}}Amostragem\end{tabular}
	& \begin{tabular}[c]{@{}>{\centering\arraybackslash}m{1.5cm}@{}}Variável Independente\end{tabular}
	& \begin{tabular}[c]{@{}>{\centering\arraybackslash}m{2.6cm}@{}}Variável Dependente\end{tabular}
	& \begin{tabular}[c]{@{}>{\centering\arraybackslash}m{4.2cm}@{}}Resultado\end{tabular} \\
	\hline
	\endhead

	\hline
	\multicolumn{6}{r}{{Continua}} \\
	\endfoot

	\hline
	\multicolumn{6}{l}{
		\tiny
		\parbox{14.8cm}{%
			\vspace{0.1cm}
			Fonte: Elaborado pelo autor. Siglas: AC - Apresentação em congresso; NS - Não significativo.
	}}
	\endlastfoot

	\hline
Painel Espacial & \citeonline{FerreiraIrffiCarneiro2024}
	& \fc[2cm]{\item NE \item 2010-2021}
	& \fc[1.5cm]{\item FDNE}
	& \fc[2cm]{\item IDM}
	& \fc[4.5cm]{
		\item IDM renda \hfill 0,019\textsuperscript{*}
		\item IDM, IDM saúde e IDM educação \hfill NS\phantom{\textsuperscript{*}}
	} \\

	\cline{2-6}

	& \citeonline{gondim2025}
	& \fc[2cm]{\item NE \item 2010-2021}
	& \fc[1.5cm]{\item FDNE (\textit{dummy})}
	& \fc[2cm]{\item $\Delta$ ln(PIB pc)}
	& \fc[4.5cm]{
		\item Efeito direto \hfill NS\phantom{\textsuperscript{*}}
		\item Efeito indireto \hfill 0,0252\textsuperscript{*}
		\item Efeito total \hfill 0,0351\textsuperscript{*}
	} \\
	\cline{5-6}
	& & &
	& \fc[2cm]{\item $\Delta$ ln(VAB pc)}
	& \fc[4.5cm]{
		\item Efeito direto \hfill 0,0409\textsuperscript{*}
		\item Efeito indireto \hfill NS\phantom{\textsuperscript{*}}
		\item Efeito total \hfill 0,6014\textsuperscript{*}
	} \\

	\hline
DID Dois Estágios & \citeonline{Carneiroetal2024a}
	& \fc[2cm]{\item NE \item 2002-2019}
	& \fc[1.5cm]{\item FDNE}
	& \fc[2cm]{\item log PIB\textit{pc}, VAB setorial}
	& \fc[4.5cm]{
		\item PIB\textit{pc} \hfill 0,063\textsuperscript{*}
		\item VAB\textit{pc} serviços \hfill 0,111\textsuperscript{*}
		\item Arrecadação \hfill 0,207\textsuperscript{*}
		\item VAB ind. e VAB agro. \hfill NS\phantom{\textsuperscript{*}}
	} \\

	\hline
DID Escalonado
	& \citeonline{BrazBastosIrffi2024}
	& \fc[2cm]{\item NE \item 2002-2021}
	& \fc[1.5cm]{\item FDNE}
	& \fc[2cm]{\item PIB\textit{pc}, emprego, renda, IDEB}
	& \fc[4.5cm]{
		\item PIB\textit{pc} \hfill 24,07\%\textsuperscript{*}
		\item Renda \hfill 4,63\%\textsuperscript{*}
		\item IDEB anos iniciais \hfill 18,67\%\textsuperscript{*}
		\item IDEB anos finais \hfill 21,49\%\textsuperscript{*}
		\item Emprego, pobreza, mortalidade infantil e saúde \hfill NS\phantom{\textsuperscript{*}}
	} \\
	\cline{2-6}

	& \citeonline{Irffietal2025}
	& \fc[2cm]{\item NE \item 1999-2022}
	& \fc[1.5cm]{\item FDNE-parques eólicos}
	& \fc[2cm]{\item PIB\textit{pc}, VAB, emprego, massa salarial}
	& \fc[4.5cm]{
		\item PIB\textit{pc} \hfill 18,97\%\textsuperscript{*}
		\item VAB indústria \hfill 71,90\%\textsuperscript{*}
		\item VAB serviço \hfill 10,24\%\textsuperscript{*}
		\item Emprego \hfill 23,41\%\textsuperscript{*}
		\item Redução temporária do VAB-agro durante construção
	} \\

	\hline
Painel Dinâmico GMM
	& \citeonline{Souza2025}
	& \fc[2cm]{\item N, NE, CO \item 2002-2021}
	& \fc[1.5cm]{\item FDNE, FDA, FDCO (+ FC + IF)}
	& \fc[2cm]{\item PIB\textit{pc}, emprego, renda}
	& \fc[4.5cm]{
		\item Heterogeneidade entre fundos, regiões e setores
		\item FDCO inconclusivo (poucos projetos)
	} \\

\end{longtable}%
\normalsize%
\setlength{\extrarowheight}{0pt}%
\setlength{\tabcolsep}{6pt}%
\linespread{1.5}\selectfont%
\renewcommand{\tablename}{Tabela}%
\renewcommand{\thetable}{\arabic{table}}%
}


\subsubsection{Modelos Quase Experimentais}

O método de diferenças em diferenças (DID) e suas extensões --- como o DID escalonado, adequado para intervenções implementadas em momentos distintos para diferentes unidades --- constituem a principal abordagem empregada para avaliar os Fundos de Desenvolvimento. Os três estudos discutidos a seguir representam as primeiras evidências empíricas de impacto causal desses fundos sobre indicadores municipais.

\citeonline{Braz2024} aplicam modelo DID escalonado aos municípios da área de atuação da SUDENE no período 2002--2021 e estimam que aqueles que receberam empreendimentos apoiados pelo FDNE experimentam elevação média de 24\% no PIB \textit{per capita} em relação ao grupo de controle. O estudo inova ao avaliar variáveis de resultado não estritamente econômicas: a renda domiciliar \textit{per capita} apresenta incremento de 4,6\%, e os indicadores de educação básica (IDEB) registram elevação de 19\% nos anos iniciais e 21\% nos anos finais. Não se obteve significância estatística para emprego formal, pobreza, mortalidade infantil e indicadores de saúde.

\citeonline{Carneiroetal2024a} avaliam simultaneamente os efeitos do FNE, do FDNE e dos incentivos fiscais da SUDENE sobre a economia dos municípios nordestinos, empregando estimadores DID em dois estágios e DID escalonado para o período 2002--2019. Para o FDNE, identificam impacto positivo e significativo sobre o PIB \textit{per capita} municipal, com coeficiente concentrado no setor de serviços e na arrecadação tributária. Os incentivos fiscais não apresentaram significância estatística, e a interação entre os instrumentos tampouco indicou sinergia, sugerindo que os efeitos operam de forma independente. Os resultados desse estudo para o FNE foram discutidos na Subseção~\ref{subsec:fc-pib}.

\citeonline{Irffietal2025} estimam o impacto da construção de parques eólicos sobre indicadores econômicos dos municípios nordestinos, empregando o método de controle sintético generalizado e estudos de eventos para o período 1999--2022. Os resultados apontam elevação de 19\% no PIB \textit{per capita}, de 72\% no valor adicionado bruto (VAB) industrial e de 10\% no VAB de serviços nos municípios com parques em operação ou construção. Identificam-se ainda incrementos de 23\% no emprego formal e de 31\% na massa salarial, além de redução temporária do VAB agropecuário durante a fase de construção, sem efeito permanente. Embora o estudo avalie todos os parques eólicos da região --- e não exclusivamente os financiados pelo FDNE ---, parcela expressiva dos empreendimentos de energia eólica no Nordeste contou com recursos desse fundo, e a proximidade dos valores obtidos com os de \citeonline{Braz2024} reforça a consistência das conclusões.

As três avaliações convergem em apontar impacto positivo do FDNE sobre o PIB \textit{per capita}, com magnitudes situadas entre 19\% e 24\%, consideravelmente superiores às estimadas para os Fundos Constitucionais --- cujos coeficientes situam-se tipicamente entre 0,03 e 0,13 p.p. de incremento no crescimento anual para cada 1 p.p. de aumento na relação Fundo/PIB. Essa diferença pode estar associada à escala dos empreendimentos financiados, com valor mínimo de R\$ 15 milhões, à concentração setorial em infraestrutura e indústria de transformação e ao horizonte mais longo de maturação desses investimentos. Cabe ressalvar que as três avaliações empregam variantes do método DID e analisam períodos parcialmente sobrepostos, o que limita a independência das estimativas.

\subsubsection{Modelos Espaciais}

Modelos de painel espacial permitem decompor o impacto de uma política em componente direto --- sobre o próprio município beneficiado --- e indireto --- sobre municípios vizinhos. Essa distinção é particularmente relevante para os Fundos de Desenvolvimento, cujos empreendimentos financiados (infraestrutura de energia, transportes, complexos industriais) podem gerar externalidades que se dispersam para além dos limites municipais.

\citeonline{FerreiraIrffiCarneiro2024} utilizam modelo SAC (\textit{Spatial Autoregressive Combined}) estático para avaliar o FDNE e os incentivos fiscais da SUDENE nos municípios nordestinos entre 2010 e 2021. A variável de resultado adotada é o Índice de Desenvolvimento Municipal (IDM) e suas componentes. Para o FDNE, o estudo identifica efeito positivo e significativo apenas sobre a componente IDM-renda (coeficiente 0,019), sem evidência sobre o IDM-geral, IDM-saúde ou IDM-educação. A adoção de variável de resultado distinta do PIB \textit{per capita} dificulta a comparação direta com os demais estudos desta subseção, embora contribua para ampliar o escopo das dimensões avaliadas.

\citeonline{gondim2025} analisam conjuntamente FNE, FDNE e incentivos fiscais da SUDENE por meio de modelo espacial generalizado com dados de 2010 a 2021. Os resultados para o FDNE revelam padrão distinto dos demais instrumentos: o efeito direto sobre o PIB \textit{per capita} é não significativo, enquanto o efeito indireto (transbordamento para municípios vizinhos) apresenta coeficiente positivo e significativo de 0,025, resultando em efeito total de 0,035. Para o VAB \textit{per capita}, o efeito direto é significativo (0,041), mas o indireto não. Esse resultado sugere que os investimentos do FDNE beneficiam não apenas o município-sede, mas também localidades vizinhas, padrão consistente com a natureza dos projetos financiados --- infraestrutura de energia e transportes cuja área de influência transcende limites municipais. Tal configuração contrasta com os Fundos Constitucionais, cujos efeitos diretos tendem a ser mais pronunciados.

\subsubsection{Avaliação Integrada e Síntese}

\citeonline{Souza2025} adotam painel dinâmico GMM com defasagens distribuídas para avaliar simultaneamente todos os instrumentos da PNDR --- incluindo FDNE, FDA e FDCO --- no período 2002--2021. O estudo constitui a única avaliação que contempla os três fundos de desenvolvimento em uma mesma especificação. Os resultados indicam heterogeneidade expressiva entre fundos, regiões e setores, com evidências inconclusivas para o FDCO, atribuídas ao reduzido número de projetos e ao curto período de operação.

Em conjunto, as evidências disponíveis apontam para impacto positivo expressivo do FDNE sobre o PIB \textit{per capita} municipal, com magnitudes entre 19\% e 24\% nas avaliações quase experimentais e presença de efeitos de transbordamento espacial nos modelos de painel. A base empírica, porém, apresenta limitações relevantes: concentra-se em estudos recentes (2024--2026) com períodos amostrais parcialmente sobrepostos e predominância do DID escalonado; não dispõe de avaliação quantitativa dedicada ao FDA nem ao FDCO; e restringe-se majoritariamente ao PIB \textit{per capita}, com evidências incipientes sobre emprego, renda domiciliar e pobreza. A diversificação metodológica e a extensão da cobertura para os três fundos de desenvolvimento constituem prioridades para a agenda de pesquisa.

\subsection{Avaliações de Impacto dos Incentivos Fiscais}
\label{subsec:if}

A avaliação empírica dos incentivos fiscais da SUDENE e da SUDAM constitui a vertente menos desenvolvida da literatura sobre a PNDR, a despeito do expressivo volume de renúncia tributária envolvido. Dentre os 34 estudos aprovados nesta revisão, 9 avaliam explicitamente esses instrumentos, refletindo a operação mais recente dos programas --- com concessões sistemáticas a partir de 2010 --- e as dificuldades de acesso a dados detalhados de renúncia fiscal em nível de empresa. Diferentemente dos Fundos Constitucionais, cujo mecanismo de transmissão opera via crédito subsidiado, os incentivos fiscais atuam pela redução de 75\% do Imposto de Renda de Pessoa Jurídica (IRPJ), diminuindo o custo tributário de operação em municípios da área de atuação da superintendência. Todos os estudos dedicados concentram-se nos benefícios concedidos pela SUDENE, não havendo avaliações específicas da SUDAM. Dois trabalhos avaliam adicionalmente o Prodepe --- programa estadual de Pernambuco --- em interação com os instrumentos federais. O Quadro~\ref{tab:estudos_if} apresenta a síntese dos estudos desta subseção.

% Arquivo: latex/tabelas/survey_artigos_if.tex
% Quadro de estudos sobre impacto dos Incentivos Fiscais
% Fonte: Revisão sistemática - dados extraídos do pipeline pndr_survey
% Adaptado da tese: removido CostaCarneiroIrffi2023 (não aprovado),
%   adicionados Oliveira2020a, Alves2024, Gondim2025, Souza2025.
%   Chaves BibTeX corrigidas para consistência com references.bib.

% IMPORTANTE: Este arquivo requer:
% - Pacotes: longtable, multirow, afterpage, enumitem, array
% - Macros: \fc e \metodo (definidas em tabelas/tabela_macros.tex)
% - Comando \citeonline (do pacote abntex2cite)
% - Contador: quadro (definido no main.tex)

\afterpage{\clearpage%
\tiny%
\setlength{\extrarowheight}{1pt}%
\setlength{\tabcolsep}{2pt}%
\linespread{0.9}\selectfont%
\renewcommand{\arraystretch}{1.2}%
\renewcommand{\tablename}{\quadroname}%
\renewcommand{\thetable}{\arabic{quadro}}%
\stepcounter{quadro}%
\begin{longtable}{|>{\raggedright\arraybackslash}p{1.4cm}
		|>{\raggedright\arraybackslash}p{1.5cm}
		|>{\raggedright\arraybackslash}p{1.5cm}
		|>{\raggedright\arraybackslash}p{2cm}
		|>{\raggedright\arraybackslash}p{1.5cm}
		|>{\raggedright\arraybackslash}p{2cm}
		|>{\raggedright\arraybackslash}p{4.2cm}|}

	\caption{Estudos de Impacto dos Incentivos Fiscais} \label{tab:estudos_if} \\

	\hline
	\begin{tabular}[c]{@{}>{\centering\arraybackslash}m{1.4cm}@{}}Método\end{tabular}
	& \begin{tabular}[c]{@{}>{\centering\arraybackslash}m{1.5cm}@{}}Artigo\end{tabular}
	& \begin{tabular}[c]{@{}>{\centering\arraybackslash}m{1.5cm}@{}}Publicação\end{tabular}
	& \begin{tabular}[c]{@{}>{\centering\arraybackslash}m{2cm}@{}}Amostragem\end{tabular}
	& \begin{tabular}[c]{@{}>{\centering\arraybackslash}m{1.5cm}@{}}Variável Independente\end{tabular}
	& \begin{tabular}[c]{@{}>{\centering\arraybackslash}m{2cm}@{}}Variável Dependente\end{tabular}
	& \begin{tabular}[c]{@{}>{\centering\arraybackslash}m{4.2cm}@{}}Resultado\end{tabular} \\
	\hline
	\endfirsthead

	\multicolumn{7}{c}%
	{{\quadroname\ \thetable{} -- Continuação}} \\
	\hline
	\begin{tabular}[c]{@{}>{\centering\arraybackslash}m{1.4cm}@{}}Método\end{tabular}
	& \begin{tabular}[c]{@{}>{\centering\arraybackslash}m{1.5cm}@{}}Artigo\end{tabular}
	& \begin{tabular}[c]{@{}>{\centering\arraybackslash}m{1.5cm}@{}}Publicação\end{tabular}
	& \begin{tabular}[c]{@{}>{\centering\arraybackslash}m{2cm}@{}}Amostragem\end{tabular}
	& \begin{tabular}[c]{@{}>{\centering\arraybackslash}m{1.5cm}@{}}Variável Independente\end{tabular}
	& \begin{tabular}[c]{@{}>{\centering\arraybackslash}m{2cm}@{}}Variável Dependente\end{tabular}
	& \begin{tabular}[c]{@{}>{\centering\arraybackslash}m{4.2cm}@{}}Resultado\end{tabular} \\
	\hline
	\endhead

	\hline
	\multicolumn{7}{r}{{Continua}} \\
	\endfoot

	\hline
	\multicolumn{7}{l}{
		\tiny
		\parbox{14.8cm}{%
			\vspace{0.1cm}
			Fonte: Elaborado pelo autor. Siglas: AC - Apresentação em congresso; NS - Não significativo.
	}}
	\endlastfoot

	\hline
DID & \citeonline{Garsous2017}
	& \textit{World Development}
	& \fc[2cm]{\item BA, RJ, ES, MG \item 1998-2009}
	& \fc[1.5cm]{\item IF-SUDENE (turismo)}
	& \fc[2cm]{\item Emprego no turismo (log)}
	& \fc[4.2cm]{
		\item Emprego turismo \hfill 20\%\textsuperscript{*}, 30\%\textsuperscript{*}
	} \\

	\cline{2-7}
	& \citeonline{Oliveira2020a}
	& AC
	& \fc[2cm]{\item PE \item 2000-2017}
	& \fc[1.5cm]{\item Prodepe + IF-SUDENE}
	& \fc[2cm]{\item Emprego, salário médio, massa salarial}
	& \fc[4.2cm]{
		\item Emprego \hfill 8,6\%\textsuperscript{*}
		\item Salário médio \hfill $-$10,3\%\textsuperscript{*}
		\item Massa salarial \hfill NS\phantom{\textsuperscript{*}}
	} \\

	\hline
DID Escalonado & \citeonline{Braz2023}
	& AC
	& \fc[2cm]{\item SUDENE \item 2002-2021}
	& \fc[1.5cm]{\item IF-SUDENE}
	& \fc[2cm]{\item Emprego, renda}
	& \fc[4.2cm]{
		\item Emprego \hfill 3,2\%\textsuperscript{*}
		\item Renda \hfill 1,2\%\textsuperscript{*}
		\item Efeito crescente e duradouro
	} \\

	\cline{2-7}
	& \citeonline{Alves2024}
	& AC
	& \fc[2cm]{\item PE \item 2000-2017}
	& \fc[1.5cm]{\item Prodepe + IF-SUDENE}
	& \fc[2cm]{\item Emprego, salário médio, massa salarial}
	& \fc[4.2cm]{
		\item Emprego \hfill 22,3\%\textsuperscript{*}
		\item Massa salarial \hfill 15,1\%\textsuperscript{*}
		\item Salário médio \hfill $-$8,2\%\textsuperscript{*}
		\item Efeito temporário e declinante
	} \\

	\hline
DEA + SFA & \citeonline{Carneiro2023a}
	& Rev. Cadernos de Finanças Públicas
	& \fc[2cm]{\item NE \item 2011-2019}
	& \fc[1.5cm]{\item IF-SUDENE}
	& \fc[2cm]{\item Eficiência produtiva (massa salarial)}
	& \fc[4.2cm]{
		\item Empresas eficientes no Semiárido
		\item Intensivas em mão de obra
		\item Ineficiência associada ao setor, não à empresa
	} \\

	\hline
DID Dois Estágios & \citeonline{Carneiroetal2024a}
	& AC
	& \fc[2cm]{\item NE, MG, ES \item 2002-2019}
	& \fc[1.5cm]{\item IF-SUDENE (+ FNE, FDNE)}
	& \fc[2cm]{\item log PIB\textit{pc}, VAB setorial}
	& \fc[4.2cm]{
		\item IF-SUDENE \hfill NS\phantom{\textsuperscript{*}}
	} \\

	\hline
Controle Sintético & \citeonline{Costa2024}
	& \textit{Applied Economics Letters}
	& \fc[2cm]{\item NE \item 2006-2019}
	& \fc[1.5cm]{\item IF-SUDENE}
	& \fc[2cm]{\item Vínculos (log)}
	& \fc[4.2cm]{
		\item Todos os setores \hfill 19,6\%\textsuperscript{*}
		\item Indústria \hfill 18,3\%\textsuperscript{*}
		\item Efeitos por até 5 anos
		\item Demais setores \hfill NS\phantom{\textsuperscript{*}}
	} \\

	\hline
Painel Espacial & \citeonline{FerreiraIrffiCarneiro2024}
	& AC
	& \fc[2cm]{\item NE \item 2010-2021}
	& \fc[1.5cm]{\item IF-SUDENE (+ FDNE)}
	& \fc[2cm]{\item IDM e subcomponentes}
	& \fc[4.2cm]{
		\item IDM \hfill 0,020\textsuperscript{*}
		\item IDM saúde \hfill 0,024\textsuperscript{*}
		\item IDM renda \hfill 0,012\textsuperscript{*}
		\item IDM educação \hfill NS\phantom{\textsuperscript{*}}
		\item Transbordamento renda \hfill 0,002\textsuperscript{*}
	} \\

	\cline{2-7}
	& \citeonline{Gondim2025}
	& Rev. Cadernos de Finanças Públicas
	& \fc[2cm]{\item NE \item 2010-2021}
	& \fc[1.5cm]{\item IF-SUDENE (+ FNE, FDNE)}
	& \fc[2cm]{\item $\Delta$ ln(PIB\textit{pc}), $\Delta$ ln(VAB\textit{pc})}
	& \fc[4.2cm]{
		\item VAB agropecuário\textit{pc} \hfill Significativo\textsuperscript{*}
		\item PIB\textit{pc} \hfill NS\phantom{\textsuperscript{*}}
	} \\

	\hline
Painel Dinâmico GMM & \citeonline{Souza2025}
	& AC
	& \fc[2cm]{\item N, NE, CO \item 2002-2021}
	& \fc[1.5cm]{\item IF-SUDENE, IF-SUDAM (+ FC, FD)}
	& \fc[2cm]{\item PIB\textit{pc}, emprego, renda}
	& \fc[4.2cm]{
		\item IF-SUDAM inconclusivo (poucas observações)
		\item Heterogeneidade entre instrumentos e regiões
	} \\

\end{longtable}%
\normalsize%
\setlength{\extrarowheight}{0pt}%
\setlength{\tabcolsep}{6pt}%
\linespread{1.5}\selectfont%
\renewcommand{\tablename}{Tabela}%
\renewcommand{\thetable}{\arabic{table}}%
}


No âmbito do mercado de trabalho, \citeonline{Garsous2017} conduzem a primeira avaliação de impacto dos incentivos fiscais da SUDENE, focalizando o setor turístico. Empregando modelo de diferenças em diferenças com \textit{matching} em municípios de Bahia, Rio de Janeiro, Espírito Santo e Minas Gerais no período 1998--2009, os autores estimam que o programa de benefícios tributários ao turismo elevou o emprego setorial em 30\% no acumulado entre 2002 e 2009. Os resultados mostram-se robustos à possibilidade de deslocamento de empregos para municípios vizinhos, embora os pesquisadores não avaliem a eficiência do programa em termos de custo fiscal por emprego gerado.

Em perspectiva mais abrangente, \citeonline{Braz2023} avaliam o efeito agregado da renúncia fiscal da SUDENE sobre o emprego formal e a renda nos municípios beneficiados, empregando estimador DID escalonado para o período 2002--2021. Os resultados indicam incremento de 3,2\% nos vínculos formais e de 1,2\% na renda nominal média, com trajetória crescente ao longo do tempo e duradoura por todo o período de vigência do benefício. Os autores destacam, porém, que os ganhos concentram-se em municípios de maior porte e dinamismo econômico, o que contrasta com o objetivo redistributivo da PNDR e levanta questões sobre a focalização territorial do instrumento.

\citeonline{Costa2024} adotam o método de controle sintético generalizado para estimar o impacto dos benefícios tributários da SUDENE sobre a criação de empregos em empresas da região Nordeste, com dados de 2006 a 2019. As evidências apontam elevação de 19,6\% nos vínculos formais no conjunto dos setores e de 18,3\% no setor manufatureiro, com persistência dos resultados por até cinco anos após o início da concessão. A magnitude consideravelmente superior à estimada por \citeonline{Braz2023} pode refletir o nível de agregação --- empresa versus município --- e a concentração no setor industrial, caracterizado por maior intensidade de mão de obra entre as atividades incentivadas.

Os três estudos convergem em apontar impacto positivo sobre o emprego formal, porém com magnitudes que variam de 3,2\% a 30\% conforme o setor analisado, o método empregado e a escala de agregação. Os ganhos mais expressivos concentram-se em segmentos específicos --- turismo e manufatura ---, enquanto o efeito agregado sobre o conjunto da economia é mais modesto. A preocupação distributiva levantada por \citeonline{Braz2023} merece atenção: se os benefícios se concentram em municípios já dinâmicos, o instrumento pode estar ampliando, em vez de reduzir, as disparidades intrarregionais.

Em abordagem complementar às avaliações de impacto, \citeonline{Carneiro2023} analisam a eficiência produtiva das empresas beneficiadas pela política de incentivos da SUDENE no período 2011--2019, combinando análise envoltória de dados (DEA) e fronteira estocástica (SFA). O estudo identifica predominância de empresas intensivas em mão de obra entre as beneficiadas e presença relevante de unidades eficientes no Semiárido. A ineficiência observada está associada mais ao setor de atividade incentivado do que às firmas individualmente, o que sugere possibilidade de aprimorar a seletividade setorial da política. Esses achados indicam que a eficiência no uso dos benefícios tributários não é uniforme e pode ser aprimorada mediante critérios diferenciados de concessão.

No que se refere ao produto e a indicadores de desenvolvimento, as evidências são mais escassas e derivam majoritariamente de estudos multi-instrumento discutidos nas subseções anteriores. \citeonline{Carneiroetal2024a}, ao avaliarem conjuntamente FNE, FDNE e incentivos fiscais, não encontram significância estatística para a renúncia fiscal da SUDENE sobre o PIB \textit{per capita} municipal nem sobre o valor adicionado setorial. Em contrapartida, \citeonline{FerreiraIrffiCarneiro2024} obtêm efeito direto significativo dos benefícios tributários sobre o IDM-geral (0,020), IDM-saúde (0,024) e IDM-renda (0,012), além de transbordamento espacial positivo sobre a componente renda nos municípios vizinhos (0,002). A divergência entre os dois estudos pode estar associada à diferença na variável de resultado e na especificação econométrica adotada.

Complementarmente, \citeonline{gondim2025} identificam efeito direto significativo da renúncia fiscal sobre o crescimento do VAB agropecuário \textit{per capita}, resultado que sugere influência do instrumento sobre setores primários não captada pelos estudos com foco no produto agregado. \citeonline{Souza2025}, ao avaliarem simultaneamente todos os instrumentos da PNDR em painel dinâmico, constituem a única referência que inclui os benefícios tributários da SUDAM, embora os resultados para esse instrumento permaneçam inconclusivos dada a escassez de observações.

Dois estudos avaliam o Programa de Desenvolvimento do Estado de Pernambuco (Prodepe), incentivo fiscal estadual que opera em complementaridade com o FNE, o FDNE e a renúncia fiscal da SUDENE. \citeonline{Oliveira2020a} empregam modelo DID com efeitos heterogêneos por tempo de exposição e dados de empresas pernambucanas no período 2000--2017, estimando incremento de 8,6\% no emprego das firmas beneficiadas. Simultaneamente, o estudo identifica redução de 10,3\% no salário médio, padrão que contrasta com a ausência de efeito salarial observada nos estudos sobre Fundos Constitucionais e sobre os incentivos federais. Não se obteve significância para a massa salarial em empresas que recebem exclusivamente o Prodepe, embora a combinação com outros instrumentos altere os resultados.

\citeonline{Alves2024} confirmam esse padrão empregando estimador DID escalonado ao mesmo recorte geográfico e temporal. As estimativas apontam elevação de 22,3\% no emprego e de 15,1\% na massa salarial, acompanhadas de redução de 8,2\% no salário médio. Os ganhos são mais pronunciados em empresas do Semiárido, de pequeno porte e do setor industrial, porém apresentam caráter temporário, com declínio progressivo ao longo dos anos de vigência do benefício. A redução salarial, consistente nos dois estudos, sugere que os incentivos estaduais atraem ou expandem atividades de menor remuneração, o que coloca em questão a qualidade dos postos de trabalho gerados.

Em síntese, as evidências sobre a renúncia fiscal revelam convergência no tocante ao emprego --- com impactos positivos variando de 3,2\% a 30\% conforme setor e método --- e divergência quanto ao produto e ao salário. Nenhum estudo identifica efeito significativo exclusivo dos incentivos fiscais sobre o PIB \textit{per capita}, embora resultados positivos sejam observados sobre indicadores multidimensionais de desenvolvimento. No mercado de trabalho, a política federal gera empregos sem alterar o salário médio, ao passo que o programa estadual eleva o emprego mas reduz a remuneração, suscitando preocupação sobre a qualidade dos postos criados. A literatura apresenta lacunas relevantes: ausência de avaliações dedicadas aos incentivos da SUDAM, inexistência de análises de custo-efetividade e predominância de métodos DID, com escassa aplicação de abordagens como RDD, PSM ou equilíbrio geral. A ampliação e diversificação metodológica desta vertente constituem prioridade para a agenda de pesquisa.

 no main.tex

A literatura empírica sobre os instrumentos da Política Nacional de Desenvolvimento Regional apresenta expressiva assimetria em termos de volume de evidências: os Fundos Constitucionais concentram a quase totalidade dos estudos, enquanto os Fundos de Desenvolvimento e os Incentivos Fiscais constituem vertentes consideravelmente menos desenvolvidas. Os resultados obtidos variam substancialmente conforme a especificação econométrica adotada, o período amostral, a escala geográfica de análise e o controle de características locais não observáveis. Esta seção apresenta os principais achados da literatura, organizados por instrumento e abordagem metodológica, buscando identificar para cada um deles convergências, divergências e condições que modulam a eficácia da política.

\subsection{Avaliações de Impacto dos Fundos Constitucionais sobre o PIB}
\label{subsec:fc-pib}

Os Fundos Constitucionais são os instrumentos da PNDR mais amplamente avaliados pela literatura empírica: dentre os 34 estudos aprovados nesta revisão, 28 investigam os efeitos desses fundos sobre variáveis de resultado relacionadas ao produto interno bruto ou ao mercado de trabalho. A questão central é se esses fundos --- que operam via crédito subsidiado de ampla capilaridade territorial --- efetivamente contribuem para o crescimento econômico dos municípios beneficiados e, em caso afirmativo, sob quais condições e com que magnitude. O FNE é o fundo com maior número de avaliações (24 estudos), seguido pelo FCO (10 estudos) e pelo FNO (8 estudos). A predominância dos Fundos Constitucionais na literatura decorre de sua maior longevidade em relação aos demais instrumentos da PNDR -- com operação contínua desde 1989 --, da disponibilidade de dados municipais anuais de PIB a partir de 2002, do elevado volume de recursos envolvidos e da capilaridade dos programas em todos os municípios das macrorregiões-alvo. A quase totalidade dos estudos adota o PIB \textit{per capita} municipal como variável de resultado, tratando-o como aproximação da renda média local, dada a escassez de alternativas em frequência anual e cobertura municipal. Os resultados obtidos pela literatura variam consideravelmente conforme a especificação econométrica adotada, o período amostral, a escala geográfica de análise e o controle de características locais não observáveis. As subseções seguintes organizam os achados por abordagem metodológica, permitindo identificar padrões gerais e divergências nas estimativas de impacto. O Quadro~\ref{tab:estudos_fc_pib} apresenta a síntese dos 21 estudos que avaliam os efeitos dos Fundos Constitucionais sobre o PIB \textit{per capita}, organizados por método de análise.

% Arquivo: latex/tabelas/tabela_macros.tex
% Macros para formatação de tabelas de resultados da revisão sistemática

% Macro principal para formatação de células com listas (label "--")
% Uso: \fc[largura]{conteúdo em itemize}
% Padrão: largura = 4.2cm (coluna Estimativa)
\newcommand{\fc}[2][4.2cm]{%
	\parbox[t]{#1}{\raggedright%
		\begin{itemize}[leftmargin=*,nosep,topsep=0pt,before=\vspace{-0.6\baselineskip},label={--},labelsep=0.4em]
			#2
		\end{itemize}%
		\vspace{0.3\baselineskip}%
	}%
}

% Macro para célula com label "--" (variáveis, amostragem)
% Uso: \fcvar[largura]{conteúdo em itemize}
% Padrão: largura = 2.2cm (colunas Amostragem e Variáveis)
\newcommand{\fcvar}[2][2.2cm]{%
	\parbox[t]{#1}{\raggedright%
		\begin{itemize}[leftmargin=*,nosep,topsep=0pt,before=\vspace{-0.6\baselineskip},label={--},labelsep=0.4em]
			#2
		\end{itemize}%
		\vspace{0.3\baselineskip}%
	}%
}

% Macro para mesclar células na coluna de método
% Parâmetros: #1 = texto, #2 = número de linhas a mesclar
\newcommand{\metodo}[2]{%
	\multirow{#2}{*}{\begin{minipage}[t]{1.4cm}\raggedright #1 \end{minipage}}%
}

% Arquivo: latex/tabelas/survey_artigos_fc_pib.tex
% Gerado automaticamente por scripts/generate_survey_tables.py
% Fonte: data/2-papers/survey_results.json

% IMPORTANTE: Este arquivo requer:
% - Pacotes: longtable, multirow, afterpage, enumitem, array
% - Macros: \fc e \metodo (definidas em tabelas/tabela_macros.tex)
% - Comando \citeonline (do pacote abntex2cite)
% - Contador: quadro (definido na classe abntex2)

\afterpage{\clearpage%
\tiny%
\setlength{\extrarowheight}{1pt}%
\setlength{\tabcolsep}{2pt}%
\linespread{0.9}\selectfont%
\renewcommand{\arraystretch}{1.2}%
\renewcommand{\tablename}{\quadroname}%
\renewcommand{\thetable}{\arabic{quadro}}%
\stepcounter{quadro}%
\begin{longtable}{|>{\raggedright\arraybackslash}p{2.2cm}
		|>{\raggedright\arraybackslash}p{2.2cm}
		|>{\raggedright\arraybackslash}p{2.3cm}
		|>{\raggedright\arraybackslash}p{3.0cm}
		|>{\raggedright\arraybackslash}p{4.2cm}|}

	\caption{Estudos de Impacto dos Fundos Constitucionais sobre o PIB \textit{per capita}} \label{tab:estudos_fc_pib} \\

\hline
	\begin{tabular}[c]{@{}>{\centering\arraybackslash}m{2.2cm}@{}}Método\end{tabular}
	& \begin{tabular}[c]{@{}>{\centering\arraybackslash}m{2.2cm}@{}}Artigo\end{tabular}
	& \begin{tabular}[c]{@{}>{\centering\arraybackslash}m{2.3cm}@{}}Amostragem\end{tabular}
	& \begin{tabular}[c]{@{}>{\centering\arraybackslash}m{3.0cm}@{}}Variáveis\end{tabular}
	& \begin{tabular}[c]{@{}>{\centering\arraybackslash}m{4.2cm}@{}}Estimativa\end{tabular} \\
	\hline
	\endfirsthead

	\multicolumn{5}{c}%
	{{\quadroname\ \thetable{} -- Continuação}} \\
	\hline
	\begin{tabular}[c]{@{}>{\centering\arraybackslash}m{2.2cm}@{}}Método\end{tabular}
	& \begin{tabular}[c]{@{}>{\centering\arraybackslash}m{2.2cm}@{}}Artigo\end{tabular}
	& \begin{tabular}[c]{@{}>{\centering\arraybackslash}m{2.3cm}@{}}Amostragem\end{tabular}
	& \begin{tabular}[c]{@{}>{\centering\arraybackslash}m{3.0cm}@{}}Variáveis\end{tabular}
	& \begin{tabular}[c]{@{}>{\centering\arraybackslash}m{4.2cm}@{}}Estimativa\end{tabular} \\
	\hline
	\endhead

	\hline
	\multicolumn{5}{r}{{Continua}} \\
	\endfoot

	\hline
	\multicolumn{5}{l}{
		\tiny
		\parbox{14.8cm}{%
			\vspace{0.1cm}
			Fonte: Elaborado pelo autor. Siglas: AC - Apresentação em congresso; TD - Trabalho para discussão; NS - Não significativo.
	}}
	\endlastfoot

Painel de Efeito Fixo
	& \citeonline{Resende2014a}
	& \fcvar[2.3cm]{\item NE \item 2004-2010}
	& \fcvar[3.0cm]{\item FC/PIB \item Cresc. PIB pc}
	& \fc[4.2cm]{
		\item FNE (munic) \hfill 0,021\textsuperscript{*}
		\item FNE (micro) \hfill 0,032\textsuperscript{*}
		\item FNE (meso) \hfill --0,288\textsuperscript{*}
		\item FNE-setorial \hfill NS\phantom{\textsuperscript{*}}
	}
	\\
	\cline{2-5}


	& \citeonline{Resende2014b}
	& \fcvar[2.3cm]{\item N \item 2004-2010}
	& \fcvar[3.0cm]{\item FC/PIB \item Cresc. PIB pc}
	& \fc[4.2cm]{
		\item FNO (munic) \hfill --0,399\textsuperscript{*}
		\item FNO (micro e meso) \hfill NS\phantom{\textsuperscript{*}}
		\item FNO setorial \hfill NS\phantom{\textsuperscript{*}}
	}
	\\
	\cline{2-5}


	& \citeonline{ResendeCravo2014}
	& \fcvar[2.3cm]{\item CO \item 2004-2010}
	& \fcvar[3.0cm]{\item FC/PIB \item Cresc. PIB pc}
	& \fc[4.2cm]{
		\item FCO (munic) \hfill 0,076\textsuperscript{*}
		\item FCO (micro e meso) \hfill NS\phantom{\textsuperscript{*}}
		\item FCO-empresarial (munic) \hfill 0,119\textsuperscript{*}
		\item FCO-rural (munic) \hfill --0,072\textsuperscript{*}
	}
	\\

	\hline
Modelo de \textit{Spillover} Espacial
	& \citeonline{Oliveira2005}
	& \fcvar[2.3cm]{\item N, CO \item 1991-2000}
	& \fcvar[3.0cm]{\item FC/PIB \item Cresc. PIB pc}
	& \fc[4.2cm]{
		\item FNO e FCO \hfill NS\phantom{\textsuperscript{*}}
	}
	\\
	\cline{2-5}


	& \citeonline{Cravo2015}
	& \fcvar[2.3cm]{\item N, NE, CO \item 2004-2010}
	& \fcvar[3.0cm]{\item FC/PIB \item Cresc. PIB pc}
	& \fc[4.2cm]{
		\item FNE (munic) \hfill NS\phantom{\textsuperscript{*}}
		\item FNO (munic) \hfill --0,313\textsuperscript{*}
		\item FCO (munic) \hfill 0,087\textsuperscript{*}
		\item FCO-empresarial (munic) \hfill +0,129\textsuperscript{*}
		\item \textit{Spillovers} espaciais \hfill NS\phantom{\textsuperscript{*}}
	}
	\\
	\cline{2-5}


	& \citeonline{Oliveira2016}
	& \fcvar[2.3cm]{\item GO \item 2004-2011}
	& \fcvar[3.0cm]{\item Log. FCO \item Log. PIB pc}
	& \fc[4.2cm]{
		\item FCO (painel espacial) \hfill 0,098\textsuperscript{*}
		\item FCO (painel clássico) \hfill 0,102\textsuperscript{*}
	}
	\\
	\cline{2-5}


	& \citeonline{Resende2017}
	& \fcvar[2.3cm]{\item NE \item 1999-2011}
	& \fcvar[3.0cm]{\item FC/PIB \item Cresc. PIB pc}
	& \fc[4.2cm]{
		\item FNE (munic. tipologia dinâmica)
		\item[\textbullet] Efeito direto \hfill 0,077\textsuperscript{*}
		\item[\textbullet] Efeito indireto \hfill 0,329\textsuperscript{*}
	}
	\\
	\cline{2-5}


	& \citeonline{Resende2018}
	& \fcvar[2.3cm]{\item N, NE, CO \item 1999-2011}
	& \fcvar[3.0cm]{\item FC/PIB \item Cresc. PIB pc}
	& \fc[4.2cm]{
		\item FC (munic. tip. dinâmica - acumulado)
		\item[\textbullet] Efeito direto \hfill 0,033\textsuperscript{*}
		\item[\textbullet] Efeito indireto \hfill 0,118\textsuperscript{*}
		\item FC (micro. tip. baixa renda - acumulado)
		\item[\textbullet] Efeito direto \hfill 0,146\textsuperscript{*}
		\item[\textbullet] Efeito indireto \hfill 0,466\textsuperscript{*}
	}
	\\
	\cline{2-5}


	& \citeonline{Silva2023}
	& \fcvar[2.3cm]{\item N, NE, CO \item 2016-2019}
	& \fcvar[3.0cm]{\item Log. FC \item Log. PIB pc}
	& \fc[4.2cm]{
		\item FNE, FNO e FCO \hfill NS\phantom{\textsuperscript{*}}
	}
	\\
	\cline{2-5}


	& \citeonline{Oliveira2026}
	& \fcvar[2.3cm]{\item NE \item 2010-2021}
	& \fcvar[3.0cm]{\item FC/PIB \item Cresc. PIB pc}
	& \fc[4.2cm]{
		\item FNE (efeito direto) \hfill 0,00001\textsuperscript{*}
		\item FNE (efeito indireto) \hfill NS\phantom{\textsuperscript{*}}
		\item FNE-serviços (efeito direto) \hfill 0,00001\textsuperscript{*}
		\item FNE-indústria (efeito direto) \hfill 0,00005\textsuperscript{*}
	}
	\\

	\hline
Modelo \textit{Threshold}
	& \citeonline{Linhares2014}
	& \fcvar[2.3cm]{\item NE \item 2000-2008}
	& \fcvar[3.0cm]{\item FC/PIB \item Cresc. PIB pc}
	& \fc[4.2cm]{
		\item FNE (PIB pc < 2.143) \hfill NS\phantom{\textsuperscript{*}}
		\item FNE (2.143 < PIB pc < 3.866) \hfill 0,078\textsuperscript{*}
		\item FNE (3.866 < PIB pc < 7.406) \hfill 0,109\textsuperscript{*}
		\item FNE (PIB pc > 7.406) \hfill NS\phantom{\textsuperscript{*}}
	}
	\\

	\hline
Primeira Diferença
	& \citeonline{Resende2014c}
	& \fcvar[2.3cm]{\item NE \item 1999-2006}
	& \fcvar[3.0cm]{\item FC/PIB \item Cresc. PIB pc}
	& \fc[4.2cm]{
		\item FNE-indústria (munic) \hfill NS\phantom{\textsuperscript{*}}
	}
	\\

	\hline
Regressão Quantílica e IV
	& \citeonline{Irffi2016}
	& \fcvar[2.3cm]{\item NE \item 2000-2010}
	& \fcvar[3.0cm]{\item FC/PIB \item Cresc. PIB pc}
	& \fc[4.2cm]{
		\item FNE \hfill (+)\textsuperscript{*}
		\item FNE-pecuária e agricultura \hfill (+)\textsuperscript{*}
	}
	\\

	\hline
Modelo de Eficiência
	& \citeonline{Carneiro2018}
	& \fcvar[2.3cm]{\item NE \item 2010-2015}
	& \fcvar[3.0cm]{\item FC/PIB \item Cresc. PIB pc}
	& \fc[4.2cm]{
		\item FNE: Maior eficiência em municípios integrados ao comércio internacional (Matopiba).
	}
	\\

	\hline
Painel Dinâmico
	& \citeonline{Cambota2019}
	& \fcvar[2.3cm]{\item NE \item 2003-2014}
	& \fcvar[3.0cm]{\item FC/PIB \item Cresc. PIB pc}
	& \fc[4.2cm]{
		\item FNE \hfill 2,96\textsuperscript{*}
		\item Efeito contemporâneo e endógeno significativo.
	}
	\\

	\hline
Equilíbrio Geral
	& \citeonline{Nascimento2017}
	& \fcvar[2.3cm]{\item BR \item 2000-2011}
	& \fcvar[3.0cm]{\item FC \item PIB}
	& \fc[4.2cm]{
		\item Remoção do FNE (alocação em despesa corrente) \hfill --0,16\%\phantom{\textsuperscript{*}}
		\item Maior concentração regional da produção
	}
	\\
	\cline{2-5}


	& \citeonline{Ribeiro2020}
	& \fcvar[2.3cm]{\item BR \item 2014-2025}
	& \fcvar[3.0cm]{\item FC/PIB \item Cresc. PIB pc}
	& \fc[4.2cm]{
		\item FNE (PIB NE) \hfill 3,51\%\textsuperscript{*}
	}
	\\
	\cline{2-5}


	& \citeonline{Goncalves2022}
	& \fcvar[2.3cm]{\item NE \item 2006-2015}
	& \fcvar[3.0cm]{\item Fundo \item Log. PIB pc}
	& \fc[4.2cm]{
		\item FNE (remoção): redução da atividade econômica no NE
	}
	\\

	\hline
RDD Geográfico
	& \citeonline{Oliveira2021}
	& \fcvar[2.3cm]{\item NE \item 1970-2010}
	& \fcvar[3.0cm]{\item Log. FNE \item Log. PIB pc}
	& \fc[4.2cm]{
		\item FNE \hfill NS\phantom{\textsuperscript{*}}
	}
	\\

	\hline
\textit{Propensity Score}
	& \citeonline{Monte2025}
	& \fcvar[2.3cm]{\item NE \item 2010-2020}
	& \fcvar[3.0cm]{\item FNE \item Log. PIB pc}
	& \fc[4.2cm]{
		\item FNE (PJ) \hfill 0,15pp\textsuperscript{*}
		\item FNE (mulheres) \hfill 0,10pp\textsuperscript{*}
		\item Efeito heterogêneo por magnitude
	}
	\\

	\hline
DID 2 Estágios, DID Escalonado
	& \citeonline{CarneiroVeloso2024}
	& \fcvar[2.3cm]{\item NE \item 2002-2019}
	& \fcvar[3.0cm]{\item FNE (\textit{dummy}) \item Log. PIB pc}
	& \fc[4.2cm]{
		\item FNE \hfill 0,08\textsuperscript{*} a 0,11\textsuperscript{*}
		\item A partir do 5º ano após o tratamento
		\item Principalmente serviços e administração pública\hfill \phantom{\textsuperscript{*}}
	}
	\\

\end{longtable}%
\normalsize%
\setlength{\extrarowheight}{0pt}%
\setlength{\tabcolsep}{6pt}%
\linespread{1.5}\selectfont%
\renewcommand{\tablename}{Tabela}%
\renewcommand{\thetable}{\arabic{table}}%
}


\subsubsection{Modelos de Efeitos Fixos}

Modelos de painel com efeitos fixos controlam características não observáveis constantes no tempo (como geografia e instituições locais) e choques temporais comuns a todas as unidades (como crises macroeconômicas), permitindo isolar de forma mais precisa o impacto da política regional. Essa abordagem foi a primeira empregada para avaliar os Fundos Constitucionais sobre o PIB municipal, aproveitando a disponibilidade de dados em painel desde 2002.

\citeonline{Resende2014a} analisa a influência do FNE sobre o crescimento do PIB \textit{per capita} dos municípios nordestinos no período 2004-2010, utilizando médias em três subperíodos (2004-2006, 2006-2008 e 2008-2010) para minimizar viés de ciclos de negócios. Os resultados apontam impacto positivo do FNE-total em nível municipal (0,021) e microrregional (0,032), mas não significativo em nível mesorregional, o que Resende (2014a) atribui à maior heterogeneidade em regiões geográficas maiores. As estimativas sugerem que o resultado positivo estaria concentrado nos empréstimos ao setor agropecuário, porém não se mostraram robustas à inclusão de controles temporais. Para o FNO, aplicando a mesma especificação econométrica à região Norte, \citeonline{Resende2014a} obtém relação inversa entre o fundo e a expansão do produto municipal, com ausência de robustez nas análises setoriais quando se controla por períodos.

\citeonline{Oliveira2017a} aplicam especificações de efeitos fixos juntamente com outros métodos (PSM e GPS) para avaliar o FCO no estado de Goiás, obtendo evidências de impacto positivo sobre a dinâmica econômica municipal, com magnitude variando conforme o método empregado.

Mais recentemente, \citeonline{Filho2024} adotam painel de efeitos fixos para avaliar conjuntamente os três fundos constitucionais, não obtendo evidências consistentes de associação positiva com o crescimento do PIB \textit{per capita} para o período 2016-2019, resultado que contrasta com os estudos anteriores.

Esses estudos ressaltam limitações metodológicas importantes. A principal dificuldade consiste na identificação causal, dado que a alocação dos fundos é endógena à demanda por crédito subsidiado, o que tende a privilegiar municípios com maior capacidade de absorção e infraestrutura socioeconômica mais desenvolvida. Adicionalmente, identificam-se dificuldades em mensurar os custos fiscais efetivos dos programas, a possibilidade de existência de peso morto (investimentos que seriam realizados mesmo sem subsídio) e potenciais deslocamentos de fatores produtivos entre municípios vizinhos. A ausência de robustez à inclusão de controles temporais sugere que parte do impacto estimado pode capturar tendências macroeconômicas comuns aos municípios tratados e não tratados, em vez do efeito causal dos fundos. Em síntese, os modelos de efeitos fixos indicam correlação positiva dos fundos constitucionais com a expansão econômica local, com impactos mais pronunciados em nível municipal e maior magnitude para o FCO, mas a fragilidade da identificação causal e a ausência de robustez em diversas especificações limitam a interpretação causal dos resultados.

\subsubsection{Modelos de Painel Espacial}

Uma vertente da literatura admite a possibilidade de interdependência econômica entre unidades espaciais geograficamente próximas, o que pode ocasionar efeitos indiretos de transbordamento da política para outras regiões, expandindo a área potencialmente beneficiada. Por outro lado, o efeito poderia consistir em vazamentos de recursos das regiões-alvo da política em direção a regiões de maior dinamismo econômico. Portanto, entender a natureza da dinâmica espacial é importante para um planejamento mais eficiente da política. \citeonline{Resende2014c} utilizam modelo de painel SDM (\textit{Spatial Durbin Model}) e estimam efeito significativo do FCO sobre o PIB \textit{per capita}, com evidência de que uma elevação de 1 ponto percentual (p.p.) na relação FCO/PIB está associada a um incremento de cerca de 0,09 p.p. na taxa média anual de crescimento do PIB \textit{per capita} municipal nos dois anos seguintes. Na análise setorial, o resultado seria atribuído à modalidade de crédito FCO-Empresarial, em que o efeito seria na ordem de 0,13 p.p. de incremento. Porém, o resultado é não significativo para o FNE e negativo para o FNO. Com respeito aos efeitos indiretos (transbordamento espacial), obteve-se resultado relevante, mas não positivo, para o FCO, o que pode estar ligado à migração de fatores de produção.

\citeonline{Oliveira2015} usam modelo espacial para avaliar efeitos do FCO nos municípios de Goiás e concluem que efeitos diretos são menores quando se admite a ocorrência de efeitos espaciais, sugerindo que modelos não espaciais tendem a superestimar o impacto real dos fundos. Entretanto, as evidências obtidas por \citeonline{Resende2014c} e \citeonline{Oliveira2015} apontam em direção contrária, com os efeitos diretos do FCO maiores em modelos espaciais do que em modelos não espaciais. Essa divergência pode estar associada a diferenças na cobertura geográfica (Goiás versus toda a região Centro-Oeste), no tipo de matriz de pesos espaciais empregada (\textit{queen} de contiguidade versus distância inversa) e no período de análise considerado. Os resultados de \citeonline{Oliveira2015} indicam que cada 1\% de aumento no valor dos financiamentos está associado a aproximadamente 0,1\% de crescimento no PIB \textit{per capita} do município, usando especificação SAR com matriz espacial \textit{queen} do tipo contiguidade e dados de 2004 a 2011.

\citeonline{ResendeSilvaFilho2017} distinguem-se dos trabalhos anteriores por incorporar na análise a heterogeneidade definida pela própria PNDR por meio das tipologias regionais. Em modelo não espacial de efeitos fixos, o efeito estimado do FNE é significativo para municípios pertencentes às tipologias dinâmica e baixa renda. Quando se admite dependência espacial, há evidências de transbordamento espacial dos investimentos do FNE quando o município diretamente beneficiado tem tipologia dinâmica. As estimativas apontam que para cada 1 p.p. de aumento na relação FNE/PIB em municípios dinâmicos há um incremento de 0,07 p.p. na taxa média anual de crescimento do PIB \textit{per capita} do próprio município nos três anos seguintes e de 0,33 p.p. na taxa média dos municípios vizinhos. Para os autores, esses resultados ressaltam a relevância de se considerar a dimensão espacial nas análises e também o caráter heterogêneo dos efeitos do FNE nos territórios. \citeonline{ResendeSilvaFilho2018} adotam a mesma abordagem econométrica, mas incluindo também as regiões Norte e Centro-Oeste, a fim de avaliar conjuntamente os fundos FNE, FNO e FCO. Os resultados são positivos para efeitos diretos e indiretos em municípios de tipologia dinâmica quando se usam valores acumulados na variável dos fundos (dois primeiros anos de cada subperíodo). Os resultados apontam que um aumento de 1 p.p. na relação FC/PIB em municípios dinâmicos está associado a 0,03 p.p. de incremento no crescimento médio do PIB \textit{per capita} do próprio município pelos três anos seguintes e 0,12 p.p. de incremento nos municípios vizinhos. \citeonline{Filho2024} não obtêm evidências de associação espacial positiva entre os financiamentos do FNE, FNO e FCO nos anos de 2016 a 2018 e o crescimento médio do PIB \textit{per capita} dos municípios de 2016 a 2019, porém, o resultado é positivo sobre o PIB \textit{per capita} de 2019. Os autores concluem que os efeitos dos fundos são em sua maior parte esporádicos e pontuais.

Em resumo, os estudos convergem para a relevância de se incorporar a dimensão espacial nas análises de efeito dos fundos constitucionais, justificada pela ocorrência de correlação espacial significativa nas variáveis envolvidas. Porém, há menor consenso quando se trata da magnitude dos efeitos. Em geral, o resultado estaria condicionado a outras características econômicas locais. O efeito direto mostrou-se significativo em municípios de tipologia dinâmica e microrregiões de baixa renda, com incrementos de 0,03 a 0,14 p.p. no crescimento anual do PIB \textit{per capita} municipal por até dois anos, para cada 1 p.p. de aumento na relação Fundo/PIB. Para o efeito indireto, as estimativas variam entre 0,1 e 0,4 p.p. de incremento no crescimento para cada 1 p.p. de aumento em Fundo/PIB, também concentrado em municípios de tipologia dinâmica e microrregiões de baixa renda. Em análise mais desagregada, evidências apontam para associação entre o FCO e taxa de crescimento, atribuída à linha de crédito empresarial. No FNE, os resultados estariam ligados ao setor industrial e de serviços. Já para o FNO não há nenhuma evidência positiva e significativa. No aspecto metodológico, nota-se a necessidade de estudos mais robustos, que combinem intervalo de tempo mais longo, especificações espaciais alternativas, correção de endogeneidade e maior número de controles (inclusive variáveis de confusão da própria PNDR).

\subsubsection{Modelos Não Lineares e de Eficiência}

\citeonline{Linharesetal2014} usam modelo \textit{threshold} e obtêm evidências de que o efeito do FNE é pouco significativo para municípios com PIB \textit{per capita} inicial abaixo de R\$ 2.143 ou acima de R\$ 7.406. Entre esses limites, os autores identificam dois grupos distintos. O primeiro é formado por municípios de PIB \textit{per capita} inicial entre R\$ 2.143 e R\$ 3.866. Nesse grupo, estima-se que 1 p.p. de aumento médio em FNE/PIB no período 2002-2006 esteja associado a um incremento de 0,08 p.p. no crescimento médio do PIB \textit{per capita} entre 2002 e 2008. No segundo grupo, de municípios com PIB \textit{per capita} inicial entre R\$ 3.866 e R\$ 7.406, o efeito estimado foi maior, na ordem de 0,11 p.p. Os resultados sugerem que o FNE tem maior efetividade em municípios com níveis intermediários de renda, reafirmando os achados de \citeonline{ResendeSilvaFilho2017} e \citeonline{ResendeSilvaFilho2018}. Em municípios de baixa renda, porém, os resultados não foram significativos, o que indica a necessidade de aprimoramento dos mecanismos da política para melhor adequação às características e necessidades dessas regiões.

\citeonline{Oliveira2015} usam modelagem quantílica para estimar a significância e heterogeneidade de efeitos do FNE ao longo da distribuição de crescimento dos municípios, empregando estimação em dois estágios e a localização de agências bancárias do BNB como instrumento para a alocação de recursos do fundo. Os efeitos foram positivos e significativos para o FNE total e para as linhas de crédito associadas à agricultura e pecuária. No entanto, os largos intervalos de confiança não permitiram constatar presença de heterogeneidade ao longo da curva de crescimento. \citeonline{Oliveira2018} analisa a relação entre a participação do FNE e a eficiência média dos municípios nordestinos. Usando modelo de análise envoltória de dados e de fronteira estocástica, o autor identifica elevada eficiência nos municípios pertencentes à região do MATOPIBA, associada à agricultura extensiva voltada à exportação. Adicionalmente, municípios do Semiárido indicaram menor eficiência no crescimento do produto médio, enquanto municípios menos densos e mais distantes das respectivas capitais e com empresas menores tendem a fazer melhor uso dos financiamentos.

\subsubsection{Modelos Dinâmicos}

\citeonline{Cambota2019} adotam modelo de painel dinâmico e instrumentos sintéticos para corrigir o viés de endogeneidade na variável PIB defasada e na variável de interesse FNE/PIB, usando dados de 2003 a 2014. Os autores confirmam a importância desse tipo de controle para a obtenção de estimativas mais realistas do impacto da política. As estimativas apontam associação positiva entre o FNE e o crescimento econômico dos municípios, na ordem de 2,96 p.p. para cada 1 p.p. de aumento em FNE/PIB no ano anterior. Cabe notar que essa magnitude é consideravelmente superior à dos demais estudos revisados, cujos coeficientes situam-se tipicamente entre 0,03 e 0,13 p.p. A discrepância pode ser atribuída à diferente especificação do modelo (painel dinâmico com instrumentos sintéticos), ao período mais longo de análise (2003-2014) ou a possível sobreestimação decorrente da proliferação de instrumentos no procedimento GMM, aspecto reconhecido na literatura como potencial fonte de viés.

\subsubsection{Modelos de Equilíbrio Geral}

Em \citeonline{Nascimento2017}, um cenário hipotético de retirada do FNE e realocação dos recursos em gastos correntes do governo resultaria a longo prazo em uma redução líquida de 0,16\% no PIB nacional e em maior concentração regional da produção. O estudo também sugere que os fundos têm papel mais relevante no crescimento de longo prazo, via elevação do estoque de capital e da capacidade produtiva local, frente à alternativa apontada. No curto prazo, ambos os cenários seriam equivalentes em termos de ativação da demanda. \citeonline{Ribeiroetal2020} indicam que os investimentos do FNE realizados em 2014 e 2015 teriam elevado o PIB da região Nordeste em 3,51\% até o ano 2025. \citeonline{Oliveira2021} constroem dois cenários alternativos envolvendo o término da política de subsídios, e ambos indicam como consequência a redução na atividade econômica na região Nordeste, puxada principalmente pela redução do produto da atividade agropecuária.

\subsubsection{Modelos Quase Experimentais}

\citeonline{Resende2014c} usa o método de primeira diferença para avaliar o efeito de empréstimos industriais do FNE sobre o crescimento do PIB \textit{per capita} municipal. Usando taxas de crescimento do PIB \textit{per capita} nos períodos 2000-2003 e 2000-2006, o autor não encontra efeito significativo das operações envolvendo o FNE-Indústria durante o ano 2000, seja em nível municipal ou microrregional. Argumenta-se que o resultado pode estar associado à elevada participação do setor agropecuário nas operações do FNE, setor em que há predomínio da informalidade, o que dificulta sua captura pelas medições do PIB.

\citeonline{Oliveira2020} adotam desenho quase-experimental que combina regressão com descontinuidade geográfica (RDD) e modelo de diferenças em diferenças. Os autores constroem um painel de dados decenal (1980, 1990, 2000 e 2010) composto por municípios dos estados de Minas Gerais e Espírito Santo que ingressaram nos critérios de elegibilidade ao FNE em períodos recentes (dados pré e pós-ingresso) e por municípios vizinhos imediatos que nunca tenham sido tratados pelo FNE. Os resultados obtidos levantam dúvidas sobre a efetividade do FNE, sugerindo que o Fundo não teve influência sobre as variáveis de produto e renda nos municípios avaliados no ano 2010. O estudo ressalta a importância de melhor vinculação do crédito ao aumento da produtividade, mediante condições diferenciadas e combinação com outros instrumentos.

Por outro lado, \citeonline{Carneiroetal2024a} obtêm evidências de efeito positivo e significativo do FNE sobre o PIB \textit{per capita} municipal. Os autores usam dados de 2002 a 2019 junto a estimadores DID dois estágios e DID escalonado, combinando três instrumentos de política regional: FNE, FDNE e incentivo fiscal da SUDENE. O estudo aponta efeito médio de tratamento de 8,8\% a 11\% no PIB \textit{per capita} do município beneficiado em relação ao grupo de controle, sendo atribuído principalmente ao setor serviços e administração pública. O estudo de eventos aponta que o efeito de tratamento é crescente e inicia no quinto ano após o primeiro tratamento. Quanto à sinergia com os demais instrumentos, o resultado é não significativo. \citeonline{MonteIrffiBastosCarneiro2025} exploram heterogeneidades de gênero, tipo de tomador e magnitude dos empréstimos (efeito dosagem) do FNE, aplicando o método \textit{Generalized Propensity Score} (GPS). O estudo encontra correlação positiva entre a participação do FNE-Mulheres no valor total do programa e a variação do PIB \textit{per capita} municipal, o que sinaliza uma eficácia maior do programa nesse grupo populacional. Efeito semelhante é obtido em relação à variável participação de tomadores pessoa jurídica no total de financiamentos. Com respeito ao efeito dosagem, os resultados apontam efeitos positivos no PIB \textit{per capita} somente na última e mais elevada dose de tratamento.

Em conjunto, as evidências sobre o efeito dos Fundos Constitucionais no PIB \textit{per capita} revelam um quadro heterogêneo e condicionado à especificação metodológica adotada. Os modelos de efeitos fixos e espaciais indicam associação positiva para o FCO e resultados ambíguos para FNE e FNO, enquanto os modelos quase experimentais apontam tanto para ausência de efeito quanto para impactos positivos significativos. Os modelos não lineares e de eficiência sugerem que a eficácia dos fundos é condicionada ao nível inicial de renda do município, com efeitos mais pronunciados em faixas intermediárias de desenvolvimento. As simulações de equilíbrio geral reafirmam a importância dos fundos para a atividade econômica regional em cenários contrafactuais de retirada da política. A constatação de que os efeitos são mais pronunciados em municípios com níveis intermediários de renda, e nulos ou inconclusivos em municípios de baixa renda, pode ser interpretada à luz da literatura sobre capacidade de absorção e armadilhas de pobreza: municípios de baixa renda podem carecer do capital humano e da infraestrutura básica necessários para converter o crédito subsidiado em crescimento econômico sustentado; municípios de renda intermediária possuiriam o mínimo de capacidade necessária para ativar o efeito multiplicador dos investimentos; e municípios de alta renda estariam próximos do efeito marginal decrescente do capital.

\subsection{Avaliações de Impacto dos Fundos Constitucionais sobre o Mercado de Trabalho}

Dezesseis estudos aprovados nesta revisão avaliam o impacto dos Fundos Constitucionais sobre variáveis do mercado de trabalho, com foco predominante no emprego formal, massa salarial e salário médio. O Quadro~\ref{tab:estudos_fc_emprego} apresenta a síntese desses estudos, organizados por método de análise.

% Arquivo: latex/tabelas/survey_artigos_fc_emprego.tex
% Gerado automaticamente por scripts/generate_survey_tables.py
% Fonte: data/2-papers/survey_results.json

% IMPORTANTE: Este arquivo requer:
% - Pacotes: longtable, multirow, afterpage, enumitem, array
% - Macros: \fc e \metodo (definidas em tabelas/tabela_macros.tex)
% - Comando \citeonline (do pacote abntex2cite)
% - Contador: quadro (definido na classe abntex2)

\afterpage{\clearpage%
\tiny%
\setlength{\extrarowheight}{1pt}%
\setlength{\tabcolsep}{2pt}%
\linespread{0.9}\selectfont%
\renewcommand{\arraystretch}{1.2}%
\renewcommand{\tablename}{\quadroname}%
\renewcommand{\thetable}{\arabic{quadro}}%
\stepcounter{quadro}%
\begin{longtable}{|>{\raggedright\arraybackslash}p{2.2cm}
		|>{\raggedright\arraybackslash}p{2.2cm}
		|>{\raggedright\arraybackslash}p{2.3cm}
		|>{\raggedright\arraybackslash}p{3.0cm}
		|>{\raggedright\arraybackslash}p{4.2cm}|}

	\caption{Estudos de Impacto dos Fundos Constitucionais sobre o Mercado de Trabalho} \label{tab:estudos_fc_emprego} \\

\hline
	\begin{tabular}[c]{@{}>{\centering\arraybackslash}m{2.2cm}@{}}Método\end{tabular}
	& \begin{tabular}[c]{@{}>{\centering\arraybackslash}m{2.2cm}@{}}Artigo\end{tabular}
	& \begin{tabular}[c]{@{}>{\centering\arraybackslash}m{2.3cm}@{}}Amostragem\end{tabular}
	& \begin{tabular}[c]{@{}>{\centering\arraybackslash}m{3.0cm}@{}}Variáveis\end{tabular}
	& \begin{tabular}[c]{@{}>{\centering\arraybackslash}m{4.2cm}@{}}Estimativa\end{tabular} \\
	\hline
	\endfirsthead

	\multicolumn{5}{c}%
	{{\quadroname\ \thetable{} -- Continuação}} \\
	\hline
	\begin{tabular}[c]{@{}>{\centering\arraybackslash}m{2.2cm}@{}}Método\end{tabular}
	& \begin{tabular}[c]{@{}>{\centering\arraybackslash}m{2.2cm}@{}}Artigo\end{tabular}
	& \begin{tabular}[c]{@{}>{\centering\arraybackslash}m{2.3cm}@{}}Amostragem\end{tabular}
	& \begin{tabular}[c]{@{}>{\centering\arraybackslash}m{3.0cm}@{}}Variáveis\end{tabular}
	& \begin{tabular}[c]{@{}>{\centering\arraybackslash}m{4.2cm}@{}}Estimativa\end{tabular} \\
	\hline
	\endhead

	\hline
	\multicolumn{5}{r}{{Continua}} \\
	\endfoot

	\hline
	\multicolumn{5}{l}{
		\tiny
		\parbox{14.8cm}{%
			\vspace{0.1cm}
			Fonte: Elaborado pelo autor. Siglas: NS - Não significativo; p.p. - Pontos percentuais.
	}}
	\endlastfoot

\textit{Propensity Score Matching}
	& \citeonline{SilvaResende2009}
	& \fcvar[2.3cm]{\item Empresas NE, N, CO \item 2000-2003}
	& \fcvar[3.0cm]{\item FNE, FNO, FCO (\textit{dummy}) \item Cresc. emprego \item Cresc. salário}
	& \fc[4.2cm]{
		\item FNE-emprego \hfill +61,2 p.p.\textsuperscript{*}
		\item FNE-salário \hfill NS\phantom{\textsuperscript{*}}
		\item FNO, FCO \hfill NS\phantom{\textsuperscript{*}}
	}
	\\
	\cline{2-5}


	& \citeonline{Soares2017}
	& \fcvar[2.3cm]{\item Empresas NE \item 2000-2006}
	& \fcvar[3.0cm]{\item FNE (\textit{dummy}) \item Cresc. emprego \item Cresc. massa salarial \item Cresc. salário}
	& \fc[4.2cm]{
		\item FNE-empr. (1 ano) \hfill +7,96 p.p.\textsuperscript{*}
		\item FNE-empr. (3 anos) \hfill +33,59 p.p.\textsuperscript{*}
		\item FNE-empr. (5 anos) \hfill +132,23 p.p.\textsuperscript{*}
		\item FNE-massa salarial \hfill (+)\textsuperscript{*}
		\item FNE-salário \hfill NS\phantom{\textsuperscript{*}}
	}
	\\

	\hline
\textit{Generalized Propensity Score}
	& \citeonline{OliveiraTerra2018}
	& \fcvar[2.3cm]{\item Empresas GO \item 2004-2011}
	& \fcvar[3.0cm]{\item FCO-Empresarial (contínuo) \item Cresc. emprego \item Cresc. salário}
	& \fc[4.2cm]{
		\item FCO-emprego \hfill (+)\textsuperscript{*}
		\item Efeito limitado a intervalo específico da dose
		\item FCO-salário \hfill NS\phantom{\textsuperscript{*}}
	}
	\\

	\hline
IV e Função Dose-Resposta
	& \citeonline{Cunha2024}
	& \fcvar[2.3cm]{\item Empresas NE \item 2000-2006}
	& \fcvar[3.0cm]{\item FNE (contínua) \item Cresc. emprego \item Cresc. massa salarial}
	& \fc[4.2cm]{
		\item Dose-resposta emprego: IV \hfill 357,54\%\textsuperscript{*}
		\item Dose-resposta emprego: OLS \hfill 133,94\%\textsuperscript{*}
		\item Dose-resposta massa: IV \hfill 220,96\%\textsuperscript{*}
		\item Dose-resposta massa: OLS \hfill 188,55\%\textsuperscript{*}
	}
	\\

	\hline
Primeira Diferença
	& \citeonline{Resende2014c}
	& \fcvar[2.3cm]{\item Empresas NE \item 1999-2006}
	& \fcvar[3.0cm]{\item FNE-Indústria (\textit{dummy}) \item Cresc. emprego}
	& \fc[4.2cm]{
		\item FNE-Ind. (anual) \hfill +10,83 a 16,11 p.p.\textsuperscript{*}
		\item FNE-Ind. (3 anos) \hfill +30 a 56 p.p.\textsuperscript{*}
	}
	\\

	\hline
Painel Espacial (SAR)
	& \citeonline{Oliveira2015}
	& \fcvar[2.3cm]{\item Municípios GO \item 2004-2011}
	& \fcvar[3.0cm]{\item Log. FCO \item Log. emprego \item Salário médio}
	& \fc[4.2cm]{
		\item FCO-emprego \hfill (+)\textsuperscript{*}
		\item FCO-salário \hfill (+)\textsuperscript{*}
		\item Sem diferença relevante vs. painel clássico
	}
	\\

	\hline
Painel Dinâmico
	& \citeonline{Ribeiro2024}
	& \fcvar[2.3cm]{\item Municípios NE \item 2010-2015}
	& \fcvar[3.0cm]{\item Log. FNE \item Log. emprego por setor}
	& \fc[4.2cm]{
		\item FNE-primário \hfill +1,10\%\textsuperscript{*}
		\item FNE-terciário \hfill +0,53\%\textsuperscript{*}
		\item FNE-secundário \hfill +0,13\%\textsuperscript{*}
		\item FNE defasado (secundário) \hfill +0,06\%\textsuperscript{*}
	}
	\\

	\hline
Equilíbrio Geral
	& \citeonline{Ribeiro2020}
	& \fcvar[2.3cm]{\item Região NE \item 2014-2025}
	& \fcvar[3.0cm]{\item FNE \item Emprego total}
	& \fc[4.2cm]{
		\item FNE (emprego NE acumulado) \hfill +2,75\%\textsuperscript{*}
		\item NE-agropecuária \hfill +8,0\%\textsuperscript{*}
		\item NE-construção \hfill +5,3\%\textsuperscript{*}
	}
	\\

	\hline
Diferenças em Diferenças
	& \citeonline{Oliveira2020}
	& \fcvar[2.3cm]{\item Empresas N (FNO) \item 2000-2010}
	& \fcvar[3.0cm]{\item FNO (\textit{dummy}) \item Emprego \item Massa salarial \item Salário médio}
	& \fc[4.2cm]{
		\item Emprego e renda \hfill (+)\textsuperscript{*}
		\item Curto e médio prazo (+), exceto grandes empresas
		\item Investimento: efeito a longo prazo
		\item Giro/custeio: efeito a curto prazo
	}
	\\

	\hline
RDD Geográfico
	& \citeonline{Oliveira2021}
	& \fcvar[2.3cm]{\item Municípios NE \item 1970-2010}
	& \fcvar[3.0cm]{\item FNE \item Emprego formal e informal (2010)}
	& \fc[4.2cm]{
		\item Agropecuária \hfill +5,4\%\textsuperscript{*}
		\item Indústria \hfill +2,8\%\textsuperscript{*}
		\item Serviços \hfill +3,5\%\textsuperscript{*}
		\item Taxa de informalidade \hfill NS\phantom{\textsuperscript{*}}
	}
	\\

\end{longtable}%
\normalsize%
\setlength{\extrarowheight}{0pt}%
\setlength{\tabcolsep}{6pt}%
\linespread{1.5}\selectfont%
\renewcommand{\tablename}{Tabela}%
\renewcommand{\thetable}{\arabic{table}}%
}


\citeonline{Silva2009} foi o primeiro estudo a investigar o impacto dos empréstimos dos FCs sobre a contratação de trabalhadores nas empresas beneficiadas. O estudo faz a interseção dos registros de empresas beneficiadas com os dados da RAIS (Relação Anual de Informações Sociais), obtendo a quantidade de vínculos de empregos formais de 449 empresas (NE=211, N=174, CO=74) nos anos de 2000 e 2003. Aplicando o método \textit{Propensity Score Matching (PSM)}, os autores concluem que nas empresas beneficiadas pelo FNE a quantidade de empregos cresceu em média 61,2 pontos percentuais a mais do que nas empresas não beneficiadas. No entanto, os valores não foram significativos para o salário médio. E ainda, para o FNO e o FCO, não foi obtido efeito positivo em nenhum dos indicadores avaliados.

\citeonline{Soares2017} também aplicam o método PSM, mas fazem uso de um conjunto mais rico de informações sobre o FNE, o que permitiu ampliar consideravelmente o número de empresas e o horizonte de tempo investigado. Além disso, os autores fazem uma análise ano a ano dos efeitos do fundo, o que permite estimativas mais eficientes, com quantidades maiores de empresas nos primeiros anos pós-tratamento. Com dados do período de 2000 a 2006, no horizonte de 3 anos pós-tratamento há 1.104 empresas e em 5 anos há 314. Os resultados obtidos reforçam as evidências obtidas por \citeonline{Silva2009}. No primeiro ano pós-tratamento o crescimento do estoque de empregos é em média 7,96 pontos percentuais maior nas empresas beneficiadas, relativamente às demais. Além disso, a trajetória é crescente nos anos subsequentes, alcançando 33,59 pontos percentuais no terceiro ano e 132,23 no quinto ano. Se confirmados, esses resultados atribuem grande relevância para a política regional, indicando que o FNE tem papel importante na geração de empregos formais nas empresas beneficiadas. Os autores também destacam que o efeito de maior monta estaria associado às empresas financiadas entre 1999 e 2001, as únicas no horizonte de 5 anos pós-tratamento, estando mais concentradas no setor industrial. O mesmo estudo ainda analisa possíveis efeitos do FNE sobre a massa salarial, obtendo resultados significativos e com dinâmica semelhante ao do efeito no emprego. Como consequência, a variável salário médio não apresentou variação significativa como decorrência do FNE, indicando que as novas contratações teriam ocorrido ao mesmo nível dos salários vigentes.

\citeonline{Oliveira2018} apresentam como inovação principal o uso de variável de tratamento do tipo contínua, permitindo analisar possíveis heterogeneidades no efeito do fundo decorrentes da intensidade do tratamento (valor do empréstimo), em contraste à abordagem dicotômica de tratamento presente nos estudos anteriores. Os autores usam o método \textit{Generalized Propensity Score Matching (GPS)} e dados de financiamentos a empresas do estado de Goiás realizados pelo FCO por meio da linha Empresarial (associada à indústria, comércio e serviços) no período de 2004 a 2011. Os resultados indicam relação positiva e significativa com o crescimento do emprego apenas em um intervalo específico da distribuição de tratamento e limitada ao período 2004-2008 (pré-crise), cenário semelhante ao obtido na abordagem PSM com variável de tratamento dicotômica. Com respeito ao salário médio, não se obteve nenhum resultado significativo. \citeonline{Oliveira2018} expandem a abordagem de \citeonline{Oliveira2018} para os demais fundos constitucionais. O dado mais relevante obtido indica efeito decrescente dos fundos sobre as variáveis de interesse, o que significa que quanto maior o valor do empréstimo, menor o efeito sobre o crescimento do emprego e do salário. Para os autores, isso se justifica do ponto de vista teórico pela propriedade de retornos marginais decrescentes da função de produção neoclássica.

\citeonline{Junior2024} contribuem com a literatura ao usar o método de função dose-resposta para captar efeitos heterogêneos e correção de viés de seleção por meio de variável instrumental. Para garantir comparabilidade, os autores optam por usar a mesma base de dados de \citeonline{Soares2017}. Na variável crescimento do emprego, o modelo IV aponta efeito médio de 357,54\%, acima de 133,94\% do modelo OLS. Para o crescimento da massa salarial, o valor médio é 220,96\% no modelo IV contra 188,55\% em OLS. Em ambas as variáveis a restrição de exogeneidade leva a subestimação do efeito médio e superestimação do ponto ótimo de tratamento, sinalizando para efeitos mais significativos em pontos mais baixos da distribuição de tratamento, quando a endogeneidade é corrigida.

\citeonline{Resende2014c} contribui para uma limitação importante dos métodos PSM e GSM: a incapacidade de controlar por efeitos não observáveis que possam influenciar tanto a probabilidade de tratamento quanto seu resultado (apenas características observáveis). Usando abordagem de primeira diferença e dados de 1999 a 2006, o autor identifica efeito positivo do FNE-Indústria sobre o crescimento do número de empregos formais nas empresas beneficiadas, com magnitude entre 10,83 e 16,11 pontos percentuais ao ano. No horizonte de três anos pós-tratamento, o valor é de 30 a 56 pontos percentuais a mais, de certa forma consistente com os resultados de \citeonline{Silva2009} e \citeonline{Soares2017}.

\citeonline{Oliveira2015} é o primeiro estudo a avaliar o efeito dos fundos constitucionais sobre variáveis agregadas do mercado de trabalho. O estudo adota os municípios de Goiás como unidade de análise e aplica o método de painel espacial autorregressivo (SAR), a fim de controlar também por características da vizinhança do município. Os resultados indicam efeitos positivos sobre emprego e salário, porém, sem diferenças relevantes em relação ao modelo de painel clássico com efeitos fixos, o que pode ser decorrência de baixa robustez dos resultados. \citeonline{Ribeiro2024} trazem duas contribuições principais: análise dos efeitos agregados do FNE por setor de atividade econômica e o uso de modelo de painel dinâmico para amenizar possível viés de endogeneidade. Os resultados apontam relação positiva entre os investimentos do FNE e o nível de emprego no município, com maior efeito sobre o setor primário (1,1\% de aumento no emprego para 1\% de aumento no investimento do FNE), seguido por terciário (0,53\%) e secundário (0,13\%). Esses resultados são observados para os investimentos correntes do FNE. Para o FNE defasado em um ano, o valor é significativo somente para o setor secundário e em menor magnitude (0,06\%). Sugere-se assim que o FNE tem efeito positivo e imediato sobre o emprego, especialmente no setor primário, setor mais intensivo em mão de obra.

Em maior nível de agregação, \citeonline{Ribeiroetal2020} aplicam um modelo dinâmico interregional de equilíbrio geral computável (CGE) para estimar os efeitos dos investimentos FNE sobre o emprego total da região Nordeste, dos estados e dos setores de atividade. As evidências apontam que os investimentos do FNE nos anos 2014 e 2015 estão associados a uma variação acumulada de 2,75\% na quantidade de empregos na região Nordeste até o ano de 2025, valor pouco menor que o estimado para a variação do PIB. Na análise por estados, Piauí, Ceará e Rio Grande do Norte também apresentam a maior variação acumulada do emprego, entre 3,19\% e 4,31\%, e Maranhão e Bahia têm as menores variações, de 0,45\% a 2,16\%. No aspecto setorial, também não há perdedores, com agropecuária (8\%) e construção (5,3\%) os setores com maior variação do emprego.

\citeonline{Oliveira2020} empregam o método diferenças-em-diferenças, que controla pela heterogeneidade não observável das empresas e possibilita uma estimativa do efeito causal do fundo. Além disso, têm como diferencial a análise dos efeitos microeconômicos do FNO por porte da empresa beneficiada, tipo de crédito concedido e horizonte de tempo. No geral, empresas beneficiadas apresentaram desempenho superior às não beneficiadas, em curto e médio prazo, em termos de geração de emprego e renda. No entanto, médias e grandes empresas beneficiadas por financiamento a capital de giro e custeio apresentam desempenho similar ao grupo de controle, no que diz respeito à geração de emprego e renda. Além disso, os autores apontam que o impacto positivo do crédito com finalidade capital de giro e custeio é em média maior no curto prazo do que no longo prazo. Por outro lado, em empresas beneficiadas pela finalidade investimento, os impactos são mais observáveis a longo prazo. Os efeitos assimétricos indicam que a finalidade investimento promove a acumulação de capital, o que aumenta a produtividade do trabalho a longo prazo, pois o impacto é positivo para o trabalho e a renda (salário médio).

\citeonline{Oliveira2021} utilizam modelo de descontinuidade geográfica (RDD) para estimar o impacto causal do FNE sobre o estoque de emprego, distinguindo-se dos estudos anteriores pela robustez metodológica e por analisar impacto sobre o emprego formal e informal, decorrência do uso de dados censitários. Os resultados indicam que o acesso ao FNE gerou impacto positivo e significativo sobre o estoque de empregos formais e informais no ano 2010, com magnitude 5,4\% (agropecuária), 2,8\% (indústria) e 3,5\% (serviços). Porém, impacto sobre a taxa de informalidade foi nulo.

Em síntese, as avaliações sobre o mercado de trabalho apresentam maior convergência do que as relativas ao PIB: há evidências consistentes de efeito positivo do FNE sobre a geração de emprego formal nas empresas beneficiadas, com magnitude crescente ao longo dos primeiros anos pós-tratamento. Os estudos também convergem em apontar efeitos mais pronunciados em micro e pequenas empresas e na modalidade de crédito para investimento. Por outro lado, a variável salário médio não apresenta variação significativa na maioria das especificações, sugerindo que os novos postos de trabalho são criados ao nível salarial vigente, sem ganhos adicionais de produtividade. Os resultados para FNO e FCO são mais limitados e inconclusivos, refletindo menor cobertura na literatura.

\subsection{Avaliações de Impacto dos Fundos de Desenvolvimento}
\label{subsec:fd}

Os Fundos de Desenvolvimento (FDNE, FDA e FDCO) constituem instrumentos mais recentes da PNDR, com primeiras contratações a partir de 2007 no caso do FDNE e de 2014 para o FDCO, o que explica a menor quantidade de avaliações empíricas em comparação aos Fundos Constitucionais. A escala dos empreendimentos financiados --- mínimo de R\$ 15 milhões, com concentração em infraestrutura e indústria --- sugere mecanismo de impacto distinto, cuja avaliação é relevante para informar a alocação de recursos entre os instrumentos da política. Dentre os estudos aprovados nesta revisão, 8 avaliam algum aspecto dos Fundos de Desenvolvimento, com concentração expressiva no FDNE. Não foram identificadas avaliações quantitativas dedicadas ao FDA nem ao FDCO, lacuna relevante dada a magnitude dos recursos envolvidos --- R\$ 5,5 bilhões e R\$ 1,8 bilhão contratados, respectivamente. As subseções seguintes organizam as evidências por abordagem metodológica. O Quadro~\ref{tab:estudos_fd} apresenta a síntese dos estudos que avaliam os Fundos de Desenvolvimento.

% Arquivo: latex/tabelas/survey_artigos_fd.tex
% Quadro de estudos sobre impacto dos Fundos de Desenvolvimento
% Fonte: Revisão sistemática - dados extraídos do pipeline pndr_survey

% IMPORTANTE: Este arquivo requer:
% - Pacotes: longtable, multirow, afterpage, enumitem, array
% - Macros: \fc e \metodo (definidas em tabelas/tabela_macros.tex)
% - Comando \citeonline (do pacote abntex2cite)
% - Contador: quadro (definido no main.tex)

\afterpage{\clearpage%
\tiny%
\setlength{\extrarowheight}{1pt}%
\setlength{\tabcolsep}{2pt}%
\linespread{0.9}\selectfont%
\renewcommand{\arraystretch}{1.2}%
\renewcommand{\tablename}{\quadroname}%
\renewcommand{\thetable}{\arabic{quadro}}%
\stepcounter{quadro}%
\begin{longtable}{|>{\raggedright\arraybackslash}p{1.6cm}
		|>{\raggedright\arraybackslash}p{1.7cm}
		|>{\raggedright\arraybackslash}p{2cm}
		|>{\raggedright\arraybackslash}p{1.5cm}
		|>{\raggedright\arraybackslash}p{2.6cm}
		|>{\raggedright\arraybackslash}p{4.2cm}|}

	\caption{Estudos de Impacto dos Fundos de Desenvolvimento} \label{tab:estudos_fd} \\

\hline
	\begin{tabular}[c]{@{}>{\centering\arraybackslash}m{1.6cm}@{}}Método\end{tabular}
	& \begin{tabular}[c]{@{}>{\centering\arraybackslash}m{1.7cm}@{}}Artigo\end{tabular}
	& \begin{tabular}[c]{@{}>{\centering\arraybackslash}m{2cm}@{}}Amostragem\end{tabular}
	& \begin{tabular}[c]{@{}>{\centering\arraybackslash}m{1.5cm}@{}}Variável Independente\end{tabular}
	& \begin{tabular}[c]{@{}>{\centering\arraybackslash}m{2.6cm}@{}}Variável Dependente\end{tabular}
	& \begin{tabular}[c]{@{}>{\centering\arraybackslash}m{4.2cm}@{}}Resultado\end{tabular} \\
	\hline
	\endfirsthead

	\multicolumn{6}{c}%
	{{\quadroname\ \thetable{} -- Continuação}} \\
	\hline
	\begin{tabular}[c]{@{}>{\centering\arraybackslash}m{1.6cm}@{}}Método\end{tabular}
	& \begin{tabular}[c]{@{}>{\centering\arraybackslash}m{1.7cm}@{}}Artigo\end{tabular}
	& \begin{tabular}[c]{@{}>{\centering\arraybackslash}m{2cm}@{}}Amostragem\end{tabular}
	& \begin{tabular}[c]{@{}>{\centering\arraybackslash}m{1.5cm}@{}}Variável Independente\end{tabular}
	& \begin{tabular}[c]{@{}>{\centering\arraybackslash}m{2.6cm}@{}}Variável Dependente\end{tabular}
	& \begin{tabular}[c]{@{}>{\centering\arraybackslash}m{4.2cm}@{}}Resultado\end{tabular} \\
	\hline
	\endhead

	\hline
	\multicolumn{6}{r}{{Continua}} \\
	\endfoot

	\hline
	\multicolumn{6}{l}{
		\tiny
		\parbox{14.8cm}{%
			\vspace{0.1cm}
			Fonte: Elaborado pelo autor. Siglas: AC - Apresentação em congresso; NS - Não significativo.
	}}
	\endlastfoot

	\hline
Painel Espacial & \citeonline{FerreiraIrffiCarneiro2024}
	& \fc[2cm]{\item NE \item 2010-2021}
	& \fc[1.5cm]{\item FDNE}
	& \fc[2cm]{\item IDM}
	& \fc[4.5cm]{
		\item IDM renda \hfill 0,019\textsuperscript{*}
		\item IDM, IDM saúde e IDM educação \hfill NS\phantom{\textsuperscript{*}}
	} \\

	\cline{2-6}

	& \citeonline{gondim2025}
	& \fc[2cm]{\item NE \item 2010-2021}
	& \fc[1.5cm]{\item FDNE (\textit{dummy})}
	& \fc[2cm]{\item $\Delta$ ln(PIB pc)}
	& \fc[4.5cm]{
		\item Efeito direto \hfill NS\phantom{\textsuperscript{*}}
		\item Efeito indireto \hfill 0,0252\textsuperscript{*}
		\item Efeito total \hfill 0,0351\textsuperscript{*}
	} \\
	\cline{5-6}
	& & &
	& \fc[2cm]{\item $\Delta$ ln(VAB pc)}
	& \fc[4.5cm]{
		\item Efeito direto \hfill 0,0409\textsuperscript{*}
		\item Efeito indireto \hfill NS\phantom{\textsuperscript{*}}
		\item Efeito total \hfill 0,6014\textsuperscript{*}
	} \\

	\hline
DID Dois Estágios & \citeonline{Carneiroetal2024a}
	& \fc[2cm]{\item NE \item 2002-2019}
	& \fc[1.5cm]{\item FDNE}
	& \fc[2cm]{\item log PIB\textit{pc}, VAB setorial}
	& \fc[4.5cm]{
		\item PIB\textit{pc} \hfill 0,063\textsuperscript{*}
		\item VAB\textit{pc} serviços \hfill 0,111\textsuperscript{*}
		\item Arrecadação \hfill 0,207\textsuperscript{*}
		\item VAB ind. e VAB agro. \hfill NS\phantom{\textsuperscript{*}}
	} \\

	\hline
DID Escalonado
	& \citeonline{BrazBastosIrffi2024}
	& \fc[2cm]{\item NE \item 2002-2021}
	& \fc[1.5cm]{\item FDNE}
	& \fc[2cm]{\item PIB\textit{pc}, emprego, renda, IDEB}
	& \fc[4.5cm]{
		\item PIB\textit{pc} \hfill 24,07\%\textsuperscript{*}
		\item Renda \hfill 4,63\%\textsuperscript{*}
		\item IDEB anos iniciais \hfill 18,67\%\textsuperscript{*}
		\item IDEB anos finais \hfill 21,49\%\textsuperscript{*}
		\item Emprego, pobreza, mortalidade infantil e saúde \hfill NS\phantom{\textsuperscript{*}}
	} \\
	\cline{2-6}

	& \citeonline{Irffietal2025}
	& \fc[2cm]{\item NE \item 1999-2022}
	& \fc[1.5cm]{\item FDNE-parques eólicos}
	& \fc[2cm]{\item PIB\textit{pc}, VAB, emprego, massa salarial}
	& \fc[4.5cm]{
		\item PIB\textit{pc} \hfill 18,97\%\textsuperscript{*}
		\item VAB indústria \hfill 71,90\%\textsuperscript{*}
		\item VAB serviço \hfill 10,24\%\textsuperscript{*}
		\item Emprego \hfill 23,41\%\textsuperscript{*}
		\item Redução temporária do VAB-agro durante construção
	} \\

	\hline
Painel Dinâmico GMM
	& \citeonline{Souza2025}
	& \fc[2cm]{\item N, NE, CO \item 2002-2021}
	& \fc[1.5cm]{\item FDNE, FDA, FDCO (+ FC + IF)}
	& \fc[2cm]{\item PIB\textit{pc}, emprego, renda}
	& \fc[4.5cm]{
		\item Heterogeneidade entre fundos, regiões e setores
		\item FDCO inconclusivo (poucos projetos)
	} \\

\end{longtable}%
\normalsize%
\setlength{\extrarowheight}{0pt}%
\setlength{\tabcolsep}{6pt}%
\linespread{1.5}\selectfont%
\renewcommand{\tablename}{Tabela}%
\renewcommand{\thetable}{\arabic{table}}%
}


\subsubsection{Modelos Quase Experimentais}

O método de diferenças em diferenças (DID) e suas extensões --- como o DID escalonado, adequado para intervenções implementadas em momentos distintos para diferentes unidades --- constituem a principal abordagem empregada para avaliar os Fundos de Desenvolvimento. Os três estudos discutidos a seguir representam as primeiras evidências empíricas de impacto causal desses fundos sobre indicadores municipais.

\citeonline{Braz2024} aplicam modelo DID escalonado aos municípios da área de atuação da SUDENE no período 2002--2021 e estimam que aqueles que receberam empreendimentos apoiados pelo FDNE experimentam elevação média de 24\% no PIB \textit{per capita} em relação ao grupo de controle. O estudo inova ao avaliar variáveis de resultado não estritamente econômicas: a renda domiciliar \textit{per capita} apresenta incremento de 4,6\%, e os indicadores de educação básica (IDEB) registram elevação de 19\% nos anos iniciais e 21\% nos anos finais. Não se obteve significância estatística para emprego formal, pobreza, mortalidade infantil e indicadores de saúde.

\citeonline{Carneiroetal2024a} avaliam simultaneamente os efeitos do FNE, do FDNE e dos incentivos fiscais da SUDENE sobre a economia dos municípios nordestinos, empregando estimadores DID em dois estágios e DID escalonado para o período 2002--2019. Para o FDNE, identificam impacto positivo e significativo sobre o PIB \textit{per capita} municipal, com coeficiente concentrado no setor de serviços e na arrecadação tributária. Os incentivos fiscais não apresentaram significância estatística, e a interação entre os instrumentos tampouco indicou sinergia, sugerindo que os efeitos operam de forma independente. Os resultados desse estudo para o FNE foram discutidos na Subseção~\ref{subsec:fc-pib}.

\citeonline{Irffietal2025} estimam o impacto da construção de parques eólicos sobre indicadores econômicos dos municípios nordestinos, empregando o método de controle sintético generalizado e estudos de eventos para o período 1999--2022. Os resultados apontam elevação de 19\% no PIB \textit{per capita}, de 72\% no valor adicionado bruto (VAB) industrial e de 10\% no VAB de serviços nos municípios com parques em operação ou construção. Identificam-se ainda incrementos de 23\% no emprego formal e de 31\% na massa salarial, além de redução temporária do VAB agropecuário durante a fase de construção, sem efeito permanente. Embora o estudo avalie todos os parques eólicos da região --- e não exclusivamente os financiados pelo FDNE ---, parcela expressiva dos empreendimentos de energia eólica no Nordeste contou com recursos desse fundo, e a proximidade dos valores obtidos com os de \citeonline{Braz2024} reforça a consistência das conclusões.

As três avaliações convergem em apontar impacto positivo do FDNE sobre o PIB \textit{per capita}, com magnitudes situadas entre 19\% e 24\%, consideravelmente superiores às estimadas para os Fundos Constitucionais --- cujos coeficientes situam-se tipicamente entre 0,03 e 0,13 p.p. de incremento no crescimento anual para cada 1 p.p. de aumento na relação Fundo/PIB. Essa diferença pode estar associada à escala dos empreendimentos financiados, com valor mínimo de R\$ 15 milhões, à concentração setorial em infraestrutura e indústria de transformação e ao horizonte mais longo de maturação desses investimentos. Cabe ressalvar que as três avaliações empregam variantes do método DID e analisam períodos parcialmente sobrepostos, o que limita a independência das estimativas.

\subsubsection{Modelos Espaciais}

Modelos de painel espacial permitem decompor o impacto de uma política em componente direto --- sobre o próprio município beneficiado --- e indireto --- sobre municípios vizinhos. Essa distinção é particularmente relevante para os Fundos de Desenvolvimento, cujos empreendimentos financiados (infraestrutura de energia, transportes, complexos industriais) podem gerar externalidades que se dispersam para além dos limites municipais.

\citeonline{FerreiraIrffiCarneiro2024} utilizam modelo SAC (\textit{Spatial Autoregressive Combined}) estático para avaliar o FDNE e os incentivos fiscais da SUDENE nos municípios nordestinos entre 2010 e 2021. A variável de resultado adotada é o Índice de Desenvolvimento Municipal (IDM) e suas componentes. Para o FDNE, o estudo identifica efeito positivo e significativo apenas sobre a componente IDM-renda (coeficiente 0,019), sem evidência sobre o IDM-geral, IDM-saúde ou IDM-educação. A adoção de variável de resultado distinta do PIB \textit{per capita} dificulta a comparação direta com os demais estudos desta subseção, embora contribua para ampliar o escopo das dimensões avaliadas.

\citeonline{gondim2025} analisam conjuntamente FNE, FDNE e incentivos fiscais da SUDENE por meio de modelo espacial generalizado com dados de 2010 a 2021. Os resultados para o FDNE revelam padrão distinto dos demais instrumentos: o efeito direto sobre o PIB \textit{per capita} é não significativo, enquanto o efeito indireto (transbordamento para municípios vizinhos) apresenta coeficiente positivo e significativo de 0,025, resultando em efeito total de 0,035. Para o VAB \textit{per capita}, o efeito direto é significativo (0,041), mas o indireto não. Esse resultado sugere que os investimentos do FDNE beneficiam não apenas o município-sede, mas também localidades vizinhas, padrão consistente com a natureza dos projetos financiados --- infraestrutura de energia e transportes cuja área de influência transcende limites municipais. Tal configuração contrasta com os Fundos Constitucionais, cujos efeitos diretos tendem a ser mais pronunciados.

\subsubsection{Avaliação Integrada e Síntese}

\citeonline{Souza2025} adotam painel dinâmico GMM com defasagens distribuídas para avaliar simultaneamente todos os instrumentos da PNDR --- incluindo FDNE, FDA e FDCO --- no período 2002--2021. O estudo constitui a única avaliação que contempla os três fundos de desenvolvimento em uma mesma especificação. Os resultados indicam heterogeneidade expressiva entre fundos, regiões e setores, com evidências inconclusivas para o FDCO, atribuídas ao reduzido número de projetos e ao curto período de operação.

Em conjunto, as evidências disponíveis apontam para impacto positivo expressivo do FDNE sobre o PIB \textit{per capita} municipal, com magnitudes entre 19\% e 24\% nas avaliações quase experimentais e presença de efeitos de transbordamento espacial nos modelos de painel. A base empírica, porém, apresenta limitações relevantes: concentra-se em estudos recentes (2024--2026) com períodos amostrais parcialmente sobrepostos e predominância do DID escalonado; não dispõe de avaliação quantitativa dedicada ao FDA nem ao FDCO; e restringe-se majoritariamente ao PIB \textit{per capita}, com evidências incipientes sobre emprego, renda domiciliar e pobreza. A diversificação metodológica e a extensão da cobertura para os três fundos de desenvolvimento constituem prioridades para a agenda de pesquisa.

\subsection{Avaliações de Impacto dos Incentivos Fiscais}
\label{subsec:if}

A avaliação empírica dos incentivos fiscais da SUDENE e da SUDAM constitui a vertente menos desenvolvida da literatura sobre a PNDR, a despeito do expressivo volume de renúncia tributária envolvido. Dentre os 34 estudos aprovados nesta revisão, 9 avaliam explicitamente esses instrumentos, refletindo a operação mais recente dos programas --- com concessões sistemáticas a partir de 2010 --- e as dificuldades de acesso a dados detalhados de renúncia fiscal em nível de empresa. Diferentemente dos Fundos Constitucionais, cujo mecanismo de transmissão opera via crédito subsidiado, os incentivos fiscais atuam pela redução de 75\% do Imposto de Renda de Pessoa Jurídica (IRPJ), diminuindo o custo tributário de operação em municípios da área de atuação da superintendência. Todos os estudos dedicados concentram-se nos benefícios concedidos pela SUDENE, não havendo avaliações específicas da SUDAM. Dois trabalhos avaliam adicionalmente o Prodepe --- programa estadual de Pernambuco --- em interação com os instrumentos federais. O Quadro~\ref{tab:estudos_if} apresenta a síntese dos estudos desta subseção.

% Arquivo: latex/tabelas/survey_artigos_if.tex
% Quadro de estudos sobre impacto dos Incentivos Fiscais
% Fonte: Revisão sistemática - dados extraídos do pipeline pndr_survey
% Adaptado da tese: removido CostaCarneiroIrffi2023 (não aprovado),
%   adicionados Oliveira2020a, Alves2024, Gondim2025, Souza2025.
%   Chaves BibTeX corrigidas para consistência com references.bib.

% IMPORTANTE: Este arquivo requer:
% - Pacotes: longtable, multirow, afterpage, enumitem, array
% - Macros: \fc e \metodo (definidas em tabelas/tabela_macros.tex)
% - Comando \citeonline (do pacote abntex2cite)
% - Contador: quadro (definido no main.tex)

\afterpage{\clearpage%
\tiny%
\setlength{\extrarowheight}{1pt}%
\setlength{\tabcolsep}{2pt}%
\linespread{0.9}\selectfont%
\renewcommand{\arraystretch}{1.2}%
\renewcommand{\tablename}{\quadroname}%
\renewcommand{\thetable}{\arabic{quadro}}%
\stepcounter{quadro}%
\begin{longtable}{|>{\raggedright\arraybackslash}p{1.4cm}
		|>{\raggedright\arraybackslash}p{1.5cm}
		|>{\raggedright\arraybackslash}p{1.5cm}
		|>{\raggedright\arraybackslash}p{2cm}
		|>{\raggedright\arraybackslash}p{1.5cm}
		|>{\raggedright\arraybackslash}p{2cm}
		|>{\raggedright\arraybackslash}p{4.2cm}|}

	\caption{Estudos de Impacto dos Incentivos Fiscais} \label{tab:estudos_if} \\

	\hline
	\begin{tabular}[c]{@{}>{\centering\arraybackslash}m{1.4cm}@{}}Método\end{tabular}
	& \begin{tabular}[c]{@{}>{\centering\arraybackslash}m{1.5cm}@{}}Artigo\end{tabular}
	& \begin{tabular}[c]{@{}>{\centering\arraybackslash}m{1.5cm}@{}}Publicação\end{tabular}
	& \begin{tabular}[c]{@{}>{\centering\arraybackslash}m{2cm}@{}}Amostragem\end{tabular}
	& \begin{tabular}[c]{@{}>{\centering\arraybackslash}m{1.5cm}@{}}Variável Independente\end{tabular}
	& \begin{tabular}[c]{@{}>{\centering\arraybackslash}m{2cm}@{}}Variável Dependente\end{tabular}
	& \begin{tabular}[c]{@{}>{\centering\arraybackslash}m{4.2cm}@{}}Resultado\end{tabular} \\
	\hline
	\endfirsthead

	\multicolumn{7}{c}%
	{{\quadroname\ \thetable{} -- Continuação}} \\
	\hline
	\begin{tabular}[c]{@{}>{\centering\arraybackslash}m{1.4cm}@{}}Método\end{tabular}
	& \begin{tabular}[c]{@{}>{\centering\arraybackslash}m{1.5cm}@{}}Artigo\end{tabular}
	& \begin{tabular}[c]{@{}>{\centering\arraybackslash}m{1.5cm}@{}}Publicação\end{tabular}
	& \begin{tabular}[c]{@{}>{\centering\arraybackslash}m{2cm}@{}}Amostragem\end{tabular}
	& \begin{tabular}[c]{@{}>{\centering\arraybackslash}m{1.5cm}@{}}Variável Independente\end{tabular}
	& \begin{tabular}[c]{@{}>{\centering\arraybackslash}m{2cm}@{}}Variável Dependente\end{tabular}
	& \begin{tabular}[c]{@{}>{\centering\arraybackslash}m{4.2cm}@{}}Resultado\end{tabular} \\
	\hline
	\endhead

	\hline
	\multicolumn{7}{r}{{Continua}} \\
	\endfoot

	\hline
	\multicolumn{7}{l}{
		\tiny
		\parbox{14.8cm}{%
			\vspace{0.1cm}
			Fonte: Elaborado pelo autor. Siglas: AC - Apresentação em congresso; NS - Não significativo.
	}}
	\endlastfoot

	\hline
DID & \citeonline{Garsous2017}
	& \textit{World Development}
	& \fc[2cm]{\item BA, RJ, ES, MG \item 1998-2009}
	& \fc[1.5cm]{\item IF-SUDENE (turismo)}
	& \fc[2cm]{\item Emprego no turismo (log)}
	& \fc[4.2cm]{
		\item Emprego turismo \hfill 20\%\textsuperscript{*}, 30\%\textsuperscript{*}
	} \\

	\cline{2-7}
	& \citeonline{Oliveira2020a}
	& AC
	& \fc[2cm]{\item PE \item 2000-2017}
	& \fc[1.5cm]{\item Prodepe + IF-SUDENE}
	& \fc[2cm]{\item Emprego, salário médio, massa salarial}
	& \fc[4.2cm]{
		\item Emprego \hfill 8,6\%\textsuperscript{*}
		\item Salário médio \hfill $-$10,3\%\textsuperscript{*}
		\item Massa salarial \hfill NS\phantom{\textsuperscript{*}}
	} \\

	\hline
DID Escalonado & \citeonline{Braz2023}
	& AC
	& \fc[2cm]{\item SUDENE \item 2002-2021}
	& \fc[1.5cm]{\item IF-SUDENE}
	& \fc[2cm]{\item Emprego, renda}
	& \fc[4.2cm]{
		\item Emprego \hfill 3,2\%\textsuperscript{*}
		\item Renda \hfill 1,2\%\textsuperscript{*}
		\item Efeito crescente e duradouro
	} \\

	\cline{2-7}
	& \citeonline{Alves2024}
	& AC
	& \fc[2cm]{\item PE \item 2000-2017}
	& \fc[1.5cm]{\item Prodepe + IF-SUDENE}
	& \fc[2cm]{\item Emprego, salário médio, massa salarial}
	& \fc[4.2cm]{
		\item Emprego \hfill 22,3\%\textsuperscript{*}
		\item Massa salarial \hfill 15,1\%\textsuperscript{*}
		\item Salário médio \hfill $-$8,2\%\textsuperscript{*}
		\item Efeito temporário e declinante
	} \\

	\hline
DEA + SFA & \citeonline{Carneiro2023a}
	& Rev. Cadernos de Finanças Públicas
	& \fc[2cm]{\item NE \item 2011-2019}
	& \fc[1.5cm]{\item IF-SUDENE}
	& \fc[2cm]{\item Eficiência produtiva (massa salarial)}
	& \fc[4.2cm]{
		\item Empresas eficientes no Semiárido
		\item Intensivas em mão de obra
		\item Ineficiência associada ao setor, não à empresa
	} \\

	\hline
DID Dois Estágios & \citeonline{Carneiroetal2024a}
	& AC
	& \fc[2cm]{\item NE, MG, ES \item 2002-2019}
	& \fc[1.5cm]{\item IF-SUDENE (+ FNE, FDNE)}
	& \fc[2cm]{\item log PIB\textit{pc}, VAB setorial}
	& \fc[4.2cm]{
		\item IF-SUDENE \hfill NS\phantom{\textsuperscript{*}}
	} \\

	\hline
Controle Sintético & \citeonline{Costa2024}
	& \textit{Applied Economics Letters}
	& \fc[2cm]{\item NE \item 2006-2019}
	& \fc[1.5cm]{\item IF-SUDENE}
	& \fc[2cm]{\item Vínculos (log)}
	& \fc[4.2cm]{
		\item Todos os setores \hfill 19,6\%\textsuperscript{*}
		\item Indústria \hfill 18,3\%\textsuperscript{*}
		\item Efeitos por até 5 anos
		\item Demais setores \hfill NS\phantom{\textsuperscript{*}}
	} \\

	\hline
Painel Espacial & \citeonline{FerreiraIrffiCarneiro2024}
	& AC
	& \fc[2cm]{\item NE \item 2010-2021}
	& \fc[1.5cm]{\item IF-SUDENE (+ FDNE)}
	& \fc[2cm]{\item IDM e subcomponentes}
	& \fc[4.2cm]{
		\item IDM \hfill 0,020\textsuperscript{*}
		\item IDM saúde \hfill 0,024\textsuperscript{*}
		\item IDM renda \hfill 0,012\textsuperscript{*}
		\item IDM educação \hfill NS\phantom{\textsuperscript{*}}
		\item Transbordamento renda \hfill 0,002\textsuperscript{*}
	} \\

	\cline{2-7}
	& \citeonline{Gondim2025}
	& Rev. Cadernos de Finanças Públicas
	& \fc[2cm]{\item NE \item 2010-2021}
	& \fc[1.5cm]{\item IF-SUDENE (+ FNE, FDNE)}
	& \fc[2cm]{\item $\Delta$ ln(PIB\textit{pc}), $\Delta$ ln(VAB\textit{pc})}
	& \fc[4.2cm]{
		\item VAB agropecuário\textit{pc} \hfill Significativo\textsuperscript{*}
		\item PIB\textit{pc} \hfill NS\phantom{\textsuperscript{*}}
	} \\

	\hline
Painel Dinâmico GMM & \citeonline{Souza2025}
	& AC
	& \fc[2cm]{\item N, NE, CO \item 2002-2021}
	& \fc[1.5cm]{\item IF-SUDENE, IF-SUDAM (+ FC, FD)}
	& \fc[2cm]{\item PIB\textit{pc}, emprego, renda}
	& \fc[4.2cm]{
		\item IF-SUDAM inconclusivo (poucas observações)
		\item Heterogeneidade entre instrumentos e regiões
	} \\

\end{longtable}%
\normalsize%
\setlength{\extrarowheight}{0pt}%
\setlength{\tabcolsep}{6pt}%
\linespread{1.5}\selectfont%
\renewcommand{\tablename}{Tabela}%
\renewcommand{\thetable}{\arabic{table}}%
}


No âmbito do mercado de trabalho, \citeonline{Garsous2017} conduzem a primeira avaliação de impacto dos incentivos fiscais da SUDENE, focalizando o setor turístico. Empregando modelo de diferenças em diferenças com \textit{matching} em municípios de Bahia, Rio de Janeiro, Espírito Santo e Minas Gerais no período 1998--2009, os autores estimam que o programa de benefícios tributários ao turismo elevou o emprego setorial em 30\% no acumulado entre 2002 e 2009. Os resultados mostram-se robustos à possibilidade de deslocamento de empregos para municípios vizinhos, embora os pesquisadores não avaliem a eficiência do programa em termos de custo fiscal por emprego gerado.

Em perspectiva mais abrangente, \citeonline{Braz2023} avaliam o efeito agregado da renúncia fiscal da SUDENE sobre o emprego formal e a renda nos municípios beneficiados, empregando estimador DID escalonado para o período 2002--2021. Os resultados indicam incremento de 3,2\% nos vínculos formais e de 1,2\% na renda nominal média, com trajetória crescente ao longo do tempo e duradoura por todo o período de vigência do benefício. Os autores destacam, porém, que os ganhos concentram-se em municípios de maior porte e dinamismo econômico, o que contrasta com o objetivo redistributivo da PNDR e levanta questões sobre a focalização territorial do instrumento.

\citeonline{Costa2024} adotam o método de controle sintético generalizado para estimar o impacto dos benefícios tributários da SUDENE sobre a criação de empregos em empresas da região Nordeste, com dados de 2006 a 2019. As evidências apontam elevação de 19,6\% nos vínculos formais no conjunto dos setores e de 18,3\% no setor manufatureiro, com persistência dos resultados por até cinco anos após o início da concessão. A magnitude consideravelmente superior à estimada por \citeonline{Braz2023} pode refletir o nível de agregação --- empresa versus município --- e a concentração no setor industrial, caracterizado por maior intensidade de mão de obra entre as atividades incentivadas.

Os três estudos convergem em apontar impacto positivo sobre o emprego formal, porém com magnitudes que variam de 3,2\% a 30\% conforme o setor analisado, o método empregado e a escala de agregação. Os ganhos mais expressivos concentram-se em segmentos específicos --- turismo e manufatura ---, enquanto o efeito agregado sobre o conjunto da economia é mais modesto. A preocupação distributiva levantada por \citeonline{Braz2023} merece atenção: se os benefícios se concentram em municípios já dinâmicos, o instrumento pode estar ampliando, em vez de reduzir, as disparidades intrarregionais.

Em abordagem complementar às avaliações de impacto, \citeonline{Carneiro2023} analisam a eficiência produtiva das empresas beneficiadas pela política de incentivos da SUDENE no período 2011--2019, combinando análise envoltória de dados (DEA) e fronteira estocástica (SFA). O estudo identifica predominância de empresas intensivas em mão de obra entre as beneficiadas e presença relevante de unidades eficientes no Semiárido. A ineficiência observada está associada mais ao setor de atividade incentivado do que às firmas individualmente, o que sugere possibilidade de aprimorar a seletividade setorial da política. Esses achados indicam que a eficiência no uso dos benefícios tributários não é uniforme e pode ser aprimorada mediante critérios diferenciados de concessão.

No que se refere ao produto e a indicadores de desenvolvimento, as evidências são mais escassas e derivam majoritariamente de estudos multi-instrumento discutidos nas subseções anteriores. \citeonline{Carneiroetal2024a}, ao avaliarem conjuntamente FNE, FDNE e incentivos fiscais, não encontram significância estatística para a renúncia fiscal da SUDENE sobre o PIB \textit{per capita} municipal nem sobre o valor adicionado setorial. Em contrapartida, \citeonline{FerreiraIrffiCarneiro2024} obtêm efeito direto significativo dos benefícios tributários sobre o IDM-geral (0,020), IDM-saúde (0,024) e IDM-renda (0,012), além de transbordamento espacial positivo sobre a componente renda nos municípios vizinhos (0,002). A divergência entre os dois estudos pode estar associada à diferença na variável de resultado e na especificação econométrica adotada.

Complementarmente, \citeonline{gondim2025} identificam efeito direto significativo da renúncia fiscal sobre o crescimento do VAB agropecuário \textit{per capita}, resultado que sugere influência do instrumento sobre setores primários não captada pelos estudos com foco no produto agregado. \citeonline{Souza2025}, ao avaliarem simultaneamente todos os instrumentos da PNDR em painel dinâmico, constituem a única referência que inclui os benefícios tributários da SUDAM, embora os resultados para esse instrumento permaneçam inconclusivos dada a escassez de observações.

Dois estudos avaliam o Programa de Desenvolvimento do Estado de Pernambuco (Prodepe), incentivo fiscal estadual que opera em complementaridade com o FNE, o FDNE e a renúncia fiscal da SUDENE. \citeonline{Oliveira2020a} empregam modelo DID com efeitos heterogêneos por tempo de exposição e dados de empresas pernambucanas no período 2000--2017, estimando incremento de 8,6\% no emprego das firmas beneficiadas. Simultaneamente, o estudo identifica redução de 10,3\% no salário médio, padrão que contrasta com a ausência de efeito salarial observada nos estudos sobre Fundos Constitucionais e sobre os incentivos federais. Não se obteve significância para a massa salarial em empresas que recebem exclusivamente o Prodepe, embora a combinação com outros instrumentos altere os resultados.

\citeonline{Alves2024} confirmam esse padrão empregando estimador DID escalonado ao mesmo recorte geográfico e temporal. As estimativas apontam elevação de 22,3\% no emprego e de 15,1\% na massa salarial, acompanhadas de redução de 8,2\% no salário médio. Os ganhos são mais pronunciados em empresas do Semiárido, de pequeno porte e do setor industrial, porém apresentam caráter temporário, com declínio progressivo ao longo dos anos de vigência do benefício. A redução salarial, consistente nos dois estudos, sugere que os incentivos estaduais atraem ou expandem atividades de menor remuneração, o que coloca em questão a qualidade dos postos de trabalho gerados.

Em síntese, as evidências sobre a renúncia fiscal revelam convergência no tocante ao emprego --- com impactos positivos variando de 3,2\% a 30\% conforme setor e método --- e divergência quanto ao produto e ao salário. Nenhum estudo identifica efeito significativo exclusivo dos incentivos fiscais sobre o PIB \textit{per capita}, embora resultados positivos sejam observados sobre indicadores multidimensionais de desenvolvimento. No mercado de trabalho, a política federal gera empregos sem alterar o salário médio, ao passo que o programa estadual eleva o emprego mas reduz a remuneração, suscitando preocupação sobre a qualidade dos postos criados. A literatura apresenta lacunas relevantes: ausência de avaliações dedicadas aos incentivos da SUDAM, inexistência de análises de custo-efetividade e predominância de métodos DID, com escassa aplicação de abordagens como RDD, PSM ou equilíbrio geral. A ampliação e diversificação metodológica desta vertente constituem prioridade para a agenda de pesquisa.

